% !TeX root = ../../Book5.tex
% 纲要标题：## 第8章*　模表示与 Brauer 特征标
% 纲要位置：工作记录/计划纲要.md，第 3981 行。
\OptionalChapter{模表示与 Brauer 特征标}{b5:ch:08}
\BookChapterNumber{b5:ch:08}

\begin{chapterabstract}
	本章研究底域特征整除群阶时的有限群表示，借投射覆盖与块分析不可分解模，以 Brauer 特征标比较模特征与特征零的表示信息，再用相对投射性、顶点与源联系子群上的模。Brauer 对应比较具有指定缺陷群的块，Green 对应比较具有指定顶点的不可分解模；矩阵计算则检验这些对象在具体例子中的结构。
\end{chapterabstract}

\begin{learninggoals}
	\begin{itemize}
		\item 能构造半单性失效的例子，区分简单模、不可分解模及其根层。
		\item 能计算循环 \(p\) 群的模、投射覆盖、基座、Cartan 矩阵及稳定同态与扩张空间。
		\item 从中心本原幂等元得到块，并在小群的群代数中验证块分解。
		\item 能在 \(p\)-正则共轭类上计算 Brauer 特征标，并由格约化确定分解数。
		\item 能按定义计算 Brauer 映射及具体块的缺陷群，明确特征标记录合成因子的范围。
		\item 能用相对迹检验相对投射性，确定循环群模的顶点与源，并以稳定同态计算一阶扩张。
		\item 能区分 Brauer 对应与 Green 对应的对象、假设及结论，在小群中算出对应项。
		\item 能将模计算的输出还原为可检查的矩阵等式，区分不可约性、绝对不可约性及不可分解性。
	\end{itemize}
	预备知识为有限群表示、群代数、合成列及有限维代数的根；投射覆盖的构造使用第二卷的幂等元提升；顶点与源的论证另调用有限长度模的局部自同态环和 Krull--Schmidt 定理，并在使用处恢复具体条件。Brauer 特征标部分另固定一个适当的分裂 \(p\)-模系统，届时会说明特征零域、离散赋值环及剩余域各自的作用。
\end{learninggoals}

设 \(G=C_p=\langle g\rangle\)，而域 \(k\) 的特征为素数 \(p\)。在群代数中令 \(u=g-1\)，则 \(u^p=g^p-1=0\)，并有 \(kG\cong k[u]/(u^p)\)。这个代数的逐次商 \(u^j kG/u^{j+1}kG\) 都是一维平凡表示，但乘 \(u\) 的作用把它们接成一个不能直接拆开的模。简单成分因而不再决定整个表示，特征零下用平均投影取得不变补空间的办法也失效。投射覆盖、块与 Brauer 特征标将分别记录这里留下的结构和数值信息。投射覆盖及有限维代数 Cartan 矩阵的基础构造，可对照\cite[Sections 9.1--9.3, pp. 209--212]{EtingofEtAl2011RepresentationTheory}。

\ProseParagraphBreak
引用 Webb 著作时，本章采用作者2016年2月23日公开稿的节号与结果编号。作者说明该稿与正式刊本的数学内容相同；稿件页码与刊本页码另有区别，以下引用使用节号。格约化与 Brauer 特征标见\cite[Sections 9.4, 10.1; Section 12.5]{Webb2016FiniteGroupRepresentations}。

\ProseParagraphBreak
当不可分解模已经不能由特征标分类时，仍可询问它能从多小的子群诱导出来，以及限制到局部子群后哪些直和项能够辨认。相对迹把这个问题化成恒等映射的分解；顶点与源记录所需的子群及其中的模。对块和不可分解模分别建立对应以后，还要在矩阵中核验实际得到的是子模、商模还是直和项。本章后半的计算保留这些区别，程序输出旁同时给出可以手算的检验。

% !TeX root = ../../Book5.tex
% 纲要标题：### 8.1 半单性的失效
% 纲要位置：工作记录/计划纲要.md，第 3983 行。

\section{半单性的失效}
\label{b5:sec:0801}


% 纲要范围：
%
% * 特征 p 整除群阶
% * 群代数的根及不可分解模
% * 循环 p 群的完整基本例子
% * 中心 p 元在简单模上的作用；一般群代数的根包含于增广理想，p 群时相等
%

当底域的特征整除群阶时，平均投影中的除法不再有意义。表示仍然可以有合成列，但这些简单商之间的扩张未必分裂；因此研究简单模还不够，还须研究不可分解模及其根层。本章以循环 \(p\) 群为反复计算的例子，始终区分直和分解与合成因子。

\subsection{\texorpdfstring{特征 \(p\) 整除群阶的情形}{特征 p 整除群阶的情形}}
\label{b5:subsec:0801:01}

\begin{definition}\label{b5:def:modular-representation}
设 \(G\) 为有限群，\(k\) 为特征 \(p>0\) 的域。\(G\) 的有限维 \(k\) 表示称为特征 \(p\) 的\term[mo-biaoshi]{模表示}[modular representation]。本章重点讨论 \(p\,\bigm\vert\,|G|\) 的情形；若 \(p\mathrel{\vcenter{\hbox{\scalebox{1.25}{\(\nmid\)}}}}|G|\)，\cref{b5:thm:maschke}仍保证完全可约。
\end{definition}
由\cref{b5:prop:maschke-converse}，当 \(p\,\bigm\vert\,|G|\) 时，增广映射
\[
\varepsilon:kG\longrightarrow k,\qquad \sum_g a_gg\longmapsto\sum_g a_g
\]
没有 \(kG\) 线性截面。具体地，一个不动元必须为 \(c\sum_g g\)，它的增广为 \(c|G|=0\)，所以不能提升目标中的一。这既给出不完全可约的正则模，也给出取不动空间不保持满射的例子。
\ProseParagraphBreak
最小的矩阵例子是 \(C_p=\langle s\rangle\) 在 \(k^2\) 上作用为
\[
s\longmapsto\begin{pmatrix}1&1\\0&1\end{pmatrix}.
\]
第一坐标轴和商都为平凡表示，但没有不变补。这个表示与平凡二维表示有相同的两个简单合成因子，却不同构：前者的生成元不是恒等。后面定义的 Brauer 特征标会记录合成因子的重数，不会额外记录这条扩张是否分裂。

\subsection{群代数的根与不可分解模}
\label{b5:subsec:0801:02}

% 跨卷引用目标：b2:thm:module-radical-jm（第二卷定理15.2.2）。
设 \(A\) 为有限维含幺 \(k\) 代数。其 Jacobson 根可写为全部简单左模的零化子的交，记为 \(J=J(A)\)。若 \(M\) 为有限维左模，全部极大子模的交称为其\term[mo-gen]{根}[radical of a module]。它等于 \(JM\)，如第二卷定理15.2.2所证：每个极大商都是被 \(J\) 零化的简单模，故 \(JM\) 包含在每个极大子模中；反过来，\(M/JM\) 半单，任一非零商类都能在某个简单分量上得到非零投影，该投影的核给出一个不含原元素的极大子模。商 \(M/JM\) 称为\term[mo-ding]{顶}[top of a module]。根层 \(J^iM/J^{i+1}M\) 记录从顶逐步向下的简单成分。它们通常不能在原模中同时选成直和项；即使全部根层逐项同构，原模也未必同构。\subsecref{b5:subsec:0804:03}将用 \(C_2\times C_2\) 的一参数模族计算这种现象：各层相同，层间的作用仍可不同。
\symidx{\(J(A),JM,M/JM\)：代数的 Jacobson 根与模的根、顶}
\begin{proposition}\label{b5:prop:finite-radical-layers}
设 \(k\) 为域，\(A\) 为有限维含幺 \(k\) 代数，\(J=J(A)\)。则 \(J\) 幂零，\(A/J\) 半单。对任意有限维左 \(A\) 模 \(M\)，全部根层 \(J^iM/J^{i+1}M\) 都是半单模，且 \(JM=M\) 蕴涵 \(M=0\)。
\end{proposition}
\begin{proof}
% 跨卷目标：b2:prop:finite-algebra-semisimple-quotient（第二卷命题15.1.3）。
取左正则模 \(A\) 的有限合成列 \(0=A_0\subset A_1\subset\cdots\subset A_\ell=A\)。\(J\) 零化每个简单商 \(A_i/A_{i-1}\)，所以 \(JA_i\subseteq A_{i-1}\)。逐次相乘便有 \(J^rA\subseteq A_{\ell-r}\) 对 \(0\le r\le\ell\) 成立，特别地 \(J^\ell A\subseteq A_0=0\)，即 \(J^\ell=0\)。有限维代数的半单商结构由第二卷命题15.1.3及该节建立的半单商分解给出：\(A/J\) 为有限个除环上矩阵代数的直积。每个根层被 \(J\) 零化，是 \(A/J\) 模，因而半单。最后若 \(JM=M\)，则 \(M=J^\ell M=0\)。
\end{proof}
上述最后一步是这里需要的 Nakayama 引理\footnote{中山正（假名：なかやま ただし，罗马音：Tadashi Nakayama，1912-1964），日本数学家，研究环论和群表示论。}。它也说明为什么顶不能为零：任何非零有限维模都有非零的顶。有限维模按不可分解项的直和分解及其唯一性可用第二卷定理12.4.2；该定理的有限长度条件在此由有限维性满足。不可分解模的各根层则可能包含多种简单模。

\ProseParagraphBreak
对任意有限群 \(G\) 和域 \(k\)，平凡一维模 \(k\) 都是简单左 \(kG\)-模，而 \(a=\sum_g a_gg\) 在它上面的作用是乘以 \(\varepsilon(a)=\sum_g a_g\)。因此
\[
J(kG)=\bigcap_{S\;\text{简单}}\operatorname{Ann}_{kG}(S)
\subseteq\operatorname{Ann}_{kG}(k)=\ker\varepsilon.
\]
增广理想只须零化平凡模，根却须零化全部简单模，所以这条包含关系未必是等号。下面的有限 \(p\)-群情形中，唯一的简单类型就是平凡模，因而两个理想相等。
\begin{proposition}\label{b5:prop:p-group-local}
设 \(p\) 为素数，\(P\) 为有限 \(p\)-群，\(k\) 为特征 \(p\) 的域，\(\varepsilon:kP\to k\) 为增广映射 \(\sum_g a_gg\mapsto\sum_g a_g\)。唯一的简单左 \(kP\)-模是平凡一维模 \(k\)，且
\[
J(kP)=\ker\varepsilon,\qquad kP/J(kP)\cong k.
\]
因此 \(kP\) 是局部代数。
\end{proposition}
\begin{proof}
按 \(|P|\) 归纳。\(P=1\) 时，简单 \(k\)-向量空间只有一维模。若 \(P\ne1\)，第一卷定理4.2.5给出 \(p\,\bigm\vert\,|Z(P)|\)：类方程中各非中心类的大小都被 \(p\) 整除。取中心中一个阶为 \(p^a>1\) 的元，再取其 \(p^{a-1}\) 次幂，便得阶恰为 \(p\) 的元 \(z\)。在简单模 \(S\) 上令 \(F(v)=(z-1)v\)。对 \(g\in P\)、\(v\in S\)，中心性给出
\[
gF(v)=g(zv-v)=z(gv)-gv=F(gv).
\]
因此 \(F\) 属于模自同态环 \(\operatorname{End}_{kP}(S)\)，而不只是一个线性算子。又有 \(F^p=0\)，因为特征为 \(p\) 且 \((z-1)^p=z^p-1=0\)。由\cref{b5:lem:schur}，简单模的每个非零模自同态都可逆；幂零算子不能可逆，所以 \(F=0\)。于是 \(z\) 平凡作用，\(S\) 的作用通过较小的 \(P/\langle z\rangle\) 给出，归纳使它平凡。平凡作用下简单性又迫使维数为一。根是这个唯一简单模的零化子，恰为增广理想。模掉根是域，所以非单位元恰为根，代数局部。
\end{proof}
中心性在上述证明中不可少。取特征二的域 \(k\)，\cref{b5:exm:s3-two-blocks}将构造 \(S_3=\langle r,s\rangle\) 的二维简单模，其作用矩阵为
\[
R=\rho(r)=\begin{pmatrix}0&1\\1&1\end{pmatrix},\qquad
S=\rho(s)=\begin{pmatrix}0&1\\1&0\end{pmatrix}.
\]
换位 \(s\) 阶为二，且 \((S-I)^2=0\)，但 \(S-I\ne0\)。在这个模上直接计算得到
\[
(S-I)R-R(S-I)=\begin{pmatrix}0&1\\1&0\end{pmatrix}\ne0,
\]
所以 \(S-I\) 不属于模自同态环，不能对它应用 Schur 引理。非中心群元素的作用可能恰好与其他作用交换，但其阶为 \(p\) 所给出的幂零性本身不保证这一点。

\ProseParagraphBreak
例如一个非循环 \(p\) 群仍然只有一个简单模，但其不可分解模远不止一种。简单类型的数目很少，并不意味着扩张和直和分解的分类也很短。

\subsection{\texorpdfstring{循环 \(p\) 群的基本例子}{循环 p 群的基本例子}}
\label{b5:subsec:0801:03}

\begin{theorem}\label{b5:thm:cyclic-p-modules}
设 \(p\) 为素数，\(k\) 为特征 \(p\) 的域，\(r\ge1\) 为整数，\(G=C_{p^r}=\langle s\rangle\)，并记 \(m=p^r\)。则
\[
kG\cong k[t]/(t^m),\qquad t\longmapsto s-1.
\]
全部有限维不可分解左 \(kG\)-模的同构类代表恰为
\[
M_d=k[t]/(t^d),\qquad 1\le d\le m.
\]
每个有限维左 \(kG\)-模唯一分解成这些模的直和，唯一性指各项同构类的重数唯一。记 \(J=J(kG)\)，则对 \(i\ge0\) 及 \(1\le a,b,d\le m\) 有
\[
J^iM_d=t^iM_d,\qquad
\operatorname{soc}(M_d)=kt^{d-1},\qquad
\dim_k\operatorname{Hom}_{kG}(M_a,M_b)=\min(a,b).
\]
其中\term[jizuo]{基座}[socle] \(\operatorname{soc}(M_d)\) 是全部简单子模之和。
\end{theorem}
\symidx{\(\operatorname{soc}(M_d)\)：模的基座，即全部简单子模之和}
\begin{proof}
特征 \(p\) 中重复使用 \((x-1)^p=x^p-1\)，得到 \((s-1)^m=0\)。由 \(s=1+t\) 可知所给代数同态满射，两侧维数都是 \(m\)，故同构。模结构因此由满足 \(T^m=0\) 的线性算子 \(T\) 决定。幂零算子的 Jordan 链分解，或第二卷定理4.3.2在 \(k[t]\) 上的应用，给出模为 \(k[t]/(t^d)\) 的直和，且 \(d\le m\)。这里多项式只有不可约因子 \(t\)，所以不要求 \(k\) 代数闭。
\(M_d\) 的自同态由一的像决定，环为 \(k[t]/(t^d)\)。这个环的常数项非零元素可用有限几何级数求逆，常数项零元素幂零，故它局部，没有非平凡幂等元。因此 \(M_d\) 不可分解；不同 \(d\) 的维数不同。任一分解中，大小至少为 \(j\) 的块数等于
\[
\dim\ker T^j-\dim\ker T^{j-1},
\]
所以块的重数唯一。根层由 \(J=(t)\) 直接给出；简单子模被 \(t\) 零化，而 \(\ker(t:M_d\to M_d)=kt^{d-1}\)，得到基座。最后，\(M_a\to M_b\) 的同态由一个满足 \(t^a v=0\) 的 \(v\in M_b\) 决定，这个核的维数为 \(\min(a,b)\)。
\end{proof}
例如在本定理的设定中，另假定 \(m=p^r\ge4\)，令 \(A=k[t]/(t^m)\)。此时 \(M_2\) 与 \(M_4\) 都是所列的 \(A\)-模。若 \(f:M_2\to M_4\) 是模同态，则必须有 \(t^2f(1)=0\)，所以
\[
f(1)\in\ker(t^2:M_4\to M_4)=kt^2\oplus kt^3.
\]
反过来，\(g:M_4\to M_2\) 所需的条件 \(t^4g(1)=0\) 对 \(M_2\) 的每个元素都成立，故
\[
g(1)\in M_2=k1\oplus kt
\]
可以任意选取。两处 Hom 空间的维数都是二，但前一方向的像必须落在较长链的末端 \(t^2M_4\)，反向则允许生成元映到短链中的任意元素。例如 \(f(1)=t^2\) 给出单射，而 \(g(1)=1\) 给出满射；维数公式的对称性没有抹去这两种方向的区别。

\ProseParagraphBreak
这些模的全部子模也成一条链。因为 \(M_d\) 是 \(k[t]\) 的循环商模，子模对应包含 \((t^d)\) 的理想，恰为 \((t^j)\)、\(0\le j\le d\)。例如 \(M_3\) 的根层各为一个平凡模，却不是三个平凡模的直和；其不动空间只有一维，后者有三维。

% !TeX root = ../../Book5.tex
% 纲要标题：### 8.2 投射与块
% 纲要位置：工作记录/计划纲要.md，第 3989 行。

\section{投射与块}
\label{b5:sec:0802}


% 纲要范围：
%
% * 不可分解投射模
% * 投射覆盖与 Cartan 矩阵
% * 块分解的基本概念
% * 三角矩阵代数的投射覆盖、简单顶与非分裂扩张；由正规子群平均构造 S3 在特征二下的块
%

非半单模往往不能拆成简单模，但仍可用投射模覆盖它。不可分解投射模由简单顶来区分，Cartan 矩阵记录它们内部的合成因子；块则把整个代数分成互不联系的部分。因而投射覆盖控制单个模的生成，Cartan 矩阵记录合成重数，块分解则作用于整个模范畴；小群的计算将展示三者如何配合。

\subsection{不可分解投射模}
\label{b5:subsec:0802:01}

设 \(k\) 为域，\(A\) 为有限维含幺 \(k\) 代数。若 \(P\) 为有限维左 \(A\)-模，且每个到 \(P\) 的满射都分裂，便称 \(P\) 投射；等价地，从 \(P\) 出发的同态可沿任意满射提升。有限生成情形中，这又等价于 \(P\) 是有限自由模的直和因子。将满射 \(A^n\to P\) 分裂即可得到后一条件，反向的提升可先在自由模上逐基选取。
\begin{definition}\label{b5:def:modular-projective-cover}
	设 \(k\) 为域，\(A\) 为有限维含幺 \(k\) 代数，\(P,M\) 为有限维左 \(A\)-模。若 \(\pi:P\to M\) 为满射模同态，\(P\) 投射，并且 \(\ker\pi\) 是多余子模，即 \(U+\ker\pi=P\) 对子模 \(U\subseteq P\) 只能在 \(U=P\) 时成立，就称 \(\pi\) 为 \(M\) 的\term[toushe-fugai]{投射覆盖}[projective cover]。
\end{definition}
对满射 \(\pi\)，条件 \(U+\ker\pi=P\) 等价于 \(\pi(U)=M\)：若右式成立，每个 \(x\in P\) 都能找到 \(u\in U\) 使 \(\pi(x)=\pi(u)\)，于是 \(x-u\in\ker\pi\)。因此多余核要求任何仍能覆盖目标的子模都已经是整个 \(P\)，它排除了不必参与覆盖的部分。

\begin{proposition}\label{b5:prop:projective-cover-top}
	设 \(k\) 为域，\(A\) 为有限维含幺 \(k\) 代数，\(J=J(A)\)，\(P,M\) 为有限维左 \(A\)-模，且 \(P\) 投射、\(\pi:P\to M\) 为满射模同态。\(\pi\) 是投射覆盖，当且仅当 \(\ker\pi\subseteq JP\)，也当且仅当它诱导的映射 \(P/JP\to M/JM\) 为同构。
\end{proposition}
\begin{proof}
	若 \(N=\ker\pi\) 多余，则任何极大子模都包含 \(N\)，否则与 \(N\) 的和为 \(P\)，矛盾。故 \(N\subseteq JP\)。反过来，若 \(N\subseteq JP\) 且 \(U+N=P\)，则 \(P/U=J(P/U)\)，由\cref{b5:prop:finite-radical-layers}为零。最后 \(\pi(JP)=JM\)，商映射的核为 \((N+JP)/JP\)，得到等价性。
\end{proof}
对任意域 \(k\)、正整数 \(m\) 及 \(1\le d\le m\)，令 \(A=k[t]/(t^m)\)、\(M_d=A/(t^d)\)。自然满射 \(A\to M_d\) 是投射覆盖，因为核 \((t^d)\) 包含在根 \((t)\) 中。也可直接看多余性：若 \(U+(t^d)=A\)，则 \(U\) 中有一个常数项非零的元素；它在 \(A\) 中可逆，而 \(U\) 是左理想，故 \(U=A\)。这里左正则模的自同态环为 \(A^{\mathrm{op}}=A\)，是局部环，故 \(A\) 没有非平凡直和分解。当 \(d<m\) 时，覆盖的核与目标都非零，所以该覆盖不能分裂，\(M_d\) 也就不是投射模。

\ProseParagraphBreak

特别地，固定 \(A=k[t]/(t^4)\)，模 \(M_1,M_2,M_3\) 的维数分别为 \(1,2,3\)，但它们的顶都同构于 \(k\)，投射覆盖的源也都是同一个 \(A\)。三个自然覆盖的核分别为
\[
	(t),\qquad (t^2),\qquad (t^3),
\]
维数分别为 \(3,2,1\)。因此，顶相同并不使目标模或覆盖的核相同；在这个例子中，顶只说明需要一个生成元，关系 \(t^d=0\) 才决定商模的长度。

\ProseParagraphBreak

对前面一般的 \(A=k[t]/(t^m)\)、\(M_d=A/(t^d)\)，投射模 \(A\oplus A\) 上的满射
\[
	f:A\oplus A\longrightarrow M_d,\qquad f(a,b)=a+(t^d)
\]
完全忽略第二项。它的核为 \((t^d)\oplus A\)，取真子模 \(U=A\oplus0\) 就有
\[
	U+\ker f=A\oplus A.
\]
所以这个满射不是投射覆盖。目标只需要第一项参与覆盖，多余核条件恰好不允许保留无用的第二个投射直和项。

\subsection{投射覆盖与 Cartan 矩阵}
\label{b5:subsec:0802:02}

\begin{theorem}\label{b5:thm:split-projective-covers}
	设 \(k\) 为域，\(A\) 为有限维含幺 \(k\) 代数，且对正整数 \(s,n_1,\ldots,n_s\) 有
	\[
		A/J(A)\cong\prod_{i=1}^sM_{n_i}(k).
	\]
	取简单左 \(A\)-模的同构类代表 \(S_1,\ldots,S_s\)，其中 \(\dim_kS_i=n_i\)，记 \([M:S_i]\) 为 \(S_i\) 在有限维左 \(A\)-模 \(M\) 的合成列中的重数。则每个 \(S_i\) 有不可分解投射覆盖 \(P_i\to S_i\)，且作为左 \(A\)-模有
	\[
		A\cong\bigoplus_iP_i^{\oplus n_i},\qquad \dim_k\operatorname{Hom}_A(P_i,M)=[M:S_i]
	\]
	对任意有限维左 \(A\) 模 \(M\) 成立。全部有限维投射模均为这些 \(P_i\) 的直和，各 \(P_i\) 的类型由顶 \(S_i\) 唯一确定。
\end{theorem}
\begin{proof}
	记 \(J=J(A)\)。利用根的幂零性及第二卷定理15.3.4，把商中各矩阵块的对角矩阵单位提升为一组完整正交幂等元 \(e_{ij}\)、\(1\le j\le n_i\)。令 \(P_{ij}=Ae_{ij}\)。它们是正则模的直和因子，顶为简单列模 \(S_i\)，且自然商的核为 \(Je_{ij}\)，由\cref{b5:prop:projective-cover-top}为覆盖。顶简单也使 \(P_{ij}\) 不可分解：若它分成两个非零项，两项的顶都非零，会分解这个简单顶。 同一个简单模 \(S\) 的任意两个投射覆盖 \(\pi:P\to S\)、\(\pi':P'\to S\)，存在与覆盖映射相容的同构 \(P\to P'\)。先由\cref{b5:prop:projective-cover-top}得到 \(\ker\pi\subseteq JP\)；另一方面，\(J\) 零化 \(S\)，故 \(JP\subseteq\ker\pi\)。因此 \(\ker\pi=JP\)，同样有 \(\ker\pi'=JP'\)，覆盖映射分别给出 \(P/JP\cong S\)、\(P'/JP'\cong S\)。投射性给出相互提升 \(u:P\to P'\)、\(v:P'\to P\)，满足 \(\pi'u=\pi\)、\(\pi v=\pi'\)。由 \(\pi vu=\pi\)，复合 \(vu\) 在 \(P/JP\) 上诱导恒等。于是 \(P=vu(P)+JP\)，商 \(P/vu(P)\) 等于其根，由\cref{b5:prop:finite-radical-layers}中的 Nakayama 结论为零。因此 \(vu\) 满射，有限维性使它可逆。同理 \(uv\) 可逆，故 \(u\) 为同构，且已有 \(\pi'u=\pi\)。这使 \(P_{ij}\) 的同构类型只依赖 \(i\)，记为 \(P_i\)，正交分解给出正则模公式。

	一般 \(M\) 的顶可写成 \(M/JM\cong\bigoplus_iS_i^{\oplus m_i}\)。令 \(Q=\bigoplus_iP_i^{\oplus m_i}\)，选取同构 \(\bar\pi:Q/JQ\xrightarrow{\sim}M/JM\)。\(Q\) 的投射性使复合 \(Q\to Q/JQ\xrightarrow{\bar\pi}M/JM\) 沿商映射 \(M\to M/JM\) 提升为 \(\pi:Q\to M\)。于是 \(M=\pi(Q)+JM\)，余核 \(C=M/\pi(Q)\) 满足 \(JC=C\)，由同一 Nakayama 结论得到 \(C=0\)。诱导的顶映射为 \(\bar\pi\)，故 \(\ker\pi\subseteq JQ\)，由\cref{b5:prop:projective-cover-top}得到 \(M\) 的投射覆盖。若 \(M\) 投射，这个覆盖分裂，而多余的核在分裂中只能为零，所以 \(M\) 就是这些 \(P_i\) 的直和。 最后 \(\operatorname{Hom}_A(P_i,S_j)\) 的维数为 \(\delta_{ij}\)，因为所有映射都通过简单顶，再用分裂假设下的 Schur 引理。函子 \(\operatorname{Hom}_A(P_i,-)\) 正合，沿 \(M\) 的合成列逐步取维数，便得到重数公式。
\end{proof}
这个证明也给出了任意有限维左 \(A\)-模的投射覆盖：若其顶为 \(\bigoplus_iS_i^{\oplus m_i}\)，覆盖的源就是 \(\bigoplus_iP_i^{\oplus m_i}\)。因而顶同构的模有同构的覆盖源。覆盖映射的核仍可不同，前面的 \(k[t]/(t^4)\) 例子正是这种情形。

\ProseParagraphBreak

重数公式中的 \(P_i\) 不能换成 \(S_i\)。简单模到 \(M\) 的同态，其像为零或简单子模，因而落在基座 \(\operatorname{soc}M\) 中；在本定理的分裂假设下，Schur 引理给出
\[
	\dim_k\operatorname{Hom}_A(S_i,M) =[\operatorname{soc}M:S_i].
\]
例如取 \(A=k[t]/(t^3)\)、\(S=A/(t)\)、\(P=A\) 及 \(M=A/(t^2)\)。一个同态 \(S\to M\) 由被 \(t\) 零化的向量决定，所以
\[
	\operatorname{Hom}_A(S,M)\cong k\bar t,\qquad \dim_k\operatorname{Hom}_A(S,M)=1.
\]
但 \(0\subset k\bar t\subset M\) 的两个合成因子都是 \(S\)，而在 \(1\) 处取值给出 \(\operatorname{Hom}_A(P,M)\cong M\)，其维数为 \(2\)。从 \(S\) 出发只数到基座中的那一层；从投射覆盖 \(P\) 出发的 Hom 函子正合，才会数到全部合成因子。
\begin{definition}\label{b5:def:modular-cartan-matrix}
	设 \(k\) 为域，\(A\) 为有限维含幺 \(k\) 代数，且 \(A/J(A)\cong\prod_{i=1}^sM_{n_i}(k)\)，其中 \(s,n_i\) 为正整数。取全部简单左 \(A\)-模的不同构代表 \(S_i\)，以及它们的有限维投射覆盖 \(P_i\to S_i\)，以 \([P_j:S_i]\) 表示合成因子的重数，记
	\[
		c_{ij}=[P_j:S_i]=\dim_k\operatorname{Hom}_A(P_i,P_j).
	\]
	矩阵 \(C_A=(c_{ij})\) 称为 \(A\) 的\term[Cartan-juzhen-daishu]{Cartan 矩阵}[Cartan matrix of an algebra]。本章按“行对应简单模，列对应投射覆盖”排列；它不是后来由根系定义的另一种 Cartan 矩阵。
\end{definition}
\symidx{\(C_A=(c_{ij})\)：代数的 Cartan 矩阵与合成因子重数}
在循环 \(p^r\) 群的例子中，唯一的 \(P_1\) 是 \(kG=M_{p^r}\)，故 Cartan 矩阵为一阶矩阵 \((p^r)\)。这个数字记录合成长度，不是不可分解投射模的种类数。

\begin{example}\label{b5:exm:triangular-projective-covers}
	设 \(k\) 为域，\(A=T_2(k)\) 为二阶上三角矩阵代数。令 \(S_1,S_2\) 为一维左 \(A\)-模，其中矩阵
	\[
		\begin{pmatrix}a&b\\0&c\end{pmatrix}
	\]
	分别按标量 \(a,c\) 作用。则 \(J(A)=kE_{12}\)，\(S_1,S_2\) 是全部简单模的不同构代表。它们的投射覆盖为
	\[
		P_1=AE_{11}=kE_{11}\cong S_1,\qquad P_2=AE_{22}=kE_{12}\oplus kE_{22}.
	\]
	模 \(P_2\) 给出不分裂的短正合列
	\[
		0\longrightarrow S_1\longrightarrow P_2 \longrightarrow S_2\longrightarrow0,
	\]
	而按行对应简单模、列对应投射覆盖的约定，
	\[
		C_A=\begin{pmatrix}1&1\\0&1\end{pmatrix}.
	\]
\end{example}
\begin{proof}
	理想 \(kE_{12}\) 的平方为零，商代数为 \(k\times k\)，故它正是 \(J(A)\)，且两个简单模就是 \(S_1,S_2\)。正交幂等元 \(E_{11},E_{22}\) 将左正则模分成 \(P_1\oplus P_2\)，所以两项投射。矩阵乘法给出
	\[
		JP_1=0,\qquad JP_2=kE_{12}\cong S_1, \qquad P_2/JP_2\cong S_2.
	\]
	两个自然商映射都在顶上诱导同构，由\cref{b5:prop:projective-cover-top}是投射覆盖。 若 \(P_2\to S_2\) 有截面，\(S_2\) 中的 \(1\) 必须提升为 \(xE_{12}+E_{22}\)。但 \(E_{12}\) 在 \(S_2\) 上作用为零，却将这个提升送到 \(E_{12}\ne0\)，与截面的 \(A\)-线性矛盾。所以短正合列不分裂。\(P_1\) 只有合成因子 \(S_1\)，\(P_2\) 的合成因子为 \(S_1,S_2\) 各一个，分别给出所列矩阵的两列。特别地，\(S_2\) 不是投射模；简单顶 \(S_2\) 的覆盖中，还必须包含不能从扩张中拆出的 \(S_1\)。
\end{proof}

\subsection{块分解的基本概念}
\label{b5:subsec:0802:03}

\begin{definition}\label{b5:def:algebra-block}
	设 \(k\) 为域，\(A\) 为有限维含幺 \(k\) 代数，\(b\in Z(A)\) 为非零幂等元。若 \(b\) 不能再写成两个非零正交中心幂等元之和，就称 \(b\) 为块幂等元，称以 \(b\) 为单位的代数 \(Ab\) 为 \(A\) 的一个\term[kuai]{块}[block]。
\end{definition}
反复分解中心幂等元，维数下降使过程终止，得到
\[
	1=b_1+\cdots+b_\ell,\qquad A\cong\prod_{\nu=1}^{\ell}Ab_\nu.
\]
这个分解唯一到重排：若另有 \(1=\sum c_\mu\)，则 \(b_\nu=\sum_\mu b_\nu c_\mu\) 是正交中心幂等分解，本原性使恰有一项非零，再对 \(c_\mu\) 同样分析，就使两组块对应相等。任何模都有 \(M=\bigoplus_\nu b_\nu M\)，不可分解模只属于其中一个块。不同块之间的模同态为零，扩张也分裂，因为对扩张中项分别乘各 \(b_\nu\) 即给出分解。

\ProseParagraphBreak

非零幂等元若不能再写成两个非零正交幂等元之和，就称为本原幂等元；这里不要求它属于中心。它与本原中心幂等元描述的分解并不相同。在 \(A=M_2(k)\) 中，
\[
	1=E_{11}+E_{22},\qquad A=AE_{11}\oplus AE_{22}
\]
分解的是左正则模。两项都是与列向量模 \(k^2\) 同构的简单模，因而不可分解，\(E_{11},E_{22}\) 均为本原幂等元；但它们都不在中心。中心为 \(kI\)，只有零与一两个中心幂等元，所以这个代数只有一个块。

\ProseParagraphBreak

非中心幂等元的像也不总是子模。例如在自然左模 \(M=k^2\) 中，\(E_{11}M=ke_1\)，却有 \(E_{21}e_1=e_2\notin ke_1\)，所以它不被 \(A\) 保持。中心幂等元 \(b\) 则对任意 \(a\in A,m\in M\) 满足 \(a(bm)=b(am)\)，故 \(bM\) 总是子模。这解释了为什么块可以直接分解每个模，而一般本原幂等元在本节用来分解正则模并构造投射覆盖。
\begin{example}\label{b5:exm:s3-two-blocks}
	设 \(k\) 为特征二的域，\(G=S_3=\langle r,s\mid r^3=s^2=1,\ srs=r^{-1}\rangle\)，\(N=\langle r\rangle\)。令
	\[
		e=\frac1{|N|}\sum_{n\in N}n=1+r+r^2, \qquad f=1-e.
	\]
	则 \(e,f\) 是两个块幂等元，且
	\[
		kG\cong kGe\times kGf\cong kC_2\times M_2(k).
	\]
	对任意有限维左 \(kG\)-模 \(V\)，有 \(eV=V^N\)。全部简单模的不同构代表为平凡模 \(S_1\) 和一个二维模 \(S_2\)，它们的投射覆盖分别为 \(P_1=kGe\)、\(P_2=S_2\)，并有
	\[
		C_{kS_3}=\begin{pmatrix}2&0\\0&1\end{pmatrix},
		\qquad J(kG)=k\!\left(\sum_{g\in G}g\right).
	\]
	特别地，\(\dim_kJ(kG)=1\)，而增广理想 \(I=\ker\varepsilon\) 的维数为 \(5\)，故 \(J(kG)\subsetneq I\)。
\end{example}
\begin{proof}
	正规子群 \(N\) 的阶为三，在 \(k\) 中可逆。其平均元满足 \(ne=e\) 对所有 \(n\in N\) 成立，故 \(e^2=e\)。正规性又使共轭置换求和中的元素，所以 \(e\in Z(kG)\)，\(f\) 也是中心幂等元。对模 \(V\)，每个 \(ev\) 都被 \(N\) 固定；反过来，若 \(v\in V^N\)，则 \(ev=v\)。因此 \(e\) 正是取出 \(N\)-固定向量的投影，\(eV=V^N\)。

	两个陪集 \(N,sN\) 的元素各自汇总为 \(e,se\)，这两元素线性无关且张成 \(kGe\)。乘法满足 \((se)^2=e\)，所以 \(kGe\cong k(G/N)\cong kC_2\)。特征二时 \(x=e+se\) 满足 \(x^2=0\)，且 \(kGe/kx\cong k\)，故这个二维代数是局部代数，根为 \(kx\)。

	另一项用二维表示来识别。取
	\[
		R=\begin{pmatrix}0&1\\1&1\end{pmatrix},\qquad
		S=\begin{pmatrix}0&1\\1&0\end{pmatrix}.
	\]
	有 \(R^3=S^2=I\)、\(SRS=R^{-1}\)，故它们定义表示映射 \(\rho:kG\to M_2(k)\)。又有 \(I+R+R^2=0\)，所以 \(\rho(e)=0\)、\(\rho(f)=I\)，而 \(f,rf,sf,rsf\) 的像分别为 \(I,R,S,RS\)。这四个矩阵线性无关：若
	\[
		aI+bR+cS+dRS
		=\begin{pmatrix}a+d&b+c\\b+c+d&a+b+d\end{pmatrix}=0,
	\]
	比较四个条目便得 \(d=a=b=c=0\)。因此 \(\rho(kGf)\) 的维数为四，等于 \(\dim_kM_2(k)\)，限制映射 \(\rho:kGf\to M_2(k)\) 满射。中心幂等分解给出 \(kG=kGe\oplus kGf\)，前面已由基 \(e,se\) 得到 \(\dim_k(kGe)=2\)，故
	\[
		\dim_k(kGf)=6-\dim_k(kGe)=4.
	\]
	满射两侧维数相等，所以 \(kGf\cong M_2(k)\)。局部代数 \(kGe\) 与矩阵代数 \(kGf\) 都没有非平凡中心幂等元，因而这确为两个块。第一块只有平凡简单模 \(S_1\)；第二块只有自然列模 \(S_2\)。在 \(S_1\) 上 \(e\) 作用为恒等，在 \(S_2\) 上 \(I+R+R^2=0\)，故 \(e\) 作用为零。两者分别具有全部与零个 \(N\)-固定向量。

	\(P_1=kGe\) 的根 \(kx\) 与顶都同构于 \(S_1\)，自然商 \(P_1\to S_1\) 由\cref{b5:prop:projective-cover-top}是投射覆盖。第二块是矩阵代数，其简单列模是左正则模的直和因子，故 \(S_2\) 本身投射，恒等映射就是它的投射覆盖。因此两个覆盖的合成因子分别为两个 \(S_1\) 与一个 \(S_2\)，得到所列 Cartan 矩阵。

	最后，直积代数的根等于两项根的直积，而矩阵块的根为零，故
	\[
		J(kG)=k(e+se)=k\!\left(\sum_{g\in G}g\right).
	\]
	增广满足 \(\varepsilon(e)=\varepsilon(se)=1\)、\(\varepsilon(f)=0\)，所以
	\[
		I=k(e+se)\oplus kGf, \qquad \dim_k I=1+4=5.
	\]
	这也直接显示根与增广理想的区别：非零矩阵块整个包含在 \(I\) 中，却对根没有贡献。
\end{proof}
这里 \(2\,\bigm\vert\,|S_3|\)，却仍有一个半单矩阵块。模特征使整个群代数不半单，并不要求每个块都非半单。

% !TeX root = ../../Book5.tex
% 纲要标题：### 8.3 Brauer 特征标
% 纲要位置：工作记录/计划纲要.md，第 3995 行。

\section{Brauer 特征标}
\label{b5:sec:0803}


% 纲要范围：
%
% * p-正则元素
% * 普通特征标与 Brauer 特征标的区别
% * 分解映射与分解数
% * 单位根提升所需的含根假设、分圆域嵌入的选择及不同整格约化；Brauer 基定理与普通正交关系的区别
%

正特征中的矩阵迹会把整数重数约化模 \(p\)，所以丢掉一部分信息。Brauer 特征标改为先提升适当的单位根，再在特征零中相加。为了让这种提升与普通表示的约化对应，必须同时说明基域、赋值环和整数格；只说“把矩阵模 \(p\)”并不足以定义分解数。

\subsection{\texorpdfstring{\(p\)-正则元素}{p-正则元素}}
\label{b5:subsec:0803:01}

\begin{definition}\label{b5:def:p-regular-element}
设 \(G\) 为有限群，\(p\) 为素数。若 \(g\in G\) 的阶与 \(p\) 互素，就称 \(g\) 为\term[p-zhengze-yuansu]{\(p\)-正则元素}[\(p\)-regular element]。全体此类元素记为 \(G_{p'}\)；其余元素称为 \(p\)-奇异元素。这个集合在共轭下稳定，但通常不是子群。
\end{definition}
\symidx{\(G_{p^{\prime}}\)：有限群的 \(p\)-正则元素集合}
设 \(k\) 的特征为 \(p\)，\(g\) 的阶为 \(m\) 且 \(p\nmid m\)。在任一表示上，\(g\) 的最小多项式整除 \(X^m-1\)；后者与导数 \(mX^{m-1}\) 互素，所以在分裂域上 \(g\) 可对角化，特征值都是阶与 \(p\) 互素的单位根。这正是定义 Brauer 特征标时选择这些元素的原因。
\begin{lemma}\label{b5:lem:modular-trace-regular-part}
设 \(G\) 为有限群，\(p\) 为素数，\(k\) 为特征 \(p\) 的域，\(V\) 为有限维左 \(kG\)-模。每个 \(g\in G\) 唯一写成 \(g=g_pg_{p'}=g_{p'}g_p\)，其中两因子都属于 \(\langle g\rangle\)，前者的阶为 \(p\) 的幂，后者的阶与 \(p\) 互素。并且
\[
\operatorname{tr}_k(g|V)=\operatorname{tr}_k(g_{p'}|V).
\]
\end{lemma}
\begin{proof}
在阶为 \(p^a m\)、\(p\nmid m\) 的循环群中，中国剩余定理将指数分成模 \(p^a\) 与模 \(m\) 的两部分，给出存在与唯一性。扩到代数闭包不改变迹。此时 \(g_{p'}\) 对角化，而 \(g_p\) 与它交换，故保持其每个特征空间。在这些空间上，\((g_p-I)^{p^a}=0\)，所以 \(g_p\) 的特征值全为一。若某特征空间上 \(g_{p'}\) 的值为 \(\lambda\)，则 \(g\) 在该空间的迹为 \(\lambda\) 乘该空间维数，与 \(g_{p'}\) 相同。相加即得。
\end{proof}
这里的等式只比较 \(k\) 中的迹，并没有恢复被模 \(p\) 丢掉的整数。比如 \(p\) 个平凡模的直和，在所有群元素上的 \(k\) 值迹都是零，却当然不是零模。

\subsection{普通特征标与 Brauer 特征标}
\label{b5:subsec:0803:02}

\begin{definition}\label{b5:def:p-modular-system}
设 \(p\) 为素数，\(\mathcal O\) 为离散赋值环，其分式域 \(K\) 的特征为零，唯一极大理想为 \((\pi)\)，剩余域 \(k=\mathcal O/(\pi)\) 的特征为 \(p\)。三元组 \((K,\mathcal O,k)\) 称为一个\term[p-mo-xitong]{\(p\)-模系统}[\(p\)-modular system]。对有限群 \(G\)，若 \(KG\) 与 \(kG\) 的半单商均为相应基域上全矩阵代数的直积，就称它为 \(G\) 的分裂 \(p\)-模系统。
\end{definition}
\symidx{\((K,\mathcal O,k)\)：\(p\)-模系统}
这里用到的环性质是：\(\mathcal O\) 是主理想整环，有限生成无挠模自由，每个非零元素等于单位乘以某个 \(\pi\) 的幂。本节把这样的系统作为给定数据。
\ProseParagraphBreak
为讨论分解矩阵，以下使用分裂 \(p\)-模系统，并进一步假设 \(K\) 含有所有 \(|G|\) 次单位根。分裂域与模系统的准备可参阅\cite[Section 9.2, Theorem 9.2.6; Section 9.4]{Webb2016FiniteGroupRepresentations}。分裂条件是这些证明的前提，任意有限域未必满足它。
\begin{lemma}\label{b5:lem:prime-to-p-root-lift}
设 \(G\) 为有限群，\((K,\mathcal O,k)\) 为 \(p\)-模系统，且 \(K\) 含有全部 \(|G|\) 次单位根。若正整数 \(m\,\bigm\vert\,|G|\) 且 \(p\nmid m\)，则约化映射给出 \(K\) 中 \(m\) 次单位根与 \(k\) 中 \(m\) 次单位根的群同构。记相应唯一提升为 \(\lambda\mapsto\widehat\lambda\)。
\end{lemma}
\begin{proof}
若 \(\zeta^m=1\)，则 \(m\,v(\zeta)=v(1)=0\)，所以 \(\zeta\) 的赋值为零，属于 \(\mathcal O^\times\)。由 \(m\,\bigm\vert\,|G|\) 及 \(K\) 含全部 \(|G|\) 次单位根的假设，\(X^m-1\) 在 \(K\) 中已经分成 \(m\) 个一次因子。这些因子的系数都在 \(\mathcal O\) 中，约化后仍给出 \(X^m-1\) 的分解。又因 \(p\nmid m\)，它在 \(k\) 中与导数 \(mX^{m-1}\) 互素，所以 \(m\) 个根的约化两两不同，并穷尽 \(k\) 中的全部根。约化保持乘法，故得到群同构。
\end{proof}
这里的根在 \(K\) 中已经存在，证明只需说明约化不会把它们合并；没有使用 \(\mathcal O\) 为 Hensel 环的假设。若删去 \(K\) 中存在全部所需单位根的条件，结论会失败。例如
\[
K=\mathbb Q,\qquad \mathcal O=\mathbb Z_{(7)},\qquad k=\mathbb F_7
\]
是一个 \(7\)-模系统。在 \(\mathbb F_7\) 中，\(2^3=4^3=1\)，所以 \(2,4\) 是非平凡三次单位根；但 \(\mathbb Q\) 中的三次单位根只有 \(1\)，因为 \(X^2+X+1\) 在 \(\mathbb Q\) 中没有根。因此它们不能提升成 \(\mathbb Z_{(7)}\) 中的三次单位根。离散赋值环的条件本身不足以保证这种提升。

\ProseParagraphBreak

记 \(E=\mathbb Q(\mu_{|G|})\subseteq K\) 为这些单位根生成的分圆域，固定一个域嵌入 \(\iota:E\to\mathbb C\)。单位根的提升 \(\widehat\lambda\) 属于 \(\mathcal O\cap E\)，相应的复单位根则是 \(\iota(\widehat\lambda)\)；所有表示和张量运算都使用同一个提升及嵌入。
\begin{definition}\label{b5:def:brauer-character}
设 \(G\) 为有限群，\(p\) 为素数，\((K,\mathcal O,k)\) 为分裂 \(p\)-模系统，\(K\) 含有全部 \(|G|\) 次单位根。令 \(E=\mathbb Q(\mu_{|G|})\subseteq K\)，并固定域嵌入 \(\iota:E\to\mathbb C\)。记 \(G_{p'}\) 为 \(p\)-正则元素集合，\(\lambda\mapsto\widehat\lambda\) 为约化映射 \(\mathcal O\to k\) 所确定的阶整除 \(|G|\) 且与 \(p\) 互素的单位根的唯一提升。对有限维左 \(kG\)-模 \(M\)，令 \(d=\dim_kM\)。对 \(g\in G_{p'}\)，将 \(g\) 在 \(M\) 上的特征值按重数写为 \(\lambda_1,\ldots,\lambda_d\)，并记提升后的和为
\[
\widetilde\varphi_M(g)=\widehat\lambda_1+\cdots+\widehat\lambda_d
\in\mathcal O\cap E.
\]
函数 \(\varphi_M:G_{p'}\to\mathbb C\)，由
\[
\varphi_M(g)=\iota\bigl(\widetilde\varphi_M(g)\bigr)
\]
确定，称为 \(M\) 的\term[Brauer-tezhengbiao]{Brauer 特征标}[Brauer character]。提升后的和约化为原来的 \(k\) 值迹，即
\[
\widetilde\varphi_M(g)\bmod\pi=\operatorname{tr}_k(g|M).
\]
其定义域只取 \(G_{p'}\)，而 \(\varphi_M\) 的值在复数域中；简单左 \(kG\)-模的 Brauer 特征标称为不可约 Brauer 特征标。
\end{definition}
\symidx{\(\varphi_M\)：模表示的 Brauer 特征标}
\begin{proposition}\label{b5:prop:brauer-character-properties}
设 \(G\) 为有限群，\(p\) 为素数，\((K,\mathcal O,k)\) 为分裂 \(p\)-模系统，\(K\) 含有全部 \(|G|\) 次单位根，并固定这些根到复数的对应。设 \(L,M,N\) 为有限维左 \(kG\)-模，\(S_i\) 为全部简单左 \(kG\)-模的不同构代表，\(\varphi_X\) 表示模 \(X\) 的 Brauer 特征标，\([M:S_i]\) 表示合成因子的重数。这些特征标在 \(p\)-正则共轭类上为常数，且对 \(p\)-正则元素 \(g\in G\) 有
\[
\varphi_M(1)=\dim_kM,\qquad
\varphi_{M\otimes N}=\varphi_M\varphi_N,\qquad
\varphi_{M^*}(g)=\varphi_M(g^{-1})=\overline{\varphi_M(g)}.
\]
对短正合列 \(0\to L\to M\to N\to0\)，有 \(\varphi_M=\varphi_L+\varphi_N\)。因而
\[
\varphi_M=\sum_i[M:S_i]\varphi_{S_i}.
\]
\end{proposition}
\begin{proof}
共轭矩阵有相同的特征值，单位元的全部特征值为一，得到前两个直接性质。张量积的特征值为所有两两乘积；单位根提升是群同态，故其提升也两两相乘，给出乘法公式。对偶把特征值取逆，而复单位根的逆为其共轭，得到对偶公式。
将一个 \(p\)-正则元素生成的循环群限制到给定正合列。因为该循环群的阶可逆，Maschke 定理使列在这个子群上分裂，所以中项的特征值多重集为两端的并。逐个提升后相加便得加法性；沿合成列归纳得到最后一式。
\end{proof}
对 \(S_3\) 与 \(p=2\)，可以把各个系数环固定为
\[
K=\mathbb Q(\zeta_3),\qquad
\mathcal O=\mathbb Z[\zeta_3]_{(2)},\qquad
k=\mathbb F_4,\qquad \pi=2,
\]
其中 \(\zeta_3=\exp(2\pi\mathrm i/3)\)。因为 \(T^2+T+1\) 在 \(\mathbb F_2[T]\) 中不可约，\(\mathbb Z[\zeta_3]/(2)\cong\mathbb F_4\)，所以 \((2)\) 是素理想，上式按此理想局部化。任意非零 \(a+b\zeta_3\) 都能提出其两个整数系数的最大共同 \(2\) 次幂；余下元素不属于 \((2)\)，在局部化中成为单位。因此 \(\mathcal O\) 是以 \(2\) 为统一化元的离散赋值环。平凡、符号和标准二维表示在 \(K\) 上给出作用同态 \(KS_3\to K\times K\times M_2(K)\)；扩域到 \(\mathbb C\) 后，由\cref{b5:tab:s3,b5:thm:complex-group-algebra}它是同构。矩阵的秩扩域不变，故原来的作用同态也是同构。又由\cref{b5:exm:s3-two-blocks}，\(kS_3\) 的半单商为 \(k\times M_2(k)\)，所以这是分裂系统；\(K\) 也含有全部六次单位根。

\ProseParagraphBreak
记 \(\alpha\in\mathbb F_4\) 为 \(\zeta_3\) 的约化。三循环 \(r=(123)\) 在上述二维 \(kS_3\)-模中的矩阵为 \(\begin{psmallmatrix}0&1\\1&1\end{psmallmatrix}\)，特征值为 \(\alpha,\alpha^2\)，所以
\[
\operatorname{tr}_k(r)=\alpha+\alpha^2=1,\qquad
\widehat\alpha+\widehat{\alpha^2}
=\zeta_3+\zeta_3^2=-1.
\]
固定的分圆域嵌入取为当前 \(K\subseteq\mathbb C\) 的包含映射，于是 Brauer 特征标在 \(r\) 处取复数值 \(-1\)。剩余域中的等式 \(\alpha+\alpha^2=1\) 不能先把右边提升成复数 \(1\)：定义提升的是两个特征值，再在特征零中相加。

\ProseParagraphBreak
同一个 \(2\)-模系统也适用于 \(C_3=\langle r\rangle\)：\(T^3-1\) 分别在 \(K\) 与 \(k\) 中分成三个不同的一次因子，所以两侧群代数都分裂。取一维 \(kC_3\)-模 \(M\)，使 \(r\) 按 \(\alpha\) 作用，则其唯一特征值的提升为 \(\widehat\alpha=\zeta_3\)。若分别采用嵌入
\[
\iota_1(\zeta_3)=\exp(2\pi\mathrm i/3),\qquad
\iota_2(\zeta_3)=\exp(4\pi\mathrm i/3),
\]
同一个模在 \(r\) 处的复值 Brauer 特征标分别为
\[
\varphi_M^{(1)}(r)=\exp(2\pi\mathrm i/3),\qquad
\varphi_M^{(2)}(r)=\exp(4\pi\mathrm i/3).
\]
第二组数值是第一组的 Galois 共轭。因此，从剩余域矩阵得到具体复数值时，嵌入的选择不能省略。这里的三个一维 \(KC_3\)-模 \(r\mapsto\zeta_3^j\)，其中 \(j=0,1,2\)，各有稳定格 \(\mathcal O\)，并分别约化为 \(r\mapsto\alpha^j\) 的简单模。每个约化模中，对应简单模的重数为一，其他重数为零；改变嵌入只改变这些特征标的复数值，不改变约化模或简单合成因子的重数。比较普通特征标与 Brauer 特征标时，两者必须采用同一个嵌入。

\ProseParagraphBreak
这些定义和性质可对照\cite[Section 10.1, Lemma 10.1.1; Proposition 10.1.3]{Webb2016FiniteGroupRepresentations}。Webb 在该章还统一要求两个域对 \(G\) 及其所有子群均为分裂域；本节则在 \(G\) 的分裂 \(p\)-模系统之外，显式要求 \(K\) 含全部 \(|G|\) 次单位根，再由\cref{b5:lem:prime-to-p-root-lift}取得定义所需的 \(k\) 中单位根及其提升。

\subsection{分解映射与分解数}
\label{b5:subsec:0803:03}

\begin{definition}\label{b5:def:stable-integral-lattice}
设 \(G\) 为有限群，\(p\) 为素数，\((K,\mathcal O,k)\) 为 \(p\)-模系统，\((\pi)\) 为 \(\mathcal O\) 的极大理想，\(V\) 为有限维左 \(KG\)-模。若 \(L\subseteq V\) 是被 \(G\) 保持的有限自由 \(\mathcal O\) 子模，且其一组 \(\mathcal O\) 基也是 \(V\) 的 \(K\) 基，就称 \(L\) 为 \(V\) 中的一个满的 \(G\)-稳定\term[zhenggeshu-ge]{格}[lattice]。商 \(L/\pi L\) 是一个有限维左 \(kG\)-模，称为该格的约化。
\end{definition}
\begin{theorem}[Brauer--Nesbitt]\label{b5:thm:lattice-reduction-independence}
设 \(G\) 为有限群，\(p\) 为素数，\((K,\mathcal O,k)\) 为 \(p\)-模系统，\((\pi)\) 为 \(\mathcal O\) 的极大理想，\(V\) 为有限维左 \(KG\)-模。则 \(V\) 有满的 \(G\)-稳定格；任意两个这种格 \(L_1,L_2\) 的约化 \(L_1/\pi L_1\)、\(L_2/\pi L_2\) 具有相同的简单左 \(kG\)-模合成因子及重数，但不必同构。
\end{theorem}
\begin{proof}
先任取 \(K\) 基 \(v_1,\ldots,v_d\)，令 \(L\) 为所有 \(gv_i\) 的 \(\mathcal O\) 线性生成空间。它有限生成、无挠且被 \(G\) 保持；由主理想整环上的模结构，它自由，清分母可知其秩恰为 \(d\)，所以是满格。

两满格可通约：分别把一组格基用另一组表示，给有限多个系数清去 \(\pi\) 的负幂，就能取整数 \(a,b\ge0\)，使
\[
\pi^aL_1\subseteq L_2\subseteq\pi^{-b}L_1.
\]
乘以 \(\pi^a\) 给出 \(L_1\to\pi^aL_1\) 的 \(\mathcal O G\)-模同构，并诱导其约化的同构。因此可先以 \(\pi^aL_1\) 替换第一个格，再比较两个嵌套格；以下将替换后的格仍记为 \(L_1\)，于是 \(L_1\subseteq L_2\)，且某个 \(\pi^N L_2\subseteq L_1\)。

商 \(L_2/L_1\) 是有限生成且被 \(\pi^N\) 零化的 \(\mathcal O\)-模，因而有有限长度。沿其 \(\mathcal O G\)-模合成列在两格间插入有限多个稳定满格，可归约到 \(L_2/L_1\) 简单的情形。乘以 \(\pi\) 是这个简单商的模自同态；若非零，其核为零、像为全体，故是自同构。但它的 \(N\) 次幂为零，不可能是自同构，因此 \(\pi\) 本身零化这个商。于是 \(\pi L_2\subseteq L_1\)，得到链
\[
L_2\supseteq L_1\supseteq\pi L_2\supseteq\pi L_1.
\]
这条链给出两个有限维左 \(kG\)-模的短正合列：
\[
0\longrightarrow L_1/\pi L_2
\longrightarrow L_2/\pi L_2
\longrightarrow L_2/L_1\longrightarrow0,
\]
\[
0\longrightarrow\pi L_2/\pi L_1
\longrightarrow L_1/\pi L_1
\longrightarrow L_1/\pi L_2\longrightarrow0.
\]
第一列的单射来自 \(L_1\subseteq L_2\)，此时已经模掉共同的 \(\pi L_2\)；它不是把 \(L_1/\pi L_1\) 直接嵌入 \(L_2/\pi L_2\)。两个约化的合成因子共有 \(L_1/\pi L_2\) 这一部分，另外两项由
\[
L_2/L_1\longrightarrow\pi L_2/\pi L_1,\qquad
x+L_1\longmapsto\pi x+\pi L_1
\]
识别。此映射良定义且满射；若 \(\pi x\in\pi L_1\)，存在 \(y\in L_1\) 使 \(\pi(x-y)=0\)，格的无挠性给出 \(x=y\)，故它也单射。沿两条正合列把合成因子相加，两约化的多重集相同；沿插入的有限链传递即得一般结论。
\end{proof}
这条格无关性称为 Brauer--Nesbitt 定理\footnote{内斯比特（Cecil James Nesbitt，1912-2001），美国数学家，与布劳尔合作研究模表示论；布劳尔已见本卷前注。}的格约化形式，见\cite[Theorem 9.4.8]{Webb2016FiniteGroupRepresentations}。必须保留“合成因子”这一表述。取
\[
G=C_2=\langle s\rangle,\qquad K=\mathbb Q,\qquad
\mathcal O=\mathbb Z_{(2)},\qquad k=\mathbb F_2,\qquad V=KG.
\]
正则格 \(L_0=\mathcal O1\oplus\mathcal Os\) 在基 \(1,s\) 下的作用矩阵为
\[
\rho_{L_0}(s)=\begin{pmatrix}0&1\\1&0\end{pmatrix}.
\]
约化后，\((\rho_{L_0}(s)-I)^2=0\)，但 \(\rho_{L_0}(s)-I\ne0\)，所以它有一条长度二的幂零链。由\cref{b5:thm:cyclic-p-modules}，\(L_0/2L_0\cong k[t]/(t^2)=M_2\) 是不可分解模。

\ProseParagraphBreak
再取 \(u=1+s\)、\(v=1-s\)，令 \(L_1=\mathcal Ou\oplus\mathcal Ov\)。基变换矩阵
\[
P=\begin{pmatrix}1&1\\1&-1\end{pmatrix},\qquad \det P=-2
\]
在 \(K\) 上可逆，故 \(L_1\) 也是满格。由于 \(su=u\)、\(sv=-v\)，它是稳定格，并且在基 \(u,v\) 下有
\[
\rho_{L_1}(s)=\begin{pmatrix}1&0\\0&-1\end{pmatrix}.
\]
约化后这个矩阵变成 \(I\)，所以 \(L_1/2L_1\cong k\oplus k\)。两种约化都有两个平凡简单合成因子，但一个不可分解，另一个是实际直和。

\ProseParagraphBreak
虽然 \(L_1\subsetneq L_0\)，包含映射诱导的 \(L_1/2L_1\to L_0/2L_0\) 并不单射：源中的基向量 \(u+2L_1\)、\(v+2L_1\) 线性无关，而它们的像都等于 \(1+s+2L_0\)。该映射的矩阵是约化后的 \(P\)，秩只有一。把特征零作用对角化需要在 \(P^{-1}\) 中除以 \(2\)，这个基变换不能在 \(\mathcal O\) 上进行，也不能带到约化模上。这说明格包含不能直接给出约化模包含；定理证明中须模掉共同的 \(\pi L_2\)，才能得到所用的第一条短正合列。
\begin{definition}\label{b5:def:decomposition-map}
设 \(G\) 为有限群，\(p\) 为素数，\((K,\mathcal O,k)\) 为分裂 \(p\)-模系统，\((\pi)\) 为 \(\mathcal O\) 的极大理想。令 \(G_0(kG)\) 为以有限维左 \(kG\)-模的同构类生成、并对每个短正合列 \(0\to L\to M\to N\to0\) 加入 \([M]=[L]+[N]\) 的交换群。Jordan--Hölder 定理\footnote{若尔当（Camille Jordan，1838-1922），法国数学家，研究群论与代数方程；赫尔德（Otto Ludwig Hölder，1859-1937），德国数学家，研究群论与分析。}使它是以全部简单左 \(kG\)-模的不同构代表类 \([S_i]\) 为基的自由交换群。令 \(R_K(G)\) 为有限维左 \(KG\)-模的表示群。对每个这样的模 \(V\)，选择其中的满的 \(G\)-稳定格 \(L\)，则对应
\[
d:R_K(G)\longrightarrow G_0(kG),\qquad
[V]\longmapsto[L/\pi L]
\]
称为\term[fenjie-yingshe]{分解映射}[decomposition map]。取全部简单左 \(KG\)-模的不同构代表 \(V_\alpha\)，以及其中的满的 \(G\)-稳定格 \(L_\alpha\)，定义
\[
d_{\alpha i}=[L_\alpha/\pi L_\alpha:S_i].
\]
这些非负整数称为\term[fenjieshu]{分解数}[decomposition numbers]，矩阵 \(D=(d_{\alpha i})\) 称为分解矩阵，行对应特征零简单模，列对应特征 \(p\) 简单模。
\end{definition}
\symidx{\(R_K(G),G_0(kG)\)：特征零表示群与模表示的 Grothendieck 群}
\symidx{\(d:R_K(G)\to G_0(kG)\)：分解映射}
\symidx{\(d_{\alpha i},D\)：分解数与分解矩阵}
前一定理保证它与格的选择无关。分解数只取决于约化模的简单合成因子，不依赖分圆域到复数域的嵌入；上面 \(C_3\) 算例中按对应的一维模排列的分解矩阵就是 \(I_3\)，两种复嵌入都给出这些相同的整数。直和格给出加法性。两侧分别以 \(K\) 上和 \(k\) 上的张量积为乘法；域上张量积在每个变量正合，所以这些乘法与短正合列的关系相容。若 \(L,L'\) 分别为 \(V,V'\) 中的稳定满格，则 \(L\otimes_{\mathcal O}L'\) 是 \(V\otimes_KV'\) 中的稳定满格，且
\[
(L\otimes_{\mathcal O}L')/\pi(L\otimes_{\mathcal O}L')
\cong (L/\pi L)\otimes_k(L'/\pi L')
\]
为 \(kG\)-模同构。因此 \(d\) 保持乘法，又将平凡模的类送到平凡模的类，故是环同态。这里 \(G_0\) 使用的是短正合关系，不是把非分裂扩张误认成实际直和。
\begin{proposition}\label{b5:prop:ordinary-brauer-restriction}
设 \(G\) 为有限群，\(p\) 为素数，\((K,\mathcal O,k)\) 为分裂 \(p\)-模系统，\((\pi)\) 为 \(\mathcal O\) 的极大理想，\(K\) 含有全部 \(|G|\) 次单位根，并固定这些根到复数的对应。设 \(V\) 为有限维左 \(KG\)-模，\(L\) 为其中的满的 \(G\)-稳定格，\(\chi_V\) 为按此对应得到的普通特征标，\(G_{p'}\) 为 \(p\)-正则元素集合，\(\varphi_M\) 表示左 \(kG\)-模 \(M\) 的 Brauer 特征标。取简单左 \(KG\)-模代表 \(V_\alpha\) 及其中的稳定满格 \(L_\alpha\)，记其普通特征标为 \(\chi_\alpha\)，取简单左 \(kG\)-模代表 \(S_i\)，记 \(d_{\alpha i}=[L_\alpha/\pi L_\alpha:S_i]\)。则
\[
\chi_V|_{G_{p'}}=\varphi_{L/\pi L},\qquad
\chi_\alpha|_{G_{p'}}=\sum_i d_{\alpha i}\varphi_{S_i}.
\]
\end{proposition}
\begin{proof}
设 \(g\in G_{p'}\) 的阶为 \(m\)。它在 \(L\) 上有系数位于 \(\mathcal O\) 的特征多项式，约化后就是在 \(L/\pi L\) 上的特征多项式。其根均为 \(m\) 次单位根，且这些单位根的约化互异；因而每个根及其重数在约化后对应到唯一的根及相同重数。把根再提升回来，其和恢复 \(\operatorname{tr}_K(g|V)\)；再经固定的分圆域嵌入，就得到 \(\varphi_{L/\pi L}(g)=\chi_V(g)\)。第二式由 Brauer 特征标对合成因子的加法性得到。
\end{proof}
例如在上面的 \(S_3\) 分裂 \(2\)-模系统中，取普通符号表示及其一维稳定格 \(L=\mathcal O\)。换位 \(s\) 在特征零上按 \(-1\) 作用，故 \(\chi_{\mathrm{sgn}}(s)=-1\)；约化到特征二后，\(-1=1\)，符号模与平凡模相同。若把约化矩阵在 \(s\) 上的唯一特征值 \(1\) 提升回去，所得却是 \(1\)，不能恢复原来的普通特征标值 \(-1\)。换位是 \(2\)-奇异元素，正好不在比较公式的定义域中；三循环则是 \(2\)-正则元素，符号表示在它上面原本就按 \(1\) 作用，约化与提升都保留这个值。

\ProseParagraphBreak
一般地，在\cref{b5:lem:modular-trace-regular-part}的证明中，\(g_p\) 在 \(g_{p'}\) 的每个特征空间上都是幂幺算子。因此 \(g=g_pg_{p'}\) 与 \(g_{p'}\) 在特征 \(p\) 下有相同的特征值多重集。若继续在 \(p\)-奇异元素上逐个提升这些特征值，只会得到其 \(p\)-正则部分已经给出的数值，无法补回约化所丢失的 \(p\) 部分。限制到 \(G_{p'}\) 才能保证与普通特征标的值对应。

\begin{example}\label{b5:exm:brauer-s3-characteristic-three}
取 \(G=S_3=\langle r,s\mid r^3=s^2=1,\ srs=r^{-1}\rangle\)，其中 \(r=(123)\)、\(s=(12)\)，并取满足本节单位根条件的分裂 \(3\)-模系统 \((K,\mathcal O,k)\)。先确定 \(kS_3\) 的简单模。对任意非零简单模 \(S\)，算子 \(r-I\) 满足 \((r-I)^3=0\)，故其核 \(S^{\langle r\rangle}\) 非零。因为 \(\langle r\rangle\) 正规，这个不动空间被 \(S_3\) 保持：对 \(x\in S^{\langle r\rangle}\)、\(g\in S_3\) 及 \(h\in\langle r\rangle\)，有
\[
h(gx)=g(g^{-1}hg)x=gx.
\]
简单性使这个不动空间等于 \(S\)，所以 \(S\) 通过 \(S_3/\langle r\rangle\cong C_2\) 作用。特征三下 \(T^2-1=(T-1)(T+1)\) 的两个根不同，\(kC_2\cong k\times k\)，故恰有两个一维简单模 \(S_+,S_-\)：\(r\) 都平凡作用，\(s\) 分别按 \(1,-1\) 作用。

两类 \(3\)-正则元素由 \(1,s\) 代表，两个不可约 Brauer 特征标按此顺序为 \(\varphi_+=(1,1)\)、\(\varphi_-=(1,-1)\)。若仿照普通特征标内积，定义
\[
(f,h)_0=\frac1{6}\sum_{g\in(S_3)_{3'}}\overline{f(g)}h(g),
\]
则因 \((S_3)_{3'}\) 包含单位元及三个换位，有
\[
(\varphi_+,\varphi_+)_0=(\varphi_-,\varphi_-)_0=\frac{1+3}{6}=\frac23,
\qquad
(\varphi_+,\varphi_-)_0=\frac{1-3}{6}=-\frac13.
\]
所以这两个函数既没有归一化，也不正交。它们仍线性无关，因为值矩阵的行列式为
\[
\det\begin{pmatrix}1&1\\1&-1\end{pmatrix}=-2,
\]
它在特征三的域中非零，提升到特征零后也非零。生成性可以从普通特征标的限制看到：平凡、符号和标准二维表示在这两类上的值依次为
\[
\chi_{\mathrm{triv}}|_{G_{3'}}=\varphi_+,\qquad
\chi_{\mathrm{sgn}}|_{G_{3'}}=\varphi_-,\qquad
\chi_{\mathrm{std}}|_{G_{3'}}=(2,0)=\varphi_++\varphi_-.
\]
前两式来自一维作用；标准表示是三点置换表示去掉常数直线，单位元上维数为二，换位的迹为不动字母数 \(1\) 减去常数直线上的迹 \(1\)，故为零。
\end{example}

一般情形中，线性无关性与生成性也各需一个依据。在特征 \(p\) 下，\(|G|\) 在 \(k\) 中可能为零，不能用除以群阶的平均投影来照搬普通正交证明；但半单商中的不同矩阵块仍可区分简单模的迹。生成性则可从特征零普通特征标的基出发，再限制到 \(p\)-正则共轭类。
\begin{theorem}\label{b5:thm:brauer-character-basis}
设 \(G\) 为有限群，\(p\) 为素数，\((K,\mathcal O,k)\) 为分裂 \(p\)-模系统，\(K\) 含有全部 \(|G|\) 次单位根，并固定这些根到复数的对应，\(G_{p'}\) 为 \(p\)-正则元素集合。不可约 Brauer 特征标构成 \(G_{p'}\) 上复类函数空间的一组基。因此简单左 \(kG\)-模的同构类型数等于 \(p\)-正则共轭类数；两个有限维左 \(kG\)-模有相同 Brauer 特征标，当且仅当它们有相同的简单合成因子及重数。
\end{theorem}
\begin{proof}
先证线性无关。简单模 \(S_i\) 的 \(k\) 值迹作为 \(kG\) 上的线性函数线性无关：在
\(kG/J\cong\prod_i\operatorname{End}_k(S_i)\)
中，取第 \(i\) 块的一个对角矩阵单位，其各个简单迹为 \(\delta_{ij}\)，再提升到群代数即可检验独立性。群元素是 \(kG\) 的基，因此这些迹限制到 \(G\) 上仍独立。由\cref{b5:lem:modular-trace-regular-part}，所有值由 \(p\)-正则元素上的值确定，故相应矩阵在 \(k\) 中有一个非零的满行秩子式。
设简单模有 \(r\) 个同构类型，选取使该子式非零的 \(p\)-正则元素 \(g_1,\ldots,g_r\)。由\cref{b5:def:brauer-character}，相应提升的行列式
\[
\Delta=\det\bigl(\widetilde\varphi_{S_i}(g_j)\bigr)_{1\le i,j\le r}
\in\mathcal O\cap E
\]
模 \(\pi\) 后等于原来的 \(k\) 值迹子式，因而非零，所以 \(\Delta\ne0\)。固定的域嵌入 \(\iota:E\to\mathbb C\) 是单射，故
\[
\det\bigl(\varphi_{S_i}(g_j)\bigr)_{1\le i,j\le r}
=\iota(\Delta)\ne0.
\]
于是复值 Brauer 特征标线性无关。这给出它们的数目不超过 \(p\)-正则共轭类数，但还没有说明它们生成全部类函数。

为证生成性，转到特征零，在这里群阶可逆，普通平均投影可以使用。先在分裂的特征零群代数 \(KG\) 上对 \(\operatorname{Hom}_K(V_\alpha,V_\beta)\) 的共轭作用取平均。平均算子的像是交织空间，分裂假设与 Schur 引理使其维数为 \(\delta_{\alpha\beta}\)，故取迹得到
\[
\frac1{|G|}\sum_{g\in G}\chi_\alpha(g^{-1})\chi_\beta(g)
=\delta_{\alpha\beta}.
\]
这些等式在固定的分圆域嵌入下仍成立，所以复值普通特征标线性无关。另一方面，\(Z(KG)\) 的共轭类和基说明其维数为共轭类数；半单分裂的矩阵块分解又说明该维数等于简单模数。因此这些特征标恰构成全部复类函数的一组基。任一 \(G_{p'}\) 上的类函数可以在其余共轭类补零，故普通不可约特征标的限制生成所需空间。\cref{b5:prop:ordinary-brauer-restriction}又把每一个这样的限制写成不可约 Brauer 特征标的线性组合，得到生成性。最后用\cref{b5:prop:brauer-character-properties}和刚证的独立性比较系数，得到合成因子的判别。
\end{proof}
对于特征二的 \(S_3\)，\cref{b5:exm:s3-two-blocks}已经给出两个简单模。两类 \(2\)-正则元素由 \(1,r\) 代表，故
\[
\begin{array}{c|rr}
&1&r\\ \hline
\varphi_{S_1}&1&1\\
\varphi_{S_2}&2&-1
\end{array}
\qquad
D=
\begin{pmatrix}
1&0\\1&0\\0&1
\end{pmatrix}.
\]
三行普通表示依次为平凡、符号和标准二维表示。前两者约化后相同；二维整系数标准表示约化后是前述简单模。由此直接计算
\(D^{\mathsf T}D=\operatorname{diag}(2,1)=C_{kS_3}\)。
这是本例的计算；一般的 Cartan--分解矩阵关系及投射模提升需要另作证明，可参阅\cite[Section 9.5, Corollary 9.5.6]{Webb2016FiniteGroupRepresentations}。

% !TeX root = ../../Book5.tex
% 纲要标题：### 8.4 模表示理论补充
% 纲要位置：工作记录/计划纲要.md，第 4001 行。

\section{模表示理论补充}
\label{b5:sec:0804}


% 纲要范围：
%
% * 本章不同于第5章 Brauer 诱导理论
% * 缺陷群和块论的学习地图
% * 不把特征零分类原样搬到正特征
% * 选定诱导格的实际约化同构与 Grothendieck 群中的相容公式；相对迹与 S3 块的 Brauer 像
% * 根层相同但不同构的二维模族；循环 p 群上稳定同态、第一合冲模与短正合扩张的直接计算
%
% ---
%

本章的几个计算已经说明，特征零表示、模表示和投射模之间有联系，却不能彼此替换。本节用约化与诱导的相容性说明第一种联系，再用一个能够直接计算的截断映射说明块的缺陷群。更深的块分类还需要其他工具，这里先把问题和所需理论说明清楚。

\subsection{模表示与 Brauer 诱导理论的区别}
\label{b5:subsec:0804:01}

\chapref{b5:ch:05}的 Brauer 诱导定理在特征零表示环中，把普通特征标写成由某些子群的一维特征标诱导而来的整数线性组合。本章的 Brauer 特征标则在特征 \(p\) 中先取表示，再把 \(p\)-正则特征值提升到特征零；分解矩阵记录整格约化的合成因子。两种理论处理不同的构造，不能因为都冠以 Brauer 之名就将其结论互换。
\begin{proposition}\label{b5:prop:reduction-induction}
	设 \(G\) 为有限群，\(H\le G\)，\(p\) 为素数，并给定对 \(G,H\) 都分裂的 \(p\)-模系统 \((K,\mathcal O,k)\)。记 \(d_G:R_K(G)\to G_0(kG)\)、\(d_H:R_K(H)\to G_0(kH)\) 为分解映射，其中表示群及 Grothendieck 群\footnote{格罗滕迪克（Alexandre Grothendieck，1928-2014），法籍数学家，重塑代数几何的概形与上同调方法。}都取有限维左模。它们与限制、诱导相容：
	\[
		d_G\operatorname{Ind}_H^G =\operatorname{Ind}_H^Gd_H,\qquad d_H\operatorname{Res}_H^G =\operatorname{Res}_H^Gd_G.
	\]
	右侧的诱导、限制作用于由短正合列定义的 \(G_0\) 群。
\end{proposition}
\begin{proof}
	限制一个稳定满格仍是稳定满格，其约化显然与限制交换。若 \(L\) 是 \(KH\)-模 \(V\) 中的满格，则 \(\mathcal OG\otimes_{\mathcal OH}L\) 是 \(\operatorname{Ind}_H^G V\) 中的稳定满格：按左陪集代表，\(\mathcal OG\) 是有限自由右 \(\mathcal OH\) 模，这个张量是若干份 \(L\) 的直和，且扩到 \(K\) 后得到所需诱导模。模 \(\pi\) 后自然同构于 \(kG\otimes_{kH}(L/\pi L)\)。 诱导因右侧自由性保持正合列，限制也保持正合列，所以二者确实诱导 \(G_0\) 上的映射。结合格选择无关性得到公式。
\end{proof}
证明对选定的诱导格给出了实际的模同构
\[
	\bigl(\mathcal OG\otimes_{\mathcal OH}L\bigr)/\pi\bigl(\mathcal OG\otimes_{\mathcal OH}L\bigr) \cong kG\otimes_{kH}(L/\pi L).
\]
而命题中关于 \(d_G,d_H\) 的公式位于 \(G_0\) 中：若先在特征零中分解诱导模，再为各直和项另外选格，其约化只保证具有相同的合成因子，未必与这个诱导格的约化同构。

\begin{example}\label{b5:exm:s3-induced-lattice}
	设 \(k\) 为特征 \(3\) 的域，\(G=S_3\)，\(H=\langle(12)\rangle\)，令 \(S_+,S_-\) 分别为平凡模与符号模。则置换模
	\[
		\operatorname{Ind}_H^G k\cong k^3
	\]
	不可分解且投射，合成因子为两个 \(S_+\) 与一个 \(S_-\)。令 \(U=\{(x,y,z)\in k^3:x+y+z=0\}\)，则 \(k^3\) 与 \(S_+\oplus U\) 的合成因子相同，却不同构。
\end{example}
\begin{proof}
	三个陪集上的置换作用就是三个坐标上的作用。限制到 \(C_3=\langle r\rangle\)、\(r=(123)\) 后，三个坐标构成一个正则轨道，故这个限制模是 \(kC_3\cong k[t]/(t^3)\)，其中 \(t=r-1\)。由\cref{b5:thm:cyclic-p-modules}它不可分解；若 \(k^3\) 有非平凡 \(kG\)-模直和分解，限制后也会分解，矛盾。 因为 \(|H|=2\) 在 \(k\) 中可逆，\(k\) 是投射 \(kH\)-模；诱导把自由模的直和因子送到自由模的直和因子，故 \(k^3\) 投射。

	在 \(U\) 中，\(D=k(1,1,1)\) 是平凡子模，\(U/D\) 是符号模：例如换位 \((12)\) 将 \((1,-1,0)\) 的类变为其负，而这个类非零。求和映射又给出 \(k^3/U\cong S_+\)，得到所列合成因子。\(r\) 在 \(U\) 上不是恒等，所以 \(0\to S_+\to U\to S_-\to0\) 不分裂；否则 \(r\) 会在两个直和项上都作用为恒等。\(S_+\oplus U\) 则有显示出的非平凡直和分解，因而不与 \(k^3\) 同构。
\end{proof}
在一个对 \(S_3,H\) 分裂的 \(3\)-模系统中，诱导格 \(\mathcal O[G/H]\) 的约化就是上述 \(k^3\)。特征零中由普通特征标或坐标求和可得
\[
	\operatorname{Ind}_H^G K\cong K\oplus U_K,\qquad U_K=\{(x,y,z)\in K^3:x+y+z=0\}.
\]
若为右侧另选直和格 \(\mathcal O\oplus\{(x,y,z)\in\mathcal O^3:x+y+z=0\}\)，其约化便是 \(S_+\oplus U\)。这两个格都位于同一个特征零表示中，分解映射将它们识别为同一个合成类，却没有将它们的约化模识别为同一个模。

\subsection{缺陷群与块论简介}
\label{b5:subsec:0804:02}

设 \(G\) 为有限群，\(p\) 为素数，\(k\) 为特征 \(p\) 的域，\(P\le G\) 为 \(p\)-子群。先把 \(P\) 的共轭作用当作置换作用来看。固定空间 \((kG)^P\) 以各个共轭轨道的轨道和为基，其中单点轨道恰好是 \(C_G(P)\) 中的元素。若一个轨道含 \(g\)，令 \(Q=C_P(g)\) 为它在 \(P\) 中的稳定子群，则其轨道和为
\[
	\sum_{x\in P/Q}xgx^{-1}.
\]
这里 \(g\in(kG)^Q\)。对任意子群 \(Q\le P\)，相对迹
\[
	\operatorname{Tr}_Q^P:(kG)^Q\longrightarrow(kG)^P,\qquad z\longmapsto\sum_{x\in P/Q}xzx^{-1}
\]
与代表元选择无关，因为 \(z\) 被 \(Q\) 固定；所得和被 \(P\) 固定。非单点轨道的稳定子群 \(Q\) 是真子群，所以该轨道和正是从较小子群来的相对迹。

\ProseParagraphBreak

反过来，若 \(Q<P\)，则 \(\operatorname{Tr}_Q^P(z)\) 在每个 \(h\in C_G(P)\) 上的系数为 \([P:Q]\) 乘以 \(z\) 在 \(h\) 上的系数；这个指数是正的 \(p\) 次幂，在 \(k\) 中为零。因此保留单点轨道、删去其余轨道的操作，其核恰为全部真子群相对迹的像之和。它保留的是固定空间中不能由这些较小子群的相对迹表示的部分：
\[
	(kG)^P\big/\sum_{Q<P} \operatorname{Tr}_Q^P\bigl((kG)^Q\bigr) \cong kC_G(P)
\]
作为 \(k\) 向量空间同构。下面还要验证，这个系数截断与乘法相容。

\begin{proposition}\label{b5:prop:brauer-truncation}
	设 \(G\) 为有限群，\(p\) 为素数，\(k\) 为特征 \(p\) 的域，\(P\le G\) 为 \(p\)-子群，\(C_G(P)\) 为其中心化子。令 \(P\) 以共轭作用于 \(kG\)，并记固定子代数为 \((kG)^P\)。映射
	\[
		\operatorname{Br}_P^G:(kG)^P\longrightarrow kC_G(P),\qquad \sum_{g\in G}a_gg\longmapsto\sum_{g\in C_G(P)}a_gg
	\]
	是保持单位的 \(k\) 代数同态，称为此处的\term[Brauer-yingshe]{Brauer 映射}[Brauer map]。
\end{proposition}
\begin{proof}
	加法、标量与单位的保持直接来自定义。取两个 \(P\) 共轭不变的元素 \(x=\sum a_gg\)、\(y=\sum b_hh\)，固定 \(z\in C_G(P)\)。乘积中 \(z\) 的系数为
	\[
		\sum_{gh=z}a_gb_h.
	\]
	\(P\) 同时共轭二元组 \((g,h)\)，保持条件 \(gh=z\)，且每个轨道内的系数乘积相同。非平凡轨道的大小是正的 \(p\) 次幂，在 \(k\) 中贡献为零。剩余的固定二元组恰好满足 \(g,h\in C_G(P)\)，其总贡献就是 \(\operatorname{Br}_P^G(x)\operatorname{Br}_P^G(y)\) 的 \(z\) 系数。逐个 \(z\) 比较即得乘法性。
\end{proof}
\begin{definition}\label{b5:def:block-defect-group}
	设 \(G\) 为有限群，\(p\) 为素数，\(k\) 为特征 \(p\) 的域，\(b\) 为 \(kG\) 的块幂等元，对 \(p\)-子群 \(D\le G\)，记 \(\operatorname{Br}_D^G:(kG)^D\to kC_G(D)\) 为保留中心化子上系数的 Brauer 映射。若 \(D\) 满足 \(\operatorname{Br}_D^G(b)\ne0\)，并且在满足这一条件的 \(p\)-子群中按包含关系极大，就称 \(D\) 为块 \(kGb\) 的一个\term[quexian-qun]{缺陷群}[defect group]。
\end{definition}
\symidx{\(\operatorname{Br}_D^G\)：群代数的 Brauer 映射}
由于 \(\operatorname{Br}_{1}^G(b)=b\ne0\)，且 \(G\) 有限，这样的极大子群存在。其进一步的共轭与结构性质属于块论；这里采用支集截断的描述以便计算。这个描述及其与一般定义的关系，可参阅\cite[Corollaries 12.5.4--12.5.5]{Webb2016FiniteGroupRepresentations}。
\ProseParagraphBreak
对\cref{b5:exm:s3-two-blocks}中的 \(S_3\) 与特征二，取 \(P=\langle s\rangle\)。中心化子 \(C_G(P)=P\)，故
\[
	\operatorname{Tr}_{1}^{P}(r)=r+srs=r+r^2.
\]
这是非单点共轭轨道的轨道和，Brauer 映射将它删去；单位元则是单点轨道，被保留下来。因此块幂等元 \(e=1+r+r^2\)、\(f=r+r^2\) 分别满足
\[
	\operatorname{Br}_P^G(1+r+r^2)=1,\qquad \operatorname{Br}_P^G(r+r^2)=0.
\]
三个非平凡 \(2\) 子群互相共轭，因此二维局部块的缺陷群为一个 Sylow \(2\) 子群，四维矩阵块的缺陷群为平凡群。后者称为缺陷零块。前者的简单模 \(S_1\) 不投射，其覆盖 \(kGe\) 有两个 \(S_1\) 合成因子，Cartan 条目为 \(2\)；后者的二维简单模 \(S_2\) 本身就投射，Cartan 条目为 \(1\)。这是本例中由块结构直接算出的对应，不代替一般缺陷零块理论的证明。在任意有限 \(p\) 群上，唯一块的幂等元为一，而对整个群的截断仍为一，所以其缺陷群就是整个群。

\ProseParagraphBreak

定义域中的共轭不变性不能省去。若把同样的支集截断定义在整个 \(kG\) 上，\(r,r^{-1}\) 都被送到零，但 \(rr^{-1}=1\) 被送到一，便不保持乘法。这两个群元素各自都不被 \(P\) 的共轭作用固定；必须先把它们组成轨道和，才能应用\cref{b5:prop:brauer-truncation}。

\subsection{特征零与正特征下的表示分类}
\label{b5:subsec:0804:03}

不可约 Brauer 特征标分类的是简单模，并由加法性确定任意模的合成因子。对于扩张信息，它有明确局限：循环 \(p\) 群上，\(M_d\) 与 \(k^{\oplus d}\) 的 Brauer 特征标都只在单位元处取整数 \(d\)，但前者不可分解，后者在 \(d>1\) 时可分解。分解矩阵同样保留约化后的合成因子，不指定某个整格约化的全部模结构。
\ProseParagraphBreak
有限群在正特征中甚至可以有无限多个不可分解表示。这里给一个不用分类定理的例子。设 \(k\) 是特征二的无限域，\(G=C_2\times C_2=\langle x,y\rangle\)，并在 \(k^2\) 上取
\[
	N=\begin{pmatrix}0&1\\0&0\end{pmatrix},\qquad
	\rho_\lambda(x)=I+N,\qquad \rho_\lambda(y)=I+\lambda N \quad(\lambda\in k).
\]
因为 \(N^2=0\)，两个矩阵交换且平方都为一，故给出表示。其自同态必须与 \(N\) 交换，恰为 \(aI+bN\)，组成局部环，所以每个表示不可分解。若两个参数的表示同构，同构矩阵 \(T\) 与 \(N\) 交换，再由 \(T(\lambda N)=\mu NT\) 得到 \((\lambda-\mu)NT=0\)。\(T\) 可逆且 \(N\ne0\)，故 \(\lambda=\mu\)。于是得到一个参数化的无限族。 记这些模为 \(V_\lambda\)，取基 \(e_1,e_2\) 使 \(Ne_2=e_1\)。由\cref{b5:prop:p-group-local}，\(A=kG\) 的根为 \(J=(x-1,y-1)\)，故对每个参数都有
\[
	JV_\lambda=NV_\lambda=ke_1,\qquad J^2V_\lambda=0, \qquad V_\lambda/JV_\lambda\cong k, \qquad \operatorname{soc}(V_\lambda)=ke_1\cong k.
\]
所以连全部根层也相同。根层各自只记录一个平凡模，而 \(y-1\) 在原模上是 \(x-1\) 的 \(\lambda\) 倍；这一比例记录两层之间的作用关系，却不出现在层商中。
\ProseParagraphBreak
稳定范畴把经过投射模分解的同态视为零。在循环 \(p\) 群上，这种商与扩张的联系也能直接算出。对有限维 \(k\) 代数 \(A\) 及其有限维左模 \(M,N\)，这里以 \(\operatorname{Ext}_A^1(M,N)\) 表示短正合列 \(0\to N\to E\to M\to0\) 的等价类：等价要求中项同构在两端诱导恒等，零类为分裂列。两个扩张先沿对角映射 \(M\to M\oplus M\) 拉回，再沿相加映射 \(N\oplus N\to N\) 推出，得到 Baer 和；沿 \(c1_N\) 推出给出标量 \(c\in k\) 的作用。下面的生成元描述将这些运算化为向量空间中的运算，不要求先掌握消解理论。

\begin{proposition}\label{b5:prop:cyclic-stable-hom-extensions}
	设 \(p\) 为素数，\(k\) 为特征 \(p\) 的域，\(r\ge1\) 为整数、\(m=p^r\)，令 \(A=k[t]/(t^m)\)、\(M_d=A/(t^d)\)，其中 \(1\le d\le m\)。对有限维左 \(A\)-模 \(M,N\)，记 \(\underline{\operatorname{Hom}}_A(M,N)\) 为 \(\operatorname{Hom}_A(M,N)\) 模掉所有可经过有限维投射 \(A\)-模分解的同态所得的商空间。对 \(1\le a,b\le m\)，有
	\[
		\underline{\operatorname{Hom}}_A(M_a,M_b) \cong\ker(t^a:M_b\to M_b)/t^{m-a}M_b.
	\]
	若 \(a<m\)，取投射覆盖 \(A\to M_a\) 的核为第一\term[hechongmo]{合冲模}[syzygy] \(\Omega M_a\)，则
	\[
		\begin{aligned}
			\Omega M_a&=t^aA\cong M_{m-a},\\
			\operatorname{Ext}_A^1(M_a,M_b) &\cong\underline{\operatorname{Hom}}_A(\Omega M_a,M_b)\\
			&\cong\ker(t^{m-a}:M_b\to M_b)/t^aM_b.
		\end{aligned}
	\]
	稳定同态空间与扩张空间的维数均为
	\[
		\min(a,b,m-a,m-b).
	\]
	\(a=m\) 时 \(M_a=A\) 投射，所有上述扩张分裂，维数仍由该式给出。
\end{proposition}
\symidx{\(\underline{\operatorname{Hom}}_A(M,N)\)：模掉经过投射模的同态所得的稳定同态空间}
\symidx{\(\Omega M\)：投射覆盖的核，即第一合冲模}
\symidx{\(\operatorname{Ext}_A^1(M,N)\)：固定两端的短正合模扩张类空间}
\begin{proof}
	同态 \(M_a\to M_b\) 由生成元一的像决定，该像恰好遍历 \(\ker(t^a:M_b\to M_b)\)。由\cref{b5:thm:split-projective-covers}，有限维投射 \(A\)-模都是有限个 \(A\) 的直和。\(M_a\to A\) 中一的像必须属于 \(t^{m-a}A\)，再经 \(A\to M_b\)，所得像恰好遍历 \(t^{m-a}M_b\)。这已是一个子空间，经过若干份 \(A\) 的映射之和不会再扩大它，得到第一式。核的维数为 \(\min(a,b)\)，分母的维数为 \(\max(0,a+b-m)\)，若 \(a+b\le m\)，则 \(m-a\ge b\)、\(m-b\ge a\)，两者之差为 \(\min(a,b)=\min(a,b,m-a,m-b)\)。若 \(a+b>m\)，可由对称性设 \(a\le b\)，差便为 \(a-(a+b-m)=m-b\)；此时 \(m-b<a\le b\)，且 \(m-b\le m-a\)，所以差仍为四者的最小值。

	对 \(a<m\)，映射 \(M_{m-a}\to t^aA\)、\(1\mapsto t^a\) 是同构。给定扩张，取商模生成元的一个提升 \(v\in E\)。因为商上 \(t^a=0\)，有 \(t^av=w\in M_b\)；又因 \(t^mE=0\)，必有 \(t^{m-a}w=0\)。反过来，给定这样的 \(w\)，在
	\[
		M_b\oplus\operatorname{span}_k\{v,tv,\ldots,t^{a-1}v\}
	\]
	上令 \(t\) 在第一项按原作用，在第二组向量上逐次后移，最后把 \(t^{a-1}v\) 送到 \(w\)。条件 \(t^{m-a}w=0\) 保证 \(t^m=0\)，所以这定义了一个具有指定核与商的 \(A\)-模扩张。 改变提升为 \(v+z\)、\(z\in M_b\)，会把参数改为 \(w+t^az\)；固定两端的扩张同构也只能这样改变提升。因此扩张类恰由 \(\ker(t^{m-a})/t^aM_b\) 参数化。Baer 和的拉回中可以取两个提升组成的向量，随后推出把两个核中的参数相加；标量推出则把参数乘以同一标量。故这是向量空间同构，零类恰为分裂扩张。

	将第一式中的 \(a\) 换成 \(m-a\)，便得到这个商空间正是 \(\underline{\operatorname{Hom}}_A(\Omega M_a,M_b)\)。也可从投射覆盖的短正合列理解分母：\(\Omega M_a\) 的生成元是 \(t^a\)，所以限制一个 \(A\to M_b\) 的同态时，该生成元的像恰好遍历 \(t^aM_b\)。扩张空间的维数为
	\[
		\min(m-a,b)-\max(0,b-a)=\min(a,b,m-a,m-b).
	\]
	若 \(a=m\)，投射性使每个到 \(M_a\) 的满射分裂，最后的断言也成立。
\end{proof}
所以这两个空间维数相同并非偶合，但一般并没有在这里把 \(\operatorname{Ext}_A^1(M_a,M_b)\) 与 \(\underline{\operatorname{Hom}}_A(M_a,M_b)\) 识别。与扩张自然相连的是第一变量换成合冲模以后的稳定同态空间。

\ProseParagraphBreak

这些计算解释了为何后续会研究投射覆盖、块、相对投射性和稳定范畴，而不只寻找另一张有限的特征标表。一般缺陷群如何控制块、何时只有有限多种不可分解模，以及不同块之间的等价，都是进一步的问题。本编接下来改从箭图研究有限维代数；那里的表示可以在任意特征上定义，所用的无环性、关系和代数闭性条件仍须逐项写明。

% !TeX root = ../../Book5.tex
% 纲要标题：### 8.5 相对投射性与稳定模范畴
% 纲要位置：工作记录/计划纲要.md（2026-10-20 新增范围）。
% * 状态：初稿已完成；新增范围尚未正式审核。
% 写作范围：群代数对称性、相对迹、顶点与源及稳定合冲模；复用8.4的循环群计算。
% #### 8.5.1 对称群代数与投射内射模
% #### 8.5.2 相对投射性与相对迹
% #### 8.5.3 顶点与源
% #### 8.5.4 稳定同态与合冲模

\section{相对投射性与稳定模范畴}
\label{b5:sec:0805}

模特征下的增广满射 \(kG\to k\) 未必分裂，因而不能再把所有模都当作投射模。不过，一个模可能在限制到某个子群以后已经具有所需的分裂，而从这个子群诱导回来可以保留哪些信息，仍能通过平均式来判断。这里的平均不除以群阶：它把一个子群同态的共轭相加，追问所得和能否恰为恒等映射。另一方面，\cref{b5:prop:cyclic-stable-hom-extensions}已经显示，经过投射模的同态与扩张中的无效选择有关。本节将这两个观察分别发展为相对投射性与稳定模范畴。

\subsection{对称群代数与投射内射模}
\label{b5:subsec:0805:01}

投射性允许沿满射提升映射，内射性允许沿单射延拓映射，二者对一般有限维代数并不相同。群代数有一个不要求群阶可逆的配对：从乘积中取出单位群元素的系数。它将正则模与其线性对偶联系起来，而对偶正则模恰好具有延拓性质。

\begin{definition}\label{b5:def:symmetric-algebra}
	设 \(k\) 为域，\(A\) 为有限维含幺 \(k\) 代数。若存在 \(k\) 线性映射 \(\lambda:A\to k\)，使双线性配对 \((a,b)\mapsto\lambda(ab)\) 非退化，且对所有 \(a,b\in A\) 有 \(\lambda(ab)=\lambda(ba)\)，则称 \(A\) 为\term[duichen-daishu]{对称代数}[symmetric algebra]。这里“对称”描述配对，不要求 \(A\) 的乘法交换；它与由张量代数取交换商所得的对称代数 \(\operatorname{Sym}(V)\) 是不同语境中的用语。
\end{definition}

\begin{proposition}\label{b5:prop:group-algebra-symmetric}
	设 \(G\) 为有限群，\(k\) 为任意域，\(A=kG\)。取 \(\lambda(\sum_g a_gg)=a_1\)，则 \(A\) 是对称代数。对有限维左 \(A\)-模，投射性与内射性等价；每个有限维左模既是有限维投射模的商，又能嵌入有限维投射模。
\end{proposition}
\begin{proof}
	若 \(a=\sum_ga_gg\)、\(b=\sum_gb_gg\)，则
	\[
		\lambda(ab)=\sum_{g\in G}a_gb_{g^{-1}}=\lambda(ba).
	\]
	若 \(a\ne0\)，选 \(g\) 使 \(a_g\ne0\)，便有 \(\lambda(g^{-1}a)=a_g\ne0\)，故配对非退化。它还满足 \(\lambda((ab)c)=\lambda(a(bc))\)，因为代数乘法结合。

	记 \(DA=\operatorname{Hom}_k(A,k)\)，赋予双模作用
	\[
		(a\varphi b)(x)=\varphi(bxa).
	\]
	映射 \(\Theta:A\to DA\)、\(\Theta(c)(x)=\lambda(xc)\) 是双模同构。非退化性及有限维性使它是向量空间同构，而
	\[
		(a\Theta(c)b)(x)=\lambda(bxac)=\lambda(xacb)=\Theta(acb)(x)
	\]
	验证了两侧作用。

	对任意左 \(A\)-模 \(M\)，在一处取值给出自然同构
	\[
		\operatorname{Hom}_A(M,DA)\longrightarrow\operatorname{Hom}_k(M,k),
		\qquad F\longmapsto\bigl(m\mapsto F(m)(1)\bigr).
	\]
	其逆把线性泛函 \(\ell\) 送到 \(F_\ell(m)(a)=\ell(am)\)：确有 \(F_\ell(bm)(a)=\ell(abm)=(bF_\ell(m))(a)\)。域上的线性泛函可从子空间延拓，所以右侧的反变 Hom 函子正合；因此 \(DA\) 内射。由 \(A\cong DA\)，正则模 \(A\) 内射，有限自由模及其直和因子也内射，故有限维投射模都内射。

	反过来，取有限维 \(M\) 的线性对偶的一组基 \(\ell_1,\ldots,\ell_d\)。映射
	\[
		M\longrightarrow(DA)^d,\qquad m\longmapsto(F_{\ell_1}(m),\ldots,F_{\ell_d}(m))
	\]
	是 \(A\) 线性单射：若像为零，在 \(a=1\) 处取值便得所有 \(\ell_i(m)=0\)，从而 \(m=0\)。目标同构于 \(A^d\)，为有限自由模。若 \(M\) 内射，这个单射有左逆，\(M\) 就是有限自由模的直和因子，因而投射。最后，任一有限维模的有限生成集给出有限自由模到它的满射，与上述嵌入一起证明末项。
\end{proof}

这个证明没有使用 \(|G|^{-1}\)。即使 \(p\,\bigm\vert\,|G|\)，投射模仍有内射性，但一般模并不因此投射。例如 \(kC_p\to k\) 仍不能分裂；失效的是平凡模的投射性，不是群代数的对称配对。

\subsection{相对投射性与相对迹}
\label{b5:subsec:0805:02}

限制一个不分裂的 \(G\)-模短正合列到子群 \(H\) 后，它可能变成分裂列。相对投射性只要求把已经存在的 \(H\)-线性分裂改进为 \(G\)-线性分裂，因而比通常的投射性弱。当 \(H=\{1\}\) 时，域上的短正合列总能线性分裂，便恢复通常的投射性。

\begin{definition}\label{b5:def:relative-projectivity}
	设 \(G\) 为有限群，\(k\) 为域，\(H\le G\)，\(M\) 为有限维左 \(kG\)-模。若每个有限维左 \(kG\)-模满射 \(q:E\to M\)，只要限制到 \(H\) 后有 \(kH\)-线性截面，就有 \(kG\)-线性截面，则称 \(M\) \term[xiangdui-toushexing]{相对于 \(H\) 投射}[relatively projective]，也简称 \(H\)-投射。
\end{definition}

给定有限维左 \(kG\)-模 \(M,N\)，群 \(G\) 在 \(\operatorname{Hom}_k(M,N)\) 上按共轭作用：\((g\cdot f)(m)=g f(g^{-1}m)\)。其 \(H\)-不动元恰为 \(H\)-线性同态。取左陪集 \(G/H\) 的代表元集 \(T\)，定义\term[xiangduiji]{相对迹}[relative trace]
\begin{equation}\label{b5:eq:relative-hom-trace}
	\operatorname{Tr}_H^G(f)(m)=\sum_{t\in T}t f(t^{-1}m),
	\qquad f\in\operatorname{Hom}_{kH}(M,N).
\end{equation}
\symidx{\(\operatorname{Tr}_H^G\)：子群同态空间上的相对迹}
若把 \(t\) 换为 \(th\)、\(h\in H\)，相应项不变，因为 \(f\) 是 \(H\) 线性的。任意 \(g\in G\) 置换左陪集，故整个和是 \(G\)-线性的。特别地，它给出 \(\operatorname{End}_{kH}(M)\to\operatorname{End}_{kG}(M)\)，但一般不是环同态。若 \(u,v\) 是可复合的 \(G\)-同态，则
\[
	\operatorname{Tr}_H^G(vfu)=v\operatorname{Tr}_H^G(f)u.
\]
所以自同态相对迹的像是 \(\operatorname{End}_{kG}(M)\) 的双边理想。这与\subsecref{b5:subsec:0804:02}中在共轭不动空间上求迹是同一种构造，只是这里取的共轭表示为 Hom 空间。

\begin{theorem}[Higman 判据]\label{b5:thm:higman-criterion}
	设 \(G\) 为有限群，\(k\) 为域，\(H\le G\)，\(M\) 为有限维左 \(kG\)-模。下列条件等价：
	\begin{equivalences}
		\item \(M\) 相对于 \(H\) 投射。
		\item 对某个有限维左 \(kH\)-模 \(U\)，\(M\) 是 \(\operatorname{Ind}_H^G U\) 的直和因子。
		\item 自然满射
		\[
			\varepsilon_M:kG\otimes_{kH}\operatorname{Res}_H^G M\longrightarrow M,
			\qquad g\otimes m\longmapsto gm
		\]
		有 \(kG\)-线性截面。
		\item 存在 \(\alpha\in\operatorname{End}_{kH}(M)\)，使 \(\operatorname{Tr}_H^G(\alpha)=1_M\)。
	\end{equivalences}
\end{theorem}
\begin{proof}
	若条件 I 成立，则 \(\varepsilon_M\) 的 \(H\)-线性截面 \(m\mapsto1\otimes m\) 可换成 \(G\)-线性截面，得到 III。条件 III 立即给出 II，取 \(U=\operatorname{Res}_H^G M\) 即可。

	为证 II 推出 I，先观察从诱导模出发的提升。设 \(q:E\to V\) 是有 \(H\)-线性截面 \(s\) 的 \(G\)-满射，\(f:\operatorname{Ind}_H^G U\to V\) 是 \(G\)-同态。定义
	\[
		\widetilde f(g\otimes u)=g s(f(1\otimes u)).
	\]
	映射 \(u\mapsto s(f(1\otimes u))\) 是 \(H\)-线性的，所以张量平衡关系使此式良定义，它也是 \(G\)-线性的，且 \(q\widetilde f=f\)。现在若有分裂包含 \(i:M\to\operatorname{Ind}_H^G U\) 和投影 \(r\) 满足 \(ri=1_M\)，对 \(q:E\to M\) 把 \(r\) 提升为 \(\widetilde r\)，则 \(\widetilde r i\) 是 \(q\) 的 \(G\)-线性截面。

	最后把截面直接写成相对迹。对任意 \(H\)-自同态 \(\alpha\)，令
	\[
		\sigma_\alpha(m)=\sum_{t\in T}t\otimes\alpha(t^{-1}m).
	\]
	代表元改变和左乘 \(G\) 的相容性按相对迹的同一计算成立，所以 \(\sigma_\alpha\) 为 \(G\)-同态，且 \(\varepsilon_M\sigma_\alpha=\operatorname{Tr}_H^G(\alpha)\)。因此 IV 推出 III。反过来，令 \(1\in T\)，并以 \(\pi_1\) 表示诱导模投影到 \(1\otimes M\) 后取出 \(M\) 坐标的映射。它是 \(H\)-线性的，因为 \(H\) 左乘保持陪集 \(H\)，也保持其余陪集的并。若 \(\sigma\) 是 \(\varepsilon_M\) 的 \(G\)-线性截面，令 \(\alpha=\pi_1\sigma\)。\(\sigma(m)\) 的 \(t\otimes M\) 坐标等于 \(\pi_1(t^{-1}\sigma(m))=\alpha(t^{-1}m)\)，所以 \(\sigma=\sigma_\alpha\)，继而 \(1_M=\varepsilon_M\sigma=\operatorname{Tr}_H^G(\alpha)\)。
\end{proof}

这个等价条件中的相对迹判别称为\term[Higman-panju]{Higman 判据}[Higman's criterion]\footnote{希格曼（Donald Gordon Higman，1928-2006），加拿大数学家，长期任教于美国密歇根大学，提出相对投射模及其判据。}；一般的相对提升形式见\cite[Proposition 11.3.4]{Webb2016FiniteGroupRepresentations}。它说明要找到的不是任意平均算子，而是一个平均后成为恒等的子群自同态。

\begin{proposition}\label{b5:prop:relative-sylow-projectivity}
	设 \(G\) 为有限群，\(k\) 为特征 \(p>0\) 的域，\(H\le G\)，\(M\) 为有限维左 \(kG\)-模。
	\begin{enumerate}
		\item 若 \(p\nmid[G:H]\)，则 \(M\) 相对于 \(H\) 投射。
		\item 若 \(P\) 为 \(G\) 的 Sylow \(p\) 子群，则 \(M\) 投射，当且仅当 \(\operatorname{Res}_P^G M\) 投射。
		\item 相对投射性对有限直和与直和因子保持。若 \(K\le H\le G\)，一个 \(H\)-模 \(U\) 相对于 \(K\) 投射，则 \(\operatorname{Ind}_H^G U\) 相对于 \(K\) 投射；一个 \(G\)-模相对于 \(K\) 投射，也相对于 \(H\) 投射。
	\end{enumerate}
\end{proposition}
\begin{proof}
	第一项由 \(\operatorname{Tr}_H^G(1_M)=[G:H]1_M\) 得到：以 \([G:H]^{-1}1_M\) 作为 Higman 判据中的 \(\alpha\)。第二项的正向利用 \(kG\) 是有限自由左 \(kP\)-模，故投射模的限制仍是自由模的直和因子。反向利用第一项与 \(M\) 是 \(\operatorname{Ind}_P^G\operatorname{Res}_P^G M\) 的直和因子；若限制模投射，诱导后仍是有限自由 \(kG\)-模的直和因子，因而 \(M\) 投射。

	有限直和与直和因子的结论直接来自诱导直和因子的刻画。若 \(U\) 是 \(\operatorname{Ind}_K^H V\) 的直和因子，则由\cref{b5:thm:induction-transitivity}，\(\operatorname{Ind}_H^G U\) 是 \(\operatorname{Ind}_K^G V\) 的直和因子。最后，一个从 \(K\) 诱导到 \(G\) 的模也可先诱导到 \(H\)，再诱导到 \(G\)，得到最后的断言。
\end{proof}

\subsection{顶点与源}
\label{b5:subsec:0805:03}

每个模都相对于 Sylow \(p\) 子群投射，但某个不可分解模可能只需更小的子群。选择不能再缩小的子群，再保留其中实际诱导出这个模的不可分解分量，就得到顶点与源。困难在于所选的子群与分量未必唯一；Mackey 分解使两个选择能够在交集子群上比较。

\begin{definition}\label{b5:def:vertex-source}
	设 \(G\) 为有限群，\(k\) 为特征 \(p>0\) 的域，\(M\) 为非零有限维不可分解左 \(kG\)-模。在使 \(M\) 相对投射的子群中，按包含关系极小的子群称为 \(M\) 的一个\term[dingdian-mo]{顶点}[vertex]。若 \(P\) 是一个顶点，\(S\) 为非零有限维不可分解左 \(kP\)-模，且 \(M\) 是 \(\operatorname{Ind}_P^G S\) 的直和因子，则称 \(S\) 为与 \(P\) 对应的一个\term[yuan-mo]{源}[source]。
\end{definition}

这里的极小要求不存在真子群仍有相对投射性，而不是先规定一个唯一的子群。子群数目有限，且 \(M\) 总是相对于 \(G\) 投射，所以顶点存在。下述证明中的关键是从一组直和因子中选出一项，这一步在非半单情形依靠局部自同态环。

% 跨卷引用目标：b2:thm:indecomposable-local-end（第二卷定理12.3.6）。
\begin{lemma}\label{b5:lem:relative-intersection}
	设 \(G\) 为有限群，\(k\) 为域，\(M\) 为非零有限维不可分解左 \(kG\)-模。若 \(M\) 同时相对于 \(H,K\le G\) 投射，则对某个 \(g\in G\)，\(M\) 相对于 \(K\cap gHg^{-1}\) 投射。
\end{lemma}
\begin{proof}
	先说明如何从有限直和中选一项。设 \(M\) 是 \(\bigoplus_iV_i\) 的直和因子，包含与投影的分量记为 \(a_i:M\to V_i\)、\(b_i:V_i\to M\)。则 \(1_M=\sum_ib_ia_i\)。有限维模具有有限长度，第二卷定理12.3.6由 Fitting 分解说明，\(\operatorname{End}_{kG}(M)\) 是局部环；它的非单位组成理想，所以这个和中至少有一项 \(b_ia_i\) 可逆。将 \(b_i\) 改为 \((b_ia_i)^{-1}b_i\)，便得到 \(M\) 是单项 \(V_i\) 的直和因子。

	由 Higman 判据，\(M\) 是 \(\operatorname{Ind}_K^G\operatorname{Res}_K^G M\) 的直和因子，也有 \(M\) 到 \(\operatorname{Ind}_H^G\operatorname{Res}_H^G M\) 的分裂包含。对后一对分裂映射先限制到 \(K\)，再诱导到 \(G\)，得到 \(M\) 是下式的直和因子：
	\[
		\operatorname{Ind}_K^G\operatorname{Res}_K^G
		\operatorname{Ind}_H^G\operatorname{Res}_H^G M
		\cong\bigoplus_{g\in K\backslash G/H}\operatorname{Ind}_{K\cap gHg^{-1}}^G W_g,
	\]
	其中 \(W_g\) 为 \({}^g(\operatorname{Res}_H^G M)\) 限制到 \(K\cap gHg^{-1}\) 的模。此同构用的是\cref{b5:thm:mackey-decomposition,b5:thm:induction-transitivity}，它们均不要求半单性。从这个有限直和中选出一项，再用 Higman 判据的诱导刻画，即得结论。
\end{proof}

\begin{theorem}\label{b5:thm:vertices-sources}
	设 \(G\) 为有限群，\(k\) 为特征 \(p>0\) 的域，\(M\) 为非零有限维不可分解左 \(kG\)-模。
	\begin{enumerate}
		\item \(M\) 的顶点都是 \(p\) 子群，并且彼此在 \(G\) 中共轭。若 \(M\) 相对于 \(H\le G\) 投射，则 \(H\) 包含某个顶点。
		\item 给定顶点 \(P\)，存在与它对应的源 \(S\)。每个这样的源都是 \(\operatorname{Res}_P^G M\) 的直和因子，且作为 \(kP\)-模具有顶点 \(P\)。
		\item 若 \(S,S'\) 是同一顶点 \(P\) 的两个源，则存在 \(g\in N_G(P)\) 使 \(S\cong{}^gS'\)。
	\end{enumerate}
\end{theorem}
\begin{proof}
	固定顶点 \(P\)，若 \(M\) 相对于 \(H\) 投射，对\cref{b5:lem:relative-intersection}取 \(K=P\)，得到 \(M\) 相对于 \(P\cap gHg^{-1}\) 投射。顶点的极小性使该交集等于 \(P\)，所以 \(g^{-1}Pg\le H\)。共轭保留相对投射性：由公式 \(a\otimes u\mapsto ag^{-1}\otimes u\)，从 \(g^{-1}Pg\) 诱导适当共轭模所得的 \(G\)-模与原诱导模同构；极小性也随之保留。因此 \(H\) 包含一个顶点。若 \(H\) 本身是顶点，其极小性使 \(g^{-1}Pg=H\)，得到共轭唯一性。再取 \(H\) 为 Sylow \(p\) 子群，\cref{b5:prop:relative-sylow-projectivity}保证 \(M\) 相对于 \(H\) 投射，故 \(P\) 也是 \(p\) 子群。

	% 跨卷引用目标：b2:thm:krull-schmidt（第二卷定理12.4.2）。
	把有限维 \(\operatorname{Res}_P^G M\) 分解为有限个不可分解模；存在性可由不断分解并比较维数得到，唯一性由第二卷定理12.4.2的 Krull--Schmidt 定理保证。\(M\) 是这个限制模诱导回去后的直和因子，上一引理证明中的局部环论证使它来自某个不可分解项 \(S\)，这就构造了一个同时为限制直和因子的源。若这个 \(S\) 相对于 \(R<P\) 投射，诱导传递性会使 \(M\) 相对于 \(R\) 投射，与顶点的极小性矛盾。因此 \(S\) 在 \(P\) 中的顶点是 \(P\)。

	现令 \(S'\) 为任意另一个源。由于 \(S\) 是 \(\operatorname{Res}_P^G M\) 的直和因子，而 \(M\) 是 \(\operatorname{Ind}_P^G S'\) 的直和因子，Mackey 分解给出 \(S\) 是
	\[
		\operatorname{Res}_P^G\operatorname{Ind}_P^G S'
		\cong\bigoplus_{g\in P\backslash G/P}
		\operatorname{Ind}_{P\cap gPg^{-1}}^P
		\operatorname{Res}_{P\cap gPg^{-1}}^{gPg^{-1}}({}^gS')
	\]
	的一项中的直和因子。于是 \(S\) 相对于 \(P\cap gPg^{-1}\) 投射。刚证明的顶点性质迫使 \(P\cap gPg^{-1}=P\)；两群阶相同，故 \(gPg^{-1}=P\)，即 \(g\in N_G(P)\)。对应的 Mackey 项就是 \({}^gS'\)，它不可分解，所以非零直和因子 \(S\) 必与它同构。这证明了第三项。

	最后 \(S'\cong{}^{g^{-1}}S\) 也是 \(\operatorname{Res}_P^G M\) 的直和因子：将 \(S\) 的分裂包含和投影共轭后，其目标为 \({}^{g^{-1}}(\operatorname{Res}_P^G M)\)，而映射 \(m\mapsto g^{-1}m\) 将后者同构到 \(\operatorname{Res}_P^G M\)。同样的诱导传递性又使 \(S'\) 不可能相对于 \(P\) 的真子群投射。于是第二项对每个源都成立。
\end{proof}

以上结论的有限维域版本不要求 \(k\) 代数闭；可与\cite[Theorem 11.6.1 and Proposition 11.6.2]{Webb2016FiniteGroupRepresentations}中的更一般系数环版本比较。顶点记录诱导所需的最小子群，源还记录这个子群上实际使用的模。只知道顶点通常不能恢复原模。

\begin{example}\label{b5:exm:trivial-module-vertex}
	设 \(G\) 为有限群，\(k\) 为特征 \(p>0\) 的域。平凡模 \(k\) 的顶点恰为 Sylow \(p\) 子群，在每个顶点上的源都是一维平凡模。非零不可分解 \(kG\)-模具有平凡顶点，当且仅当它投射。
\end{example}
\begin{proof}
	\(\operatorname{End}_{kH}(k)=k\)，而 \(\operatorname{Tr}_H^G\) 在此空间上就是乘以 \([G:H]\)。Higman 判据说明 \(k\) 相对于 \(H\) 投射，当且仅当 \(p\nmid[G:H]\)，也就是 \(H\) 包含一个 Sylow \(p\) 子群。极小的这样的子群正是 Sylow 子群。源是限制模 \(k\) 的非零直和因子，只能是它本身。另一方面，从平凡群诱导的有限维模是有限自由 \(kG\)-模，所以相对于平凡群投射恰是通常的投射性。
\end{proof}

例如 \(G=C_p\) 时，\(M_1,\ldots,M_{p-1}\) 都有顶点 \(G\)，相应源就是各模本身；唯一的不可分解投射模 \(M_p\) 的顶点为平凡群。这个最小例子已经表明，同一个顶点可以对应许多不同的源。

\subsection{稳定同态与合冲模}
\label{b5:subsec:0805:04}

\cref{b5:prop:cyclic-stable-hom-extensions}把循环群上的扩张写成 \(\underline{\operatorname{Hom}}(\Omega M,N)\)。一般群上的同一公式，需要先解释为什么“经过投射模”能够对同态取商，以及选择投射满射、选择提升时的歧义为什么在商中消失。

\begin{definition}\label{b5:def:stable-group-module-category}
	设 \(G\) 为有限群，\(k\) 为域。对有限维左 \(kG\)-模 \(M,N\)，令 \(\mathcal P(M,N)\) 为所有能分解为 \(M\to Q\to N\)、其中 \(Q\) 是有限维投射模的同态所成的子空间。以有限维左模为对象、以
	\[
		\underline{\operatorname{Hom}}_{kG}(M,N)
		=\operatorname{Hom}_{kG}(M,N)/\mathcal P(M,N)
	\]
	为态射空间，并以代表映射的复合为复合，得到\term[wending-mo-fanchou]{稳定模范畴}[stable module category]，记为 \(\underline{\operatorname{mod}}(kG)\)。
\end{definition}
\symidx{\(\underline{\operatorname{mod}}(kG)\)：有限维群代数模的稳定范畴}
有限个分解经过投射模的同态之和，可以经过这些投射模的直和来分解，故 \(\mathcal P(M,N)\) 确为子空间；复合任意模同态仍有这种分解，故取商后复合良定义。一个模在稳定范畴中同构于零，当且仅当它投射：若 \(1_M\) 经过投射模分解，\(M\) 就是该投射模的直和因子。特别地，稳定范畴保留的是忽略投射直和项以后的对象与映射信息，而不是把原模中的投射子模逐个删除。

\begin{proposition}\label{b5:prop:stable-syzygy-functor}
	设 \(G\) 为有限群，\(k\) 为域。对每个有限维左 \(kG\)-模 \(M\) 选择有限维投射满射 \(\pi_M:P_M\to M\)，记其核为 \(\Omega M\)。这给出稳定模范畴上的函子
	\[
		\Omega:\underline{\operatorname{mod}}(kG)
		\longrightarrow\underline{\operatorname{mod}}(kG),
	\]
	不同选择得到的函子自然同构。它是一个等价函子，逆函子可由嵌入有限维内射模后的余核构造。
\end{proposition}
\begin{proof}
	给定 \(f:M\to N\)，\(P_M\) 的投射性使 \(f\pi_M\) 提升为 \(\widetilde f:P_M\to P_N\)。限制到核得到下图的左侧映射：
	\[
		\begin{tikzcd}[column sep=small]
			0\arrow[r]&\Omega M\arrow[r,hook]\arrow[d,"\Omega f"]&P_M\arrow[r,"\pi_M"]\arrow[d,"\widetilde f"]&M\arrow[r]\arrow[d,"f"]&0\\
			0\arrow[r]&\Omega N\arrow[r,hook]&P_N\arrow[r,"\pi_N"]&N\arrow[r]&0.
		\end{tikzcd}
	\]
	两个提升之差的像包含在 \(\Omega N\) 中，所以它们在 \(\Omega M\) 上的限制之差经过 \(P_M\) 分解，在稳定范畴中相等。若 \(f\) 分解为 \(M\xrightarrow{u}Q\xrightarrow{v}N\)，其中 \(Q\) 投射，可将 \(v\) 提升为 \(\widetilde v:Q\to P_N\)，并选 \(\widetilde f=\widetilde v u\pi_M\)。这个提升在 \(\Omega M\) 上为零。因此 \(\Omega f\) 只依赖 \(f\) 的稳定类。提升的复合可以作为复合映射的提升，恒等映射也有恒等提升，故这确实定义函子。

	若另选 \(\pi'_M:P'_M\to M\)，把 \(1_M\) 沿两个满射相互提升，得到两个核之间的稳定映射。其复合与恒等映射都是 \(1_M\) 的提升在核上的限制，故稳定类相同。这给出稳定同构；对任意 \(f\)，两个比较路径也都是同一个映射的提升在核上的限制，因此同构自然。对象层面的差别还可显式写出。拉回
	\[
		X=\{\,(u,v)\in P_M\oplus P'_M\mid \pi_M(u)=\pi'_M(v)\,\}
	\]
	分别投影到 \(P'_M,P_M\)，核分别为 \(\Omega M,\Omega' M\)。两个目标均投射，所以两条满射均分裂，从而
	\begin{equation}\label{b5:eq:syzygy-schanuel}
		\Omega M\oplus P'_M\cong X\cong\Omega'M\oplus P_M.
	\end{equation}
	因此不同的核确实只相差投射直和项。

	由\cref{b5:prop:group-algebra-symmetric}，可选单射 \(M\to I_M\)，其中 \(I_M\) 同时投射与内射，令 \(\Sigma M=I_M/M\)。把前面对满射的提升改为沿单射的延拓，便在稳定范畴上得到函子 \(\Sigma\)：对 \(f:M\to N\)，将 \(M\to N\to I_N\) 延拓到 \(I_M\)，再取余核；两个延拓之差在 \(M\) 上为零，因而诱导的余核映射之差经过 \(I_N\) 分解，所以稳定类与延拓选择无关。若 \(f\) 经过投射模 \(Q\) 分解，因 \(Q\) 内射，可将 \(M\to Q\) 延拓到 \(I_M\)，使所选延拓的像包含在 \(N\subseteq I_N\) 中；其余核映射便为零。这也保证 \(\Sigma\) 对稳定类良定义。将 \(0\to M\to I_M\to\Sigma M\to0\) 与所选的 \(P_{\Sigma M}\to\Sigma M\) 比较，\eqref{b5:eq:syzygy-schanuel}给出 \(\Omega\Sigma M\cong M\) 于稳定范畴。对单射作对偶的推出比较，或对 \(0\to\Omega M\to P_M\to M\to0\) 重复余核比较，得到 \(\Sigma\Omega M\cong M\)。前述提升与延拓的唯一性同时保证两种同构自然，故 \(\Sigma\) 是 \(\Omega\) 的逆等价。
\end{proof}

取投射覆盖时，\(\Omega M\) 就是\subsecref{b5:subsec:0804:03}中使用的第一合冲模；上面的构造也允许任意投射满射，不必先把多余投射分量消去。稳定范畴中的等价 \(\Omega\) 并不是所有实际模的同构类上的置换，因为投射模的覆盖可以取恒等，其合冲模为零。

\begin{proposition}\label{b5:prop:stable-hom-ext}
	设 \(G\) 为有限群，\(k\) 为域，\(M,N\) 为有限维左 \(kG\)-模。取正合列 \(0\to\Omega M\xrightarrow{i}P\xrightarrow{\pi}M\to0\)，其中 \(P\) 有限维且投射。短正合扩张的 Baer 和给出自然的 \(k\) 向量空间同构
	\[
		\operatorname{Ext}_{kG}^1(M,N)
		\cong\underline{\operatorname{Hom}}_{kG}(\Omega M,N).
	\]
\end{proposition}
\begin{proof}
	给定扩张 \(0\to N\xrightarrow{j}E\xrightarrow{q}M\to0\)，将 \(\pi:P\to M\) 提升为 \(s:P\to E\)。由于 \(qsi=0\)，存在唯一 \(f:\Omega M\to N\) 满足 \(jf=si\)。改变 \(s\) 所产生的 \(f\) 之差恰为某个同态 \(P\to N\) 在 \(\Omega M\) 上的限制。

	反过来，给定 \(f:\Omega M\to N\)，构造推出
	\[
		E_f=(N\oplus P)/\{\,(f(x),-i(x))\mid x\in\Omega M\,\}.
	\]
	以 \(n\mapsto[(n,0)]\) 嵌入 \(N\)，以 \([(n,y)]\mapsto\pi(y)\) 映到 \(M\)。前者的核为零，因为 \(i\) 单射；后者满射，且当 \(\pi(y)=0\) 时，写 \(y=i(x)\)，便有 \([(n,y)]=[(n+f(x),0)]\)。所以得到所需短正合列。从扩张及提升 \(s\) 得到的 \(f\)，反构造后由 \([(n,y)]\mapsto j(n)+s(y)\) 恢复原扩张。若 \(f'=f+hi\)，则 \([(n,y)]\mapsto[(n-h(y),y)]\) 给出固定两端的 \(E_f\cong E_{f'}\)。反过来，固定两端的扩张同构把两组提升相差的映射写成 \(jh\)，故对应参数也只相差 \(hi\)。于是扩张类组成的空间识别为
	\[
		\operatorname{Hom}_{kG}(\Omega M,N)
		\big/\operatorname{im}\bigl(\operatorname{Hom}_{kG}(P,N)\to\operatorname{Hom}_{kG}(\Omega M,N)\bigr).
	\]
	Baer 和的拉回、推出把两个参数相加，沿标量的推出把参数乘该标量，因而此识别为线性同构；沿两端模同态拉回与推出也分别对应参数的复合，所以它自然。

	最后比较分母。每个限制 \(hi\) 都经过投射模 \(P\)。反过来，若 \(f\) 经过投射模 \(Q\) 分解为 \(\Omega M\xrightarrow{u}Q\xrightarrow{v}N\)，则 \(Q\) 由\cref{b5:prop:group-algebra-symmetric}也是内射模。故 \(u\) 沿 \(i\) 延拓为 \(\widetilde u:P\to Q\)，于是 \(f=(v\widetilde u)i\)。分母正是全部经过投射模的同态，得到稳定 Hom 公式。
\end{proof}

投射模的内射性在最后一步有明确作用：它保证经过任意投射模的映射，都能延拓到所选的 \(P\)。对一般有限维代数，仅有 \(P\) 投射不能推出这一点，因此不能无条件把上述扩张商写成稳定 Hom。群代数的对称性使这个公式成立，也解释了\cref{b5:prop:cyclic-stable-hom-extensions}中循环群算例背后的共同机制。第四卷附录 D 进一步在有足够多投射与内射且二者相同的正合范畴中构造三角结构；本节已经给出了这里所需的模、稳定映射及合冲模等价。

% !TeX root = ../../Book5.tex
% 纲要标题：### 8.6 Brauer 对应与 Green 对应
% 纲要位置：工作记录/计划纲要.md（2026-10-20 新增骨架）。
% * 状态：初稿已完成；新增范围尚未正式审核。
% 写作范围：依次标明 Brauer 主定理及 Green 对应的假设、陈述、证明概要与引用范围，不承担深层块分类及研究猜想。
% #### 8.6.1 缺陷群的共轭性与 Brauer 对应
% #### 8.6.2 Brauer 三个主定理的陈述与证明范围
% #### 8.6.3 Green 对应
% #### 8.6.4 小群的块与对应计算

\section{Brauer 对应与 Green 对应}
\label{b5:sec:0806}

对一个块取 Brauer 像，得到的是较小群代数中的幂等元；把一个不可分解模限制到子群，得到的模却可能分裂。这两个操作都有选择问题：哪个局部块应当与原块对应，限制模中的哪一项应当与原模对应？缺陷群为第一个问题指定子群，顶点为第二个问题指定直和项，后者由 Green 对应实现\footnote{格林（James Alexander Green，1926-2014），英国数学家，研究有限群表示，建立不可分解模的 Green 对应及有限一般线性群的特征标理论。}。本节分别说明这两个选择怎样成立，并用 \(S_3\) 比较两种对应。涉及块的主定理时，假设特征 \(p\) 的域 \(k\) 对所讨论群及其全部子群都是分裂域；取 \(k\) 代数闭可满足这一要求。单纯关于顶点和 Green 对应的结论不需要分裂假设。

\subsection{缺陷群的共轭性与 Brauer 对应}
\label{b5:subsec:0806:01}

\cref{b5:def:block-defect-group}用 \(\operatorname{Br}_D^G(b)\ne0\) 的极大性定义缺陷群。它便于从块幂等元的系数计算，却没有立即说明两个极大子群为什么共轭。把块同时当作左模和右模，就可以利用顶点的唯一性。具体地，令 \(G\times G\) 在 \(kG\) 上作用为
\[
	(g,h)\cdot a=gah^{-1},\qquad
	\Delta H=\{\,(h,h)\mid h\in H\,\}\le G\times G.
\]
\(G\times G\) 的子模恰是双侧理想；块 \(kGb\) 因而是不可分解的 \(k[G\times G]\)-模。这里一定要保留两侧作用：块作为左 \(kG\)-模本来就是投射模，其左模顶点为一，不能记录缺陷群。

\begin{theorem}\label{b5:thm:defect-vertices-traces}
	设 \(G\) 为有限群，\(k\) 为特征 \(p>0\) 的域，\(b\) 为 \(kG\) 的块幂等元。存在 \(p\)-子群 \(D\le G\)，使 \(\Delta D\) 为上述双模 \(kGb\) 的顶点。对任意这样的 \(D\)，有
	\[
		b=\operatorname{Tr}_D^G(a)\quad\text{对某个}\;a\in(kGb)^D,
	\]
	且对每个 \(p\)-子群 \(Q\le G\)，
	\[
		\operatorname{Br}_Q^G(b)\ne0
		\quad\Longleftrightarrow\quad
		Q\le {}^gD\quad\text{对某个}\;g\in G.
	\]
	特别地，\(b\) 的缺陷群恰为 \(D\) 的 \(G\)-共轭子群。
\end{theorem}
\begin{proof}
	映射 \((g,h)\otimes1\mapsto gh^{-1}\) 给出
	\[
		\operatorname{Ind}_{\Delta G}^{G\times G}k\cong kG.
	\]
	因此 \(kGb\) 相对 \(\Delta G\) 投射。由\cref{b5:thm:vertices-sources}，可以在 \(\Delta G\) 中选择它的一个顶点，必为 \(\Delta D\)。若 \(\Delta D_1=(g,h)\Delta D(g,h)^{-1}\)，投影到第一坐标就有 \(D_1=gDg^{-1}\)，所以这些 \(D\) 互相共轭。

	对任意 \(H\le G\)，先证明相对迹判据
	\[
		kGb\text{ 相对 }\Delta H\text{ 投射}
		\quad\Longleftrightarrow\quad
		b=\operatorname{Tr}_H^G(a)\text{，其中 }a\in(kGb)^H.
	\]
	正向先限制到 \(\Delta G\)。Mackey 分解中出现的交集为 \(\Delta G\cap(g,h)\Delta H(g,h)^{-1}\)；将 \((g,h)\) 写成 \((g,g)(1,g^{-1}h)\)，可见每个交集都包含于某个 \(\Delta G\)-共轭的 \(\Delta H\)。所以该限制模相对 \(\Delta H\) 投射。Higman 判据给出 \(1_{kGb}=\operatorname{Tr}_H^G(\alpha)\)，其中 \(\alpha\) 对 \(H\) 的共轭作用线性。将等式作用到被 \(G\) 固定的 \(b\)，便得 \(b=\operatorname{Tr}_H^G(\alpha(b))\)。反向若 \(b=\operatorname{Tr}_H^G(a)\)，则左乘 \(a\) 是共轭 \(H\) 模自同态，其相对迹为左乘 \(b\)，也就是恒等。因此限制到 \(\Delta G\) 的模相对 \(\Delta H\) 投射。原模已相对 \(\Delta G\) 投射，再利用诱导的传递性，即得原模相对 \(\Delta H\) 投射。

	取顶点 \(\Delta D\)，便有所需的迹表达式，且 \(D\) 对此性质极小。若 \(\operatorname{Br}_D^G(b)=0\)，由\subsecref{b5:subsec:0804:02}中按共轭轨道求得的核，可写
	\[
		b=\sum_{Q<D}\operatorname{Tr}_Q^D(u_Q),\qquad u_Q\in(kG)^Q.
	\]
	与 \(b=\operatorname{Tr}_D^G(a)\) 相乘，利用 \(a\) 的 \(D\)-不变性和迹的传递性，再乘一次中心元 \(b\)，得到
	\[
		b=\sum_{Q<D}\operatorname{Tr}_Q^G(ba u_Q).
	\]
	各项都在交换局部代数 \(bZ(kG)\) 中。这里 \(bZ(kG)\) 就是不可分解双模 \(kGb\) 的自同态环：双模同态由 \(b\) 的像决定，该像须与整个群代数交换；局部性因此由\cref{b5:lem:relative-intersection}证明中回顾的局部自同态环结论给出。一个局部环中若若干元素之和为单位，至少一项是单位。于是某个 \(\operatorname{Tr}_Q^G(ba u_Q)\) 在 \(bZ(kG)\) 中可逆，乘其中心逆元就把 \(b\) 写成从真子群 \(Q<D\) 来的相对迹，矛盾。因此 \(\operatorname{Br}_D^G(b)\ne0\)。

	对任意 \(p\)-子群 \(Q\)，将 \(\operatorname{Tr}_D^G(a)\) 按 \(Q\backslash G/D\) 分组，得到
	\[
		b=\sum_{x\in Q\backslash G/D}
		\operatorname{Tr}_{Q\cap{}^xD}^{Q}(xax^{-1}).
	\]
	若所有交集都是 \(Q\) 的真子群，Brauer 映射杀掉右侧每一项。因此像非零迫使某个 \(Q\le{}^xD\)。反过来，若 \(Q\le{}^xD\)，则 \(C_G({}^xD)\subseteq C_G(Q)\)；\(b\) 在前者上已有非零系数，在后者上的截断也非零。于是极大的 \(Q\) 正是 \(D\) 的各个共轭。
\end{proof}

这个证明还使缺陷群与块中的各个模发生联系。
\begin{corollary}\label{b5:cor:block-modules-relative-defect}
	设 \(G\) 为有限群，\(k\) 为特征 \(p>0\) 的域，块 \(kGb\) 的缺陷群为 \(D\)。若有限维左 \(kG\)-模 \(M\) 满足 \(bM=M\)，则 \(M\) 相对 \(D\) 投射。特别地，不可分解的 \(M\) 有一个顶点包含于 \(D\)。
\end{corollary}
\begin{proof}
	取 \(b=\operatorname{Tr}_D^G(a)\)。元素 \(a\) 在 \(M\) 上的作用是 \(kD\)-自同态，而它的相对迹是 \(b\) 的作用，即 \(1_M\)。由\cref{b5:thm:higman-criterion}得到相对投射性，再用顶点的包含性质即得后一句。
\end{proof}

现在固定 \(p\)-子群 \(D\)，令 \(N=N_G(D)\)。对 \(z\in Z(kG)\)，\(\operatorname{Br}_D^G(z)\) 不仅在 \(kC_G(D)\) 中，而且被 \(N\) 共轭固定，因此也属于 \(Z(kN)\)。若 \(1=\sum_i b_i\) 是 \(kG\) 的块分解，那么 \(1=\sum_i\operatorname{Br}_D^G(b_i)\) 是 \(Z(kN)\) 中的正交幂等分解，部分项可以为零。
\begin{definition}\label{b5:def:brauer-correspondent}
	设 \(G\) 为有限群，\(k\) 为特征 \(p>0\) 的域，\(D\le G\) 为 \(p\)-子群，\(N=N_G(D)\)，\(c\) 为 \(kN\) 的块幂等元。在 \(kG\) 的块幂等元中，恰有一个 \(b\) 满足
	\[
		c\operatorname{Br}_D^G(b)=c.
	\]
	记此 \(b\) 为 \(c^G\)，称为 \(c\) 的 \term[Brauer-duiying-kuai]{Brauer 对应块}[Brauer correspondent]。这里固定使用子群 \(N_G(D)\) 及映射 \(\operatorname{Br}_D^G\)。
\end{definition}
\symidx{\(c^G\)：子群块的 Brauer 对应块}
存在与唯一性只用 \(c\) 的本原性：\(c=\sum_i c\operatorname{Br}_D^G(b_i)\) 是正交中心幂等分解，恰有一项非零。此时尚不能断言 \(c\mapsto c^G\) 为双射，也不能断言两块具有相同缺陷群；第一主定理需要明确加入缺陷群恰为 \(D\) 的条件。

\subsection{Brauer 三个主定理的陈述与证明范围}
\label{b5:subsec:0806:02}

\term[Brauer-diyi-zhudingli]{Brauer 第一主定理}[Brauer's first main theorem]把缺陷群恰为 \(D\) 的块与其正规化子中的相应块配对。这里固定实际子群 \(D\)，块的缺陷群列表则仍按共轭理解。

\begin{theorem}[Brauer 第一主定理]\label{b5:thm:brauer-first-main}
	设 \(G\) 为有限群，\(k\) 为特征 \(p>0\) 的域，并且对 \(G\) 及其全部子群都是分裂域。给定 \(p\)-子群 \(D\le G\)，记 \(N=N_G(D)\)。则
	\[
		b\longmapsto\operatorname{Br}_D^G(b)
	\]
	给出缺陷群为 \(D\) 的 \(kG\) 块与缺陷群为 \(D\) 的 \(kN\) 块之间的双射，逆映射为 \(c\mapsto c^G\)。特别地，在这一范围内，\(\operatorname{Br}_D^G(b)\) 本身就是一个块幂等元。
\end{theorem}
\begin{proofsketch}
	相对迹与截断的相容公式是
	\[
		\operatorname{Br}_D^G\operatorname{Tr}_D^G(a)
		=\operatorname{Tr}_D^N\operatorname{Br}_D^G(a)
		\qquad(a\in(kG)^D).
	\]
	按 \(D\backslash G/D\) 分组后，\(D\cap{}^xD<D\) 的项都被截断杀掉；剩下的代表元恰来自 \(N/D\)。若 \(c\) 是 \(kN\) 中缺陷为 \(D\) 的块，令 \(b=c^G\)。利用这一公式和\cref{b5:thm:defect-vertices-traces}，可以把 \(c\) 放入
	\[
		\operatorname{Br}_D^G\bigl(b\operatorname{Tr}_D^G((kG)^D)\bigr)
	\]
	中。括号内是 \(bZ(kG)\) 的理想。它的像含非零幂等元，故该理想不可能是幂零理想；有限维局部代数的真理想全在幂零根中，所以它必须等于 \(bZ(kG)\)。于是 \(b\) 来自 \(D\) 的相对迹，而 \(\operatorname{Br}_D^G(b)\ne0\)，从而 \(D\) 正是它的缺陷群。同一个局部代数的像至多有一个非零幂等元，进一步给出 \(\operatorname{Br}_D^G(b)=c\)。

	这段论证还使用两个块论结论：正规 \(p\)-子群包含在每个块的缺陷群中，且块幂等元的支集包含在该正规子群的中心化子中。前者保证 \(kN\) 的块的缺陷群包含 \(D\)，后者允许在相对迹中先做截断。它们也使反向构造穷尽 \(kG\) 中缺陷为 \(D\) 的块。这里没有证明这两个一般结论；其证明分别见\cite[Corollary 12.3.4 and Proposition 12.5.7]{Webb2016FiniteGroupRepresentations}，上述局部理想论证及双射的完整证明见该书 Theorem 12.6.4。
\end{proofsketch}

\term[Brauer-dier-zhudingli]{Brauer 第二主定理}[Brauer's second main theorem]联系的是普通特征标在 \(p\)-奇异元素处的值与中心化子的 Brauer 特征标。设 \((K,\mathcal O,k)\) 是对 \(G\) 及其子群分裂的 \(p\)-模系统，\(\mathcal O\) 完备，且 \(K\) 含有 \(|G|\) 次单位根。给定 \(p\)-元素 \(u\in G\)，令 \(C=C_G(u)\)。对普通不可约特征标 \(\chi\)，函数 \(v\mapsto\chi(uv)\) 是 \(C\) 的 \(p\)-正则元素上的类函数，因而可唯一写成
\begin{equation}\label{b5:eq:generalized-decomposition-numbers}
	\chi(uv)=\sum_{\varphi\in\operatorname{IBr}(C)}d^u_{\chi\varphi}\varphi(v),
	\qquad v\in C_{p'}.
\end{equation}
其中使用\cref{b5:thm:brauer-character-basis}。普通特征标、Brauer 特征标及单位根都沿用\secref{b5:sec:0803}中固定的分圆域嵌入，统一取复数值。系数称为\term[guangyi-fenjie-shu]{广义分解数}[generalized decomposition number]。它们一般不是非负整数：将表示限制到 \(C\) 后，中心元 \(u\) 的各特征空间均是 \(KC\)-子模；每个子模的分解数为整数，再乘 \(u\) 的相应特征值并相加，就得到 \(d^u_{\chi\varphi}\)。因此这些系数是 \(p\) 次幂单位根的整数线性组合；\(u=1\) 时恢复普通分解数。
\symidx{\(d^u_{\chi\varphi}\)：广义分解数}

还须说明普通特征标“属于一个模特征块”的含义。由于 \(\mathcal O\) 完备，\(kG\) 的每个中心幂等元 \(b\) 唯一提升为 \(Z(\mathcal OG)\) 的幂等元 \(\widehat b\)。这里的提升可直接在中心中进行：共轭类和给出中心的基，因此约化映射 \(Z(\mathcal OG)\to Z(kG)\) 满射；从任一近似提升 \(a\) 出发，因 \((2a-1)^2\equiv1\pmod\pi\)，\(2a-1\) 在完备代数中可逆；以 \(a-(a^2-a)(2a-1)^{-1}\) 反复修正，误差的 \(\pi\)-进阶数加倍，完备性给出极限。两个同余的中心幂等提升之差为零，因为它们的交叉项 \(e(1-e')\)、\(e'(1-e)\) 都是落在 \(\pi\mathcal OG\) 中的幂等元，必须落在每个 \(\pi^n\mathcal OG\) 中，因而为零。各块的提升仍正交且和为一，所以一个不可约 \(KG\)-模上恰有一个 \(\widehat b\) 作用为恒等；它的普通特征标便称为属于块 \(b\)。

\begin{theorem}[Brauer 第二主定理]\label{b5:thm:brauer-second-main}
	设 \(G\) 为有限群，\((K,\mathcal O,k)\) 为对 \(G\) 及其全部子群分裂的 \(p\)-模系统，其中 \(\mathcal O\) 完备，且 \(K\) 含有 \(|G|\) 次单位根。取 \(p\)-元素 \(u\in G\)、普通不可约特征标 \(\chi\)，以及 \(C=C_G(u)\) 的不可约 Brauer 特征标 \(\varphi\)。若 \(b\) 为 \(\chi\) 所属块在 \(kG\) 中的幂等元，\(c\) 为 \(\varphi\) 所属 \(kC\) 块的幂等元，则
	\[
		d^u_{\chi\varphi}\ne0
		\quad\Longrightarrow\quad
		c\operatorname{Br}_{\langle u\rangle}^G(b)=c.
	\]
	也就是说，\eqref{b5:eq:generalized-decomposition-numbers}只有位于该 Brauer 像之下的局部块可能贡献非零项。
\end{theorem}
\begin{proofsketch}
	在 \(\chi\) 的表示中选择稳定满格，并按 \(C\) 的块分解其限制。关键的局部分解结论说，不属于 \(\operatorname{Br}_{\langle u\rangle}^G(b)\) 的那些块分量的每个不可分解项，都相对某个不包含 \(\langle u\rangle\) 的子群 \(Q\le C\) 投射。因为 \(u\) 在 \(C\) 中心，也不属于任何 \(C\)-共轭的 \(Q\)。Green 的相对投射特征标消失定理于是使这些项在 \(uv\)、\(v\in C_{p'}\) 处的特征标值为零。对整个诱导模，消失可从无固定陪集看出；对其不可分解直和项仍消失，则是这里真正需要的附加定理，不能由迹的可加性直接推出。这样，\(\chi(uv)\) 完全由 Brauer 像之下的其余块分量贡献；各块的普通特征标约化只产生本块的 Brauer 特征标，再用线性独立性就得到断言。

	这里省去的是满格的局部分解及上述直和项的特征标消失证明，二者均在特征零中进行，不能以 \(k\) 值矩阵迹为零代替。Nagao 的原文用相对迹理想与本原幂等元建立局部分解，在 Section 3 的 Theorem 1 得到所需的特征标等式，再在 Section 4 的 Theorem 2 推出广义分解数的限制；见\cite[Sections 3--4, Theorems 1--2]{Nagao1963GeneralizedDecompositionNumbers}。
\end{proofsketch}

不能将第二主定理读成“非零 Brauer 像所允许的广义分解数都非零”，也不能把 \(\langle u\rangle\) 自动当作原块的缺陷群。它只限制一张广义分解矩阵中哪些块外的条目必须为零。

\term[Brauer-disan-zhudingli]{Brauer 第三主定理}[Brauer's third main theorem]涉及主块的对应。下面陈述并证明其 Sylow 正规化子形式。
\begin{theorem}[Brauer 第三主定理的 Sylow 正规化子形式]\label{b5:thm:brauer-third-main-sylow}
	设 \(G\) 为有限群，\(k\) 为特征 \(p>0\) 的域，并且对 \(G\) 及其全部子群都是分裂域。令 \(P\) 为 Sylow \(p\)-子群，\(N=N_G(P)\)。记 \(b_0(G)\)、\(b_0(N)\) 为包含平凡模的主块幂等元，则 \(b_0(G)\) 的缺陷群为 \(P\)，且
	\[
		\operatorname{Br}_P^G(b_0(G))=b_0(N).
	\]
\end{theorem}
\begin{proof}
	若 \(D\) 是主块的缺陷群，由\cref{b5:cor:block-modules-relative-defect}，平凡模 \(k\) 相对 \(D\) 投射。在 \(\operatorname{End}_k(k)=k\) 上共轭作用平凡，所以 Higman 判据在此就是要求乘数 \([G:D]\) 非零。因 \(D\) 是 \(p\)-群，这等价于 \(D\) 为 Sylow \(p\)-子群。于是第一主定理保证 \(c=\operatorname{Br}_P^G(b_0(G))\) 是 \(kN\) 的块幂等元。

	增广 \(\varepsilon:kG\to k\) 给出平凡模的作用。块幂等元在平凡模上只能为零或一，且恰有主块取一。另一方面，\(P\) 的每个非单点共轭轨道的大小在 \(k\) 中为零，所以对 \(z\in(kG)^P\) 有 \(\varepsilon(\operatorname{Br}_P^G(z))=\varepsilon(z)\)。因此 \(\varepsilon(c)=\varepsilon(b_0(G))=1\)，只能有 \(c=b_0(N)\)。
\end{proof}
这个版本比较主块与 Sylow 正规化子的主块。一般第三主定理还描述任意 \(p\)-子群的中心化子中的主块及其 Brauer 对应；这里没有把这项更强的结论作为上述短证明的推论。

\subsection{Green 对应}
\label{b5:subsec:0806:03}

\term[Green-duiying]{Green 对应}[Green correspondence]选择的是不可分解模。它利用的识别条件并非简单性或维数，而是顶点。为避免混淆共轭范围，下列陈述中“顶点为 \(Q\)”表示可以选取 \(Q\) 本身为顶点；在 \(kH\) 模中，顶点的唯一性只到 \(H\)-共轭。
\begin{theorem}[Green 对应]\label{b5:thm:green-correspondence}
	设 \(G\) 为有限群，\(k\) 为特征 \(p>0\) 的域，\(Q\le G\) 为 \(p\)-子群，\(N_G(Q)\le H\le G\)。所有模均为有限维左模。
	\begin{enumerate}
		\item 若不可分解 \(kG\)-模 \(M\) 的顶点为 \(Q\)，则
		\[
			\operatorname{Res}_H^G M\cong f(M)\oplus X,
		\]
		其中恰有一个不可分解直和项 \(f(M)\) 的顶点为 \(Q\)，重数为一。\(X\) 的每个不可分解直和项都相对某个 \(H\cap{}^xQ\)、\(x\notin H\)，投射。
		\item 若不可分解 \(kH\)-模 \(V\) 的顶点为 \(Q\)，则
		\[
			\operatorname{Ind}_H^G V\cong g(V)\oplus Y,
		\]
		其中恰有一个不可分解直和项 \(g(V)\) 的顶点为 \(Q\)，重数为一。\(Y\) 的每个不可分解直和项都相对某个 \(Q\cap{}^xQ\)、\(x\notin H\)，投射。
		\item 这两个选择在同构类上互逆：\(g(f(M))\cong M\)、\(f(g(V))\cong V\)。
	\end{enumerate}
\end{theorem}
\symidx{\(f(M),g(V)\)：固定子群与顶点后的 Green 对应}
\begin{proof}
	若 \(x\notin H\)，则 \(Q\cap{}^xQ<Q\)，因为相等会迫使 \(x\in N_G(Q)\subseteq H\)。另一个交集 \(H\cap{}^xQ\) 却可能与 \(Q\) 同阶；真正需要的是它不与 \(Q\) 在 \(H\) 中共轭。若 \(H\cap{}^xQ={}^hQ\)、\(h\in H\)，两边同阶便有 \({}^xQ={}^hQ\)，从而 \(h^{-1}x\in N_G(Q)\subseteq H\)，仍是矛盾。

	先从顶点为 \(Q\) 的 \(kH\)-模 \(V\) 出发，取其源 \(S\)，写 \(\operatorname{Ind}_Q^H S=V\oplus W\)。由 Mackey 分解，\(\operatorname{Res}_H^G\operatorname{Ind}_Q^G S\) 中，代表元落在 \(H\) 的一项就是 \(\operatorname{Ind}_Q^H S=V\oplus W\)，其余项都相对 \(H\cap{}^xQ\)、\(x\notin H\)，投射。另一方面，\(\operatorname{Res}_H^G\operatorname{Ind}_H^G V\) 具有由单位双陪集给出的直和项 \(V\)，对 \(W\) 也如此。把这两式相加，并在 Krull--Schmidt 分解中消去 \(V\oplus W\)，便知
	\[
		\operatorname{Res}_H^G\operatorname{Ind}_H^G V\cong V\oplus V',
	\]
	且 \(V'\) 的全部不可分解项都相对上述 \(H\cap{}^xQ\) 投射。它们不可能有顶点 \(Q\)，故这里的 \(V\) 是唯一的该顶点项。

	分解 \(\operatorname{Ind}_H^G V\)，至少有一个不可分解项 \(M\) 的限制含 \(V\)。它相对 \(Q\) 投射；若它的顶点能选在 \(Q\) 的真子群中，限制后的 Mackey 分解会使 \(V\) 也相对一个阶更小的子群投射，矛盾。因此 \(M\) 的顶点是 \(Q\)。对另一项 \(M'\)，其限制的每个不可分解项都在 \(V'\) 中。取 \(M'\) 的顶点 \(Q'\le Q\)，并在限制到 \(H\) 后选取进一步限制到 \(Q'\) 含其源的一项 \(L\)。记 \(L\) 的一个顶点为 \(R\)。由 \(M'\) 相对 \(Q'\) 投射，再对其限制使用 Mackey 分解，得到 \(|R|\le|Q'|\)。另一方面，源是 \(\operatorname{Res}_{Q'}^H L\) 的直和因子；把 \(L\) 的相对 \(R\) 投射性代入 Mackey 分解，便使该源相对某个 \(Q'\cap{}^hR\)、\(h\in H\)，投射。源的顶点为 \(Q'\)，所以 \(Q'\le{}^hR\)；结合阶数不等式，得到 \(Q'={}^hR\)。因此可以选 \(Q'\) 本身为 \(L\) 的顶点。由于 \(L\) 是 \(V'\) 中的一项，某个 \(h\in H\) 满足
	\[
		{}^hQ'\le H\cap{}^yQ\quad(y\notin H),\qquad
		Q'\le Q\cap{}^{h^{-1}y}Q.
	\]
	后者是真子群。这证明诱导方向中的存在、唯一性及余项范围。

	反过来，设 \(M\) 是顶点为 \(Q\) 的不可分解 \(kG\)-模，取源 \(S\)。从 \(M\mid\operatorname{Ind}_Q^G S\) 及诱导传递性，在 \(\operatorname{Ind}_Q^H S\) 中选一项 \(V\)，使 \(M\mid\operatorname{Ind}_H^G V\)。\(V\) 只能有顶点 \(Q\)，否则 \(M\) 的顶点也会变小。由已证部分，\(M=g(V)\)。其限制到 \(H\) 后有一项进一步限制到 \(Q\) 含源 \(S\)，故该项具有顶点 \(Q\)；它只能是前面唯一的 \(V\)。其余项具有第一问指定的相对投射性。两方向的唯一性同时给出互逆关系。
\end{proof}

在限制方向中，余项的顶点不能一概说成“比 \(Q\) 小”：有些可能与 \(Q\) 在 \(G\) 中共轭，却不在 \(H\) 中共轭。诱导方向才由 \(Q\cap{}^xQ<Q\) 得到严格变小。这个区别也说明，不能只比较顶点的阶来挑出限制后的对应项。以上证明采用\cite[Theorem 11.6.4]{Webb2016FiniteGroupRepresentations}的源与 Mackey 分解方法。

\ProseParagraphBreak

Brauer 对应匹配块，Green 对应匹配指定顶点的不可分解模；两者都不是“限制整个模范畴就得到等价”的断言。特别地，诱导和限制一般仍带有额外直和项，简单模也可能对应到不简单的模。下面的计算直接追踪这些额外项。

\subsection{小群的块与对应计算}
\label{b5:subsec:0806:04}

先取特征二的 \(k\) 与 \(G=S_3\)，沿用\cref{b5:exm:s3-two-blocks}的 \(r=(123)\)、\(s=(12)\)、\(e=1+r+r^2\)、\(f=r+r^2\)。对 \(P=\langle s\rangle\)，有 \(N_G(P)=C_G(P)=P\)。已算出的
\[
	\operatorname{Br}_P^G(e)=1,\qquad \operatorname{Br}_P^G(f)=0
\]
现在给出实际的 Brauer 对应：缺陷为 \(P\) 的主块 \(kGe\) 对应 \(kP\) 的唯一块。矩阵块 \(kGf\) 的缺陷为一，不能把它也配给 \(kP\)；对它应取 \(N_G(1)=G\)，此时对应为恒等。

\ProseParagraphBreak

\(kP\cong k[t]/(t^2)\) 只有两个不可分解模 \(k\)、\(kP\)，其中只有 \(k\) 的顶点为 \(P\)。另一方面，\(kGe\cong k[t]/(t^2)\)，而矩阵块只有投射的简单模 \(S_2\)；所以 \(kG\) 中也只有平凡模的顶点为 \(P\)。Green 对应在这里把两边的平凡模配对，但诱导仍不等于这个对应项：
\[
	\operatorname{Ind}_P^{S_3}k\cong k\oplus S_2.
\]
验证可在三个陪集上的置换模 \(k^3\) 中进行。常向量构成平凡子模，坐标和为零的平面构成补模，因为 \(3=1\ne0\)；这个平面正是前面构造的 \(S_2\)。额外项 \(S_2\) 投射，顶点为一，符合 \(P\cap{}^xP=1\)、\(x\notin P\) 的结论。

\begin{example}\label{b5:exm:s3-three-block-local}
	设 \(k\) 为特征三的域，\(G=S_3=\langle r,s\mid r^3=s^2=1,\ srs=r^{-1}\rangle\)，\(P=\langle r\rangle\)。令 \(S_+\)、\(S_-\) 为平凡模和符号模，\(t=r-1\)、\(e_\pm=(1\pm s)/2\)。则 \(kG\) 只有主块，其缺陷群为 \(P\)，且
	\[
		J(kG)=(t),\quad J(kG)^3=0,\qquad P_\pm=kGe_\pm
	\]
	分别是 \(S_\pm\) 的投射覆盖。\(P_+\) 从顶到基座的根层为 \(S_+,S_-,S_+\)，\(P_-\) 的相应根层为 \(S_-,S_+,S_-\)。因此
	\[
		C=\begin{pmatrix}2&1\\1&2\end{pmatrix}.
	\]
\end{example}
\begin{proof}
	因为 \(t^3=r^3-1=0\)，且 \(sts=r^{-1}-1=-t+t^2\)，所以由 \(t\) 生成的双侧理想是 \(tkP\oplus tkP s\)，它的三次幂为零。商代数为 \(k(G/P)\cong kC_2\cong k\times k\)，是半单的；故该幂零理想恰为 \(J(kG)\)，商的两个简单模正是 \(S_+,S_-\)。

	\(e_+,e_-\) 是正交幂等元，和为一，所以 \(P_\pm\) 投射，且各以 \(e_\pm,te_\pm,t^2e_\pm\) 为基。\(J^jP_\pm/J^{j+1}P_\pm\) 由 \(t^je_\pm\) 的类张成；\(r=1+t\) 在它上面作用为一，而
	\[
		s t^je_\pm=(-t+t^2)^j s e_\pm
		\equiv \pm(-1)^j t^je_\pm\pmod{J^{j+1}P_\pm}.
	\]
	于是得到三层交替的符号。它们的顶均简单，故由\cref{b5:prop:projective-cover-top}，\(P_\pm\) 是相应投射覆盖且不可分解。一个不可分解模只能属于一个块，且它的全部合成因子都在同一块中。\(P_+\) 同时含两个简单模，所以这两个简单模属于同一块；每个非零块都含简单模，故全代数只有一个块，幂等元为一。对 \(P\) 有 \(\operatorname{Br}_P^G(1)=1\)，且 \(P\) 已为 Sylow 子群，因而它就是缺陷群。根层中的次数给出 Cartan 矩阵。
\end{proof}

在这个例子中 \(P\) 正规，故 \(N_G(P)=G\)。以 \(P\) 为顶点或缺陷群的对应都留在 \(G\) 本身，不能仅因“对应到了正规化子”就期待群阶下降。同一个群在特征二和特征三中的块数、缺陷群及对应计算均不同；决定这些区别的是特征与子群的位置，而不是普通特征标表的行数。

% * 状态：初稿已完成；尚未正式审核。

% !TeX root = ../../Book5.tex
% 纲要标题：### 8.7 模表示的计算与核验
% * 状态：初稿已完成；新增范围尚未正式审核。
% #### 8.7.1 Brauer 特征标与分解矩阵的计算数据
% #### 8.7.2 MeatAxe 的输入与输出
% #### 8.7.3 GAP 计算与独立核验

\section{模表示的计算与核验}
\label{b5:sec:0807}

给出群生成元的矩阵以后，可以求不变子空间、分解出简单因子，再计算这些因子的 Brauer 特征标。但这几个任务保留的信息不同：一张特征标表不记录扩张的连接方式，一串因子维数也不能证明原模已经分成直和。要使计算结果可以检查，必须同时保存工作域、生成元次序、换基矩阵及所用的简单模标识。

\subsection{Brauer 特征标与分解矩阵的计算数据}
\label{b5:subsec:0807:01}

固定\cref{b5:def:p-modular-system}中的分裂 \(p\)-模系统，并满足\cref{b5:def:brauer-character}的单位根及复嵌入条件。记 \(g_1,\ldots,g_\ell\) 为 \(p\)-正则共轭类代表，\(S_1,\ldots,S_\ell\) 为全部不同构简单模；二者数目相同是\cref{b5:thm:brauer-character-basis}的结论。普通不可约特征标记为 \(\chi_1,\ldots,\chi_c\)，不可约 Brauer 特征标记为 \(\varphi_1,\ldots,\varphi_\ell\)。计算中首先要固定这两组行标，以及同一组共轭类列标。

\begin{proposition}\label{b5:prop:decomposition-data-recovery}
	设 \(G\) 为有限群，\((K,\mathcal O,k)\) 为分裂 \(p\)-模系统，\(K\) 含全部 \(|G|\) 次单位根，并固定定义 Brauer 特征标所用的复嵌入。令 \(g_1,\ldots,g_\ell\) 为全部 \(p\)-正则类代表，\(\chi_1,\ldots,\chi_c\) 为全部普通不可约特征标，\(S_1,\ldots,S_\ell\) 为全部不同构简单 \(kG\)-模，\(\varphi_i\) 为 \(S_i\) 的 Brauer 特征标。令
	\[
		F_{\alpha j}=\chi_\alpha(g_j),\qquad
		\Phi_{ij}=\varphi_i(g_j).
	\]
	则 \(\Phi\) 可逆，分解矩阵由
	\[
		F=D\Phi,\qquad D=F\Phi^{-1}
	\]
	唯一确定。若 \(d_i=\dim_kS_i\)，则第 \(\alpha\) 行还满足
	\[
		\chi_\alpha(1)=\sum_iD_{\alpha i}d_i.
	\]
\end{proposition}
\begin{proof}
	不可约 Brauer 特征标构成全部 \(p\)-正则类函数的一组基，所以其值矩阵 \(\Phi\) 可逆。\cref{b5:prop:ordinary-brauer-restriction}在每个 \(g_j\) 上给出 \(\chi_\alpha(g_j)=\sum_iD_{\alpha i}\varphi_i(g_j)\)，这正是矩阵等式。最后在单位元处求值，普通和 Brauer 特征标的值分别为相应维数，得到维数等式。
\end{proof}

这里的计算发生在特征零的分圆数域中，可以用精确的代数数运算进行。若把 \(\Phi\) 换成模特征中的迹矩阵，即使矩阵仍可逆，算出的也只是整数重数模 \(p\) 后的信息。反过来，一张候选表通过了整数性、非负性与维数检查，也还需要证明所列模确实简单且没有遗漏；这些前提不能由矩阵乘法补出。

\ProseParagraphBreak
例如对 \(S_3\) 和 \(p=3\)，按平凡、符号、标准表示排列普通行，按平凡、符号模排列简单模，并按单位元、换位排列列。\cref{b5:exm:brauer-s3-characteristic-three}给出
\[
	F=\begin{pmatrix}1&1\\1&-1\\2&0\end{pmatrix},
	\qquad
	\Phi=\begin{pmatrix}1&1\\1&-1\end{pmatrix},
	\qquad
	D=\begin{pmatrix}1&0\\0&1\\1&1\end{pmatrix}.
\]
两个简单模均为一维，所以 \(D\) 的行和依次为 \(1,1,2\)。\cref{b5:exm:s3-three-block-local}另从投射覆盖的根层算出
\[
	C=\begin{pmatrix}2&1\\1&2\end{pmatrix}=D^{\mathsf T}D.
\]
两条计算的对象不同：前者要知道投射模的合成因子，后者只使用普通表示约化后的重数；不能先把 \(D^{\mathsf T}D\) 当作任意有限维代数的 Cartan 矩阵定义。

\subsection{MeatAxe 的输入与输出}
\label{b5:subsec:0807:02}

矩阵输入通常取一个有限域 \(k=\mathbb F_q\)、一个正整数 \(d\)，以及群生成元 \(g_1,\ldots,g_s\) 在 \(k^d\) 上的可逆矩阵 \(A_1,\ldots,A_s\)。还须确认这些矩阵满足群的定义关系；否则它们至多给出自己所生成的矩阵群的表示。此处一律以列向量表示左模，即 \(g_iv=A_iv\)。

\ProseParagraphBreak
\term[MeatAxe]{MeatAxe}[MeatAxe]是一组有限域上的模计算方法，用于寻找不变子空间、判断不可约性和计算合成因子等。不同程序可用不同的搜索与不可约性判据。一个基本操作是从少量向量出发，反复加上它们在生成矩阵下的像；这个操作本身有可以完全说明的停止条件。

\begin{proposition}\label{b5:prop:module-spin-up}
	设 \(k\) 为域，有限群 \(G=\langle g_1,\ldots,g_s\rangle\) 在 \(M=k^d\) 上的矩阵为 \(A_i\)，\(W_0\) 为给定子空间。递归令
	\[
		W_{j+1}=W_j+A_1W_j+\cdots+A_sW_j.
	\]
	若 \(W_{j+1}=W_j\)，则此后不再变化，所得空间是包含 \(W_0\) 的最小 \(kG\)-子模。每次尚未停止时维数严格增加，所以至多经过 \(d-\dim W_0\) 次严格增加就会停止。
\end{proposition}
\begin{proof}
	维数相等的相邻空间因包含关系而相等；此时每个 \(A_i\) 保持该空间。\(A_i\) 可逆，其在有限维子空间上的限制单射，故像等于该子空间，逆矩阵也保持它。于是所有群元素都保持它，得到子模。任意包含 \(W_0\) 的子模由归纳包含每个 \(W_j\)，说明所得子模最小。最后，维数是上界为 \(d\) 的整数，尚未停止时至少增加一。
\end{proof}

在实际计算中，把当前基和所有 \(A_i\) 作用后的向量排在一起消元，就能取得下一轮的基。发现一个维数严格介于 \(0\) 与 \(d\) 之间的结果，便找到一个非平凡子模；然而某一个非零向量生成了整个模，并不说明模不可约，正则模本身就由单位元生成。在有限域上穷尽非零向量并逐个作此计算，确实能判断不可约性，但向量数为 \(q^d-1\)，通常很快变得不可行。MeatAxe 的作用正在于使用更有效的搜索和判据；随机尝试没有找到子模，也不能单凭这一事实宣布不可约。

\begin{proposition}\label{b5:prop:matrix-submodule-certificate}
	设 \(k\) 为域，有限群 \(G=\langle g_1,\ldots,g_s\rangle\) 在 \(k^d\) 上的矩阵为 \(A_i\)。若 \(C\in M_{d\times r}(k)\) 列满秩，则其列空间为子模，当且仅当存在 \(B_i\in M_r(k)\) 使
	\[
		A_iC=CB_i\qquad(1\le i\le s).
	\]
	把 \(C\) 的列补成一组基，所得换基矩阵记为 \(T\)。上述条件等价于每个矩阵均有形式
	\[
		T^{-1}A_iT=
		\begin{pmatrix}B_i&Z_i\\0&D_i\end{pmatrix}.
	\]
	其中 \(B_i,D_i\) 分别给出子模和商模的作用。这个子模存在不变补空间，当且仅当有同一个 \(X\in M_{r\times(d-r)}(k)\) 满足
	\[
		B_iX+Z_i=XD_i\qquad(1\le i\le s).
	\]
\end{proposition}
\begin{proof}
	列空间不变恰好表示 \(A_iC\) 的每一列都可用 \(C\) 的列表示；系数矩阵就是 \(B_i\)，且因列独立而唯一。补基以后，下方左块为零正是这一不变性，上方左块给出子模作用，下方右块给出商作用。

	在此分块坐标中，任意与第一个直和因子互补的线性子空间，唯一写成某个线性映射的图像
	\[
		\Gamma_X=\{\,(Xv,v)\mid v\in k^{d-r}\,\}.
	\]
	生成矩阵把 \((Xv,v)\) 送到 \((B_iXv+Z_iv,D_iv)\)。它仍位于同一图像中，当且仅当 \(B_iX+Z_i=XD_i\)。所有生成元都满足这些等式，才得到群不变补空间。
\end{proof}

因此，一份子模计算可以附上 \(C,B_i\)，供人直接检查秩和矩阵乘法；一份合成列计算可以附上适应该列的基，以及各个简单商的识别。上三角的分块矩阵保留了 \(Z_i\)；只列出对角块会丢掉扩张信息。要求实际直和分解时，还需解出同一个 \(X\)，或给出另一份能够消去所有 \(Z_i\) 的换基。

\begin{proposition}\label{b5:prop:full-matrix-algebra-certificate}
	设 \(k\) 为域，有限群 \(G\) 在 \(k^d\) 上的生成矩阵为 \(A_1,\ldots,A_s\)。若这些矩阵生成的含幺 \(k\) 子代数等于 \(M_d(k)\)，则表示在任意扩域后仍不可约。特别地，若能给出 \(d^2\) 个由生成矩阵相乘所得的字，其展平后的列向量线性无关，就得到这一结论。
\end{proposition}
\begin{proof}
	扩域后这组矩阵仍生成全部矩阵代数。若一个不变子空间含有非零向量 \(v\)，任选 \(v\) 的一个非零坐标，利用矩阵单位可将它送到每个坐标方向；于是该子空间包含全部标准基。故任一扩域上的非零子模都是全空间。第二个条件给出 \(M_d(k)\) 中 \(d^2\) 个独立元素，它们已经张成全部空间，因而足够。
\end{proof}

在 \(k=\mathbb F_2\) 上取
\[
	A=\begin{pmatrix}0&1\\1&1\end{pmatrix},
	\qquad
	B=\begin{pmatrix}0&1\\1&0\end{pmatrix}.
\]
有 \(A^3=B^2=I\)、\(BAB=A^{-1}\)，故得到 \(S_3\) 的二维表示。将 \(I,A,B,AB\) 按同一顺序展平，所得四阶矩阵的行列式为 \(1\)，所以它们张成 \(M_2(\mathbb F_2)\)。这证明该模绝对不可约，而不只是说明某次子模搜索没有成功。

\subsection{GAP 计算与独立核验}
\label{b5:subsec:0807:03}

GAP 的 MeatAxe 接口用行向量右乘矩阵；本节的左模使用列向量。因此输入本节的 \(A_i\) 时，应传入 \(A_i^{\mathsf T}\)。对于一个不交换的乘积，转置同时反转矩阵次序；这只是把同一个线性算子作用改写为行坐标，不能把转置后的列表又按原来的列作用约定读回。

\ProseParagraphBreak
取 \(k=\mathbb F_3\)、\(G=S_3=\langle r,s\mid r^3=s^2=1,\ srs=r^{-1}\rangle\)。三点置换模中令
\[
	U=\{\,(x,y,z)\in k^3\mid x+y+z=0\,\},\qquad
	u=(1,-1,0),\quad v=(0,1,-1).
\]
在基 \(u,v\) 下，\(r=(123)\)、\(s=(12)\) 的列作用矩阵为
\begin{equation}\label{b5:eq:s3-modular-matrix-input}
	A=\begin{pmatrix}0&2\\1&2\end{pmatrix},
	\qquad
	B=\begin{pmatrix}2&1\\0&1\end{pmatrix}.
\end{equation}
直接相乘给出 \(A^3=B^2=I\)、\(BAB=A^2\)，所以这些输入满足群关系。相应计算可采用 GAP 4.16.1 的 MeatAxe 接口\footnote{GAP 4.16.1 官方参考手册，第69章，69.1-1、69.5-1及69.7-10：\url{https://docs.gap-system.org/doc/ref/chap69_mj.html}。模由生成矩阵建立，合成因子以模对象列表返回；此处读取各因子的维数。}：
\begin{verbatim}
F := GF(3);;
a := [[0,2],[1,2]] * One(F);;
b := [[2,1],[0,1]] * One(F);;
m := GModuleByMats(
  [TransposedMat(a),TransposedMat(b)], F);;
MTX.IsIrreducible(m);
cf := MTX.CompositionFactors(m);;
List(cf, MTX.Dimension);
\end{verbatim}
这里第一个判断应为 \texttt{false}，最后的维数列表应为 \texttt{[ 1, 1 ]}；因子的次序和选取的基不必固定。维数列表只说明有两个一维合成因子，并没有给出它们的名称，也没有给出两个因子在原模中的连接方式。

\ProseParagraphBreak
下面直接从输入得到可检查的子模和商模。令 \(w=u+2v=(1,1,1)\)，取
\[
	T=\begin{pmatrix}1&0\\2&1\end{pmatrix},\qquad
	T^{-1}AT=\begin{pmatrix}1&2\\0&1\end{pmatrix},\qquad
	T^{-1}BT=\begin{pmatrix}1&1\\0&2\end{pmatrix}.
\]
第一条基向量生成平凡模 \(kw\)，商上的 \(r,s\) 分别作用为 \(1,-1\)，故商是符号模。两者均为一维，已经简单，于是得到合成列 \(0\subset kw\subset U\)。若有不变补直线，\cref{b5:prop:matrix-submodule-certificate}中的未知矩阵 \(X\) 只是一个标量 \(x\)。对 \(r\) 的方程却为 \(x+2=x\)，无解。因此
\[
	0\longrightarrow S_+\longrightarrow U\longrightarrow S_-\longrightarrow0
\]
不分裂；GAP 的合成因子列表不能改写成 \(U\cong S_+\oplus S_-\)。

\ProseParagraphBreak
这个例子还区分了自同态计算与不可约性判断。在上述三角基中，和 \(A\) 交换的矩阵具有形式
\[
	X=\begin{pmatrix}a&b\\0&a\end{pmatrix}.
\]
再要求它与 \(B\) 交换，便得 \(b=0\)。所以 \(\operatorname{End}_{kG}(U)=k\)，但 \(U\) 明明含有子模 \(kw\)。Schur 引理给出的“简单模的自同态环是除环”不能反向使用；检验自同态空间只有一维，也不能代替不可约性证明。

\ProseParagraphBreak
最后，在 \(3\)-正则元素 \(s\) 处，矩阵 \(B\) 的特征值为 \(1,-1\)，按固定模系统提升后仍为 \(1,-1\)，所以 \(\varphi_U(s)=0\)、\(\varphi_U(1)=2\)。它与 \(S_+\oplus S_-\) 的 Brauer 特征标相同；非分裂性则保存在上三角块和补空间方程中。把特征标、合成列与这些矩阵一起保留，才同时说明了计算得到哪些简单成分，以及它们怎样组成原模。
% ---


\begin{chaptersummary}
	在模特征下，群代数的根和模的根层保留了非半单部分的信息，但连全部根层也未必确定原模。循环 \(p^r\) 群的群代数是 \(k[t]/(t^{p^r})\)，其不可分解模由幂零链的长度区分；它们虽然只有一种简单合成因子，直和结构仍不相同。投射覆盖按简单顶排列，Cartan 矩阵记录投射模中的简单重数，块则把整个模范畴分成彼此没有同态和扩张的部分。循环 \(p\) 群上的短正合扩张还可由投射覆盖的核到目标模的稳定同态直接计算，其中模掉经过投射模的同态恰好对应改变扩张提升所消去的参数。

	\ProseParagraphBreak

	Brauer 特征标只在 \(p\)-正则元素上定义，把剩余域中的特征值提升到特征零后求和。这样避免了 \(k\) 值迹把整数重数约化模 \(p\) 的损失；在本章的分裂条件下，整个 Brauer 特征标函数确定简单合成因子及其整数重数，但不能区分具有相同合成因子的不同扩张。分解数描述普通不可约表示约化后的简单成分；块的缺陷群则借助 Brauer 映射把块与 \(p\) 子群联系起来。

	\ProseParagraphBreak
	群代数的对称配对使有限维投射模同时内射，因而合冲模在稳定范畴中成为等价，一阶扩张由 \(\underline{\operatorname{Hom}}(\Omega M,N)\) 给出。相对迹的恒等映射判据把诱导直和因子与子群联系起来，顶点只确定到共轭，源还须记录子群上的模。Brauer 对应处理块，Green 对应处理不可分解模，两者均有指定局部子群与顶点或缺陷群的条件，不能由此推出任意两块的模范畴等价。

	\ProseParagraphBreak
	计算时，特征标表、分解矩阵及合成因子列表首先记录数值；不变子空间和分裂性则需要矩阵证据。等式 \(A_iC=CB_i\) 核验子模，换基后的非对角块记录扩张，补空间存在与否由一组共同的线性方程决定。这样既可以利用软件找到候选结果，也能说明结果为何成立，以及它还没有确定哪些结构。
\end{chaptersummary}

\BookExercises

\begin{optionalexercise}\label{b5:ex:08-jordan-data}
	设 \(\operatorname{char}k=2\)，\(C_8=\langle s\rangle\) 在十维空间 \(V\) 上作用。令 \(T=s-1\)，已知
	\[
		\dim\ker T^j=3,5,7,8,9,10\qquad(j=1,\ldots,6).
	\]
	求不可分解直和分解，以及顶、基座、不动空间和 \(T^5\) 的秩。
\end{optionalexercise}
\begin{solution}
	相邻核维数之差为 \(3,2,2,1,1,1\)，分别给出长度至少为 \(1,\ldots,6\) 的幂零块数。故块长为 \(6,3,1\)， \(V\cong M_6\oplus M_3\oplus M_1\)。 三个直和项各有一维顶和一维基座，故两者维数均为 \(3\)。不动空间为 \(\ker(s-1)=\ker T\)，也为三维。最后 \(\operatorname{rank}T^5=10-9=1\)。只有简单合成因子及其总数不能恢复这些块长，而各次核维数可以。
\end{solution}

\begin{optionalexercise}\label{b5:ex:08-norm-invariants}
	设 \(G=C_m\)，其中 \(m=p^r\)，\(\operatorname{char}k=p\)，\(t=s-1\)。对 \(M_d=k[t]/(t^d)\)、\(1\leq d\leq m\)，计算不动空间、余不动空间 \(M_d/tM_d\) 及两者之间的自然映射。再证明群的求和算子 \(\sum_{j=0}^{m-1}s^j\) 等于乘 \(t^{m-1}\)，并计算它在这些模上的像。
\end{optionalexercise}
\begin{solution}
	不动空间是 \(kt^{d-1}\)，余不动空间由 \(1\) 的类生成。自然映射由包含不动空间再取商给出：\(d=1\) 时为同构，\(d>1\) 时因 \(t^{d-1}\in tM_d\) 而为零。 在多项式环中，
	\[
		\sum_{j=0}^{m-1}(1+t)^j =\frac{(1+t)^m-1}{t}=t^{m-1},
	\]
	其中商式是几何级数恒等式，而特征 \(p\) 及 \(m=p^r\) 给出 \((1+t)^m=1+t^m\)。故在 \(d<m\) 时求和算子为零；在 \(d=m\) 时，其像为 \(kt^{m-1}\)，并诱导从一维余不动空间到一维不动空间的同构。
\end{solution}

\begin{optionalexercise}\label{b5:ex:08-elementary-abelian-radical}
	设 \(P=C_p\times C_p=\langle s,u\rangle\)，\(\operatorname{char}k=p\)。证明
	\[
		kP\cong k[x,y]/(x^p,y^p),\qquad x=s-1,\ y=u-1.
	\]
	求根 \(J\)、每层 \(J^i/J^{i+1}\) 的维数、最小的 \(N\) 使 \(J^N=0\)，以及正则模的基座。
\end{optionalexercise}
\begin{solution}
	两生成元交换，且 \(x^p=y^p=0\)。所给商以 \(x^a y^b\)、\(0\leq a,b<p\) 为基，共 \(p^2\) 个；到 \(kP\) 的映射满射且维数相等，所以同构。 根为 \(J=(x,y)\)，第 \(i\) 层以全次数为 \(i\) 的上述单项式为基。因此层维数依次为
	\[
		\dim J^i/J^{i+1}=
		\begin{cases}
			i+1,&0\leq i\leq p-1,\\
			2p-1-i,&p\leq i\leq2p-2,\\
			0,&i\geq2p-1.
		\end{cases}
	\]
	最高非零单项式为 \(x^{p-1}y^{p-1}\)，故 \(J^{2p-2}\neq0\)、\(J^{2p-1}=0\)。一个元素同时被 \(x,y\) 零化，当且仅当只含这个最高单项式；因此基座为 \(kx^{p-1}y^{p-1}\)。
\end{solution}

\begin{optionalexercise}\label{b5:ex:08-elementary-abelian-family}
	设 \(k\) 为特征 \(p\) 的无限域。记 \(V_\lambda\) 为 \(C_p\times C_p\) 的二维表示，两生成元分别作用为 \(I+N\)、\(I+\lambda N\)，其中 \(N^2=0,N\neq0\)。求任意参数 \(\lambda,\mu\) 之间的同态空间及非零同态的秩。由此说明不同构的不可分解模之间也可以有非零同态。
\end{optionalexercise}
\begin{solution}
	两生成元的矩阵交换且 \(p\) 次幂为恒等，故确为表示。取基使
	\(N=\left(\begin{smallmatrix}0&1\\0&0\end{smallmatrix}\right)\)。
	同态矩阵先须与 \(N\) 交换，因而为 \(f=aI+bN\)；第二生成元的条件为 \((\lambda-\mu)fN=0\)，而 \(fN=aN\)。 所以当 \(\lambda=\mu\) 时，同态空间为 \(kI+kN\)，其中 \(a\neq0\) 的映射可逆、\(a=0,b\neq0\) 的映射秩为一。当 \(\lambda\neq\mu\) 时，同态空间仅为 \(kN\)，所有非零映射秩为一，且两个这样的跨参数映射的复合为零。 每个自同态环同构于局部环 \(k[t]/(t^2)\)，所以模不可分解；不同参数之间没有可逆同态，却仍有一维同态空间。非半单情形下，不能把简单模的 Schur 引理直接改成不可分解模的断言。
\end{solution}

\begin{optionalexercise}\label{b5:ex:08-hom-through-projectives}
	设 \(p\) 为素数，\(k\) 为特征 \(p\) 的域，\(r\ge1\) 为整数，\(m=p^r\ge4\)、\(A=k[t]/(t^m)\)，并令 \(M_d=A/(t^d)\)。计算稳定自同态环 \(\underline{\operatorname{End}}_A(M_2)\)。再取同态
	\[
		i:M_1\to M_2,\quad i(1)=t,\qquad q:M_2\to M_1,\quad q(1)=1,
	\]
	计算两个方向的稳定同态空间及 \(qi,iq\) 的稳定类。说明稳定范畴中非零态射的复合未必非零。
\end{optionalexercise}
\begin{solution}
	每个 \(M_2\) 的自同态都是乘以 \(a+bt\)，复合对应多项式相乘。由\cref{b5:prop:cyclic-stable-hom-extensions}，经过投射模的映射由 \(t^{m-2}M_2\) 给出；因 \(m\ge4\)，这个空间为零。因此 \(\underline{\operatorname{End}}_A(M_2)\cong k[\epsilon]/(\epsilon^2)\)，其中 \(\epsilon\) 是乘 \(t\) 的稳定类。 从 \(M_1\) 到 \(M_2\) 的同态由被 \(t\) 零化的元素 \(kt\) 给出，故由 \(i\) 生成；反向同态由一的任意像给出，故由 \(q\) 生成。两空间中经过投射模的映射分别对应 \(t^{m-1}M_2\)、\(t^{m-2}M_1\)，也都为零，所以两个稳定类均非零。 但 \(qi=0\)，而 \(iq\) 是 \(M_2\) 上的乘 \(t\) 映射，稳定类为非零的 \(\epsilon\)。即使两端的同态空间都是一维，两个方向的复合仍可有不同结果。
\end{solution}

\begin{optionalexercise}\label{b5:ex:08-cyclic-extension-parameters}
	设 \(k\) 为特征 \(2\) 的域，\(A=k[t]/(t^4)\)、\(M_d=A/(t^d)\)。对扩张 \(0\to M_2\to E_c\to M_1\to0\)，取指定子模的生成元 \(u\) 与商生成元的提升 \(v\)，令 \(tv=cu\)、\(c\in k\)。说明每个扩张类都可唯一这样表示，并从 Baer 和的拉回、推出构造算出 \([E_c]+[E_d]\)。确定中间项 \(E_c\) 的同构类型，说明两个中间项都不可分解的扩张之和仍可分裂。
\end{optionalexercise}
\begin{solution}
	由\cref{b5:prop:cyclic-stable-hom-extensions}，这里的参数空间是 \(M_2/tM_2\cong k\)。若最初 \(tv=cu+btu\)，把 \(v\) 改为 \(v-bu\) 就消去第二项，而进一步改变提升不能改变常数项 \(c\)。因此 \(c\) 唯一确定扩张类。 两个扩张的直和沿对角映射 \(M_1\to M_1\oplus M_1\) 拉回后，可取 \((v_c,v_d)\) 为商生成元的提升；乘 \(t\) 得到核中的 \((cu_c,du_d)\)。再沿 \(M_2\oplus M_2\to M_2\) 的相加映射推出，这个向量成为 \((c+d)u\)。因此 \([E_c]+[E_d]=[E_{c+d}]\)。 \(c=0\) 时为 \(M_2\oplus M_1\)；\(c\ne0\) 时，\(v,tv,t^2v\) 是一组基，且 \(t^2v=ctu\ne0\)、\(t^3v=0\)，所以 \(E_c\cong M_3\)。特别地，\([E_1]+[E_1]=[E_0]\)，两条非分裂扩张之和分裂。扩张的加法运算作用于带有指定核、商的扩张类，不是对中间模取直和。
\end{solution}

\begin{optionalexercise}\label{b5:ex:08-c4-extension-middle}
	设 \(k\) 为特征 \(2\) 的域，\(A=kC_4\cong k[t]/(t^4)\)，\(M_d=A/(t^d)\)、\(1\le d\le4\)。 对左 \(A\)-模短正合列 \(0\to M_2\to E\to M_2\to0\)，记指定子模的生成元为 \(u\)、商模生成元的提升为 \(v\)，把扩张参数写成 \(t^2v=\alpha u+\beta tu\)。分别讨论 \(\alpha\neq0\)、\(\alpha=0,\beta\neq0\)、\(\alpha=\beta=0\) 时中间项的不可分解直和类型。
\end{optionalexercise}
\begin{solution}
	由\cref{b5:prop:cyclic-stable-hom-extensions}，\(m=4,a=b=2\) 时整个 \(M_2\) 都可作为参数，且 \(t^2M_2=0\)，所以参数确为 \(\alpha+\beta t\)，不同参数给出不同扩张类。 若 \(\alpha\neq0\)，有 \(t^3v=\alpha tu\neq0\)。向量 \(v,tv,t^2v,t^3v\) 线性无关，故 \(E\cong M_4\)。 若 \(\alpha=0,\beta\neq0\)，有 \(t^2v=\beta tu\neq0\)、\(t^3v=0\)，故四维模含有一个长度三的幂零块，剩余一维给出 \(E\cong M_3\oplus M_1\)。 若参数为零，\(v\) 生成一个与指定子模互补的 \(M_2\)，故 \(E\cong M_2\oplus M_2\)。 扩张类的两个标量参数与中间项的三种同构类型记录的是不同的分类要求。
\end{solution}

\begin{optionalexercise}\label{b5:ex:08-projective-sylow-dimension}
	设 \(G\) 为有限群，\(k\) 的特征为 \(p\)，\(P\) 为一个 Sylow \(p\) 子群。证明每个有限维投射 \(kG\)-模的维数都被 \(|P|\) 整除，并给出维数可被 \(|P|\) 整除却不投射的模。
\end{optionalexercise}
\begin{solution}
	将 \(G\) 分为互不相交的集合 \(Pg\)，可知左 \(kP\)-模 \(kG\) 自由。因此投射 \(kG\)-模限制到 \(P\) 后，是有限自由 \(kP\)-模的直和因子，仍投射。 \(kP\) 是局部代数，唯一简单模为 \(k\)，其投射覆盖为 \(kP\)。由\cref{b5:thm:split-projective-covers}，有限维投射 \(kP\)-模都是有限自由模，维数为 \(|P|\) 的倍数。 反例取 \(G=C_p\) 及 \(p\) 维平凡模 \(k^{\oplus p}\)。如果它投射，便应为秩一自由 \(kC_p\)-模；但自由模的生成元 \(s\) 作用并非恒等，而平凡模中作用恒等，矛盾。
\end{solution}

\begin{optionalexercise}\label{b5:ex:08-projective-tensor}
	设 \(G\) 为有限群，\(k\) 为任意域，\(P\) 为有限维投射 \(kG\)-模，\(V\) 为任意有限维 \(kG\)-模。证明带对角群作用的 \(P\otimes_k V\) 仍投射。
\end{optionalexercise}
\begin{solution}
	先看 \(kG\otimes_kV\)。从对角作用的版本到只在第一因子作用的版本，定义
	\[
		g\otimes v\longmapsto g\otimes g^{-1}v.
	\]
	它的逆为 \(g\otimes v\mapsto g\otimes gv\)。对 \(h\in G\)，对角作用后的像为 \(hg\otimes(hg)^{-1}(hv)=hg\otimes g^{-1}v\)，恰为目标上的左作用，故这是模同构。目标按 \(V\) 的一组基为 \(\dim V\) 份正则模，因而自由。 取 \(P\) 在某个 \((kG)^r\) 中的直和补，张量保持此直和分解，所以 \(P\otimes V\) 为自由模 \((kG\otimes V)^r\) 的直和因子，因而投射。
\end{solution}

\begin{optionalexercise}\label{b5:ex:08-two-block-tensor}
	设 \(G=C_p\)，\(\operatorname{char}k=p\)，\(M_2\) 中生成元的矩阵为 \(I+N\)，其中 \(N^2=0,N\neq0\)。分别在 \(p=2\) 和 \(p>2\) 时求对角作用的 \(M_2\otimes_kM_2\) 的不可分解分解。
\end{optionalexercise}
\begin{solution}
	取基 \(u,v\) 满足 \(Nu=0,Nv=u\)。记张量基为 \(a=u\otimes u,b=u\otimes v,c=v\otimes u,d=v\otimes v\)。 若 \(D=s-1\)，则
	\[
		Da=0,\qquad Db=Dc=a,\qquad Dd=b+c+a.
	\]
	因此 \(\operatorname{rank}D=2\)，且 \(D^2d=2a\)，其余基向量的二次像为零，\(D^3=0\)。 若 \(p=2\)，\(D^2=0\)，四维空间上的秩二幂零算子由两个长度二块组成，故张量积为 \(M_2^{\oplus2}\)。 若 \(p>2\)，二次幂秩为一，块长为 \(3,1\)，故张量积为 \(M_3\oplus M_1\)。\(p=2\) 时 \(M_2=kC_2\) 投射，其张量积的结果也符合上一题。
\end{solution}

\begin{optionalexercise}\label{b5:ex:08-nonsymmetric-cartan}
	设 \(k\) 为域，\(A=T_3(k)/(kE_{13})\)，其中 \(T_3(k)\) 是上三角矩阵代数。以同一记号表示矩阵单位在商中的像，取 \(e_i=E_{ii}\) 和简单模 \(S_i=Ae_i/J(A)e_i\)。按本章“行是简单模、列是投射覆盖”的约定计算 Cartan 矩阵，说明一般有限维代数的 Cartan 矩阵不必对称。 再计算各不可分解投射模的根与基座。
\end{optionalexercise}
\begin{solution}
	\(kE_{13}\) 是双边理想，商中的严格上三角部分 \(J=kE_{12}\oplus kE_{23}\) 平方为零，商 \(A/J\cong k^3\) 半单，故这就是根。\(P_i=Ae_i\) 是投射模，顶为 \(S_i\)，所以是相应的不可分解投射覆盖。\(P_1=kE_{11}=S_1\)；\(P_2\) 以 \(E_{12},E_{22}\) 为基，根 \(kE_{12}\cong S_1\)，顶为 \(S_2\)；\(P_3\) 以 \(E_{23},E_{33}\) 为基，根 \(kE_{23}\cong S_2\)，顶为 \(S_3\)。因此
	\[
		C_A=\begin{pmatrix}1&1&0\\0&1&1\\0&0&1\end{pmatrix}.
	\]
	它不对称。三个投射模的根分别为 \(0,S_1,S_2\)；被 \(J\) 零化的部分就是基座，所以基座分别为 \(S_1,S_1,S_2\)。这里 \(P_3\) 不再包含 \(S_1\) 合成因子，因为关系 \(E_{12}E_{23}=E_{13}=0\) 切断了这两次作用的复合。
\end{solution}

\begin{optionalexercise}\label{b5:ex:08-c6-blocks}
	设 \(\operatorname{char}k=3\)，\(G=C_3\times C_2=\langle r\rangle\times\langle s\rangle\)。 求 \(kG\) 的块、简单模、不可分解投射模及 Cartan 矩阵。再证明与 \(s\mapsto-1,r\mapsto1\) 的一维表示作张量会交换两个块中的模。
\end{optionalexercise}
\begin{solution}
	因为 \(2\) 可逆，中心幂等元 \(e_+=(1+s)/2\)、\(e_-=(1-s)/2\) 正交且和为一。每个角理想 \(kGe_\pm\) 都同构于局部代数 \(kC_3\cong k[t]/(t^3)\)，故它们恰为两个块。 简单模 \(S_+,S_-\) 均为一维，\(r\) 平凡作用，\(s\) 分别按 \(1,-1\) 作用。投射覆盖为 \(P_\pm=kGe_\pm\)，每个有三个相同的简单合成因子，所以 \(C=\operatorname{diag}(3,3)\)。 与指定一维表示作张量后，\(r\) 的矩阵不变而 \(s\) 的矩阵乘 \(-1\)，因此 \(e_+\)、\(e_-\) 的作用互换。块分解使同态和扩张分开，却不表示张量任意表示都保持原块。
\end{solution}

\begin{optionalexercise}\label{b5:ex:08-s3-augmentation}
	设 \(k\) 为特征 \(2\) 且含有本原三次单位根 \(\alpha\) 的域，\(G=S_3=\langle r,s\rangle\)，其中 \(r^3=s^2=1\)、\(srs=r^{-1}\)。令 \(I=\ker(\varepsilon:kG\to k)\)。分别将 \(I\) 限制到 \(\langle r\rangle\) 与 \(\langle s\rangle\)，求其不可分解直和类型及不动空间维数，判断 \(I\) 是否投射。
\end{optionalexercise}
\begin{solution}
	限制到 \(C_3=\langle r\rangle\) 后，\(kG\) 是两个正则 \(kC_3\)-模的直和。记 \(L_j\) 为 \(r\) 作用为 \(\alpha^j\) 的一维模。由 Maschke 定理，增广列在 \(C_3\) 上分裂，故
	\[
		\operatorname{Res}_{C_3}^G I\cong L_0\oplus L_1^{\oplus2}\oplus L_2^{\oplus2}, \qquad \dim_k I^{C_3}=1.
	\]
	限制到 \(C_2=\langle s\rangle\)，把六个群元素配成三对 \(g_i,sg_i\)，令 \(T=s-1\)。对向量 \(\sum_i(a_ig_i+b_isg_i)\)，有 \(T\bigl(\sum_i(a_ig_i+b_isg_i)\bigr)=\sum_i(a_i+b_i)(g_i+sg_i)\)。属于 \(I\) 的条件是 \(\sum_i(a_i+b_i)=0\)，故 \(T|_I\) 的秩为 \(2\)。又 \(T^2=0\)、\(\dim I=5\)，得到
	\[
		\operatorname{Res}_{C_2}^G I\cong M_2^{\oplus2}\oplus M_1, \qquad \dim_k I^{C_2}=3.
	\]
	\(I\) 不投射：投射 \(kG\)-模限制到 \(C_2\) 后仍投射，而 \(kC_2\) 的有限维投射模都自由，维数必须为偶数；这里维数为 \(5\)。同一个模限制到不同子群后，既可能完全可约，也可能保留非分裂的幂零链。
\end{solution}

\begin{optionalexercise}\label{b5:ex:08-regular-set-and-parts}
	证明 \(S_3\) 的 \(3\)-正则元素不构成子群。再设 \(g\) 的阶为 \(12\)，按 \(p=2\) 写出其交换的 \(p\) 部分与 \(p'\) 部分，并核验两部分的阶。
\end{optionalexercise}
\begin{solution}
	恒等元及三个换位的阶与 \(3\) 互素，但两个不同换位的乘积是三循环，所以该集合不封闭。 对阶为 \(12\) 的 \(g\)，取
	\[
		g_2=g^9,\qquad g_{2'}=g^4.
	\]
	前者的阶为 \(12/\gcd(12,9)=4\)，后者的阶为 \(3\)，且二者都在 \(\langle g\rangle\) 中，所以交换；乘积为 \(g^{13}=g\)。指数 \(9\) 在模 \(4\) 时为 \(1\)、模 \(3\) 时为 \(0\)，指数 \(4\) 则相反，这就是相应的中国剩余分解。
\end{solution}

\begin{optionalexercise}\label{b5:ex:08-trace-brauer-distinction}
	设 \(k\) 的特征为 \(p\)，\(G=C_p\)。比较正则模 \(M_p=kG\)、\(p\) 份平凡模的直和和零模的 \(k\) 值迹函数、Brauer 特征标及不动空间。哪些资料能够区分这三个对象？
\end{optionalexercise}
\begin{solution}
	\(G\) 的唯一 \(p\)-正则元素是恒等元。前两个模的 Brauer 特征标在该元素上都取整数 \(p\)，零模取 \(0\)，所以 Brauer 特征标区分它们与零模，却不能区分这两个非零模。 三者的 \(k\) 值迹函数都为零：在两个 \(p\) 维模上，每个群元素都幂幺，迹为域中的 \(p=0\)。不动空间则分别为一维、\(p\) 维及零维，因而能够区分三者。普通的模特征矩阵迹会进一步丢掉整数重数，Brauer 特征标保留重数，但仍不记录扩张是否分裂。
\end{solution}

\begin{optionalexercise}\label{b5:ex:08-cyclic-mixed-brauer}
	设 \(G=C_{p^r}\times C_\ell=\langle s\rangle\times\langle h\rangle\)，其中 \(p\nmid\ell\)，并固定本章的分裂 \(p\)-模系统。令 \(\alpha\in k\) 为本原 \(\ell\) 次单位根，\(\zeta=\widehat\alpha\)。求全部简单模、不可分解模及它们的 Brauer 特征标，再求分解矩阵和 Cartan 矩阵。
\end{optionalexercise}
\begin{solution}
	中心幂等元
	\[
		e_j=\frac1\ell\sum_{v=0}^{\ell-1}\alpha^{-jv}h^v \qquad(0\leq j<\ell)
	\]
	将群代数分成 \(\ell\) 个块，每个块同构于 \(kC_{p^r}\)。第 \(j\) 块的唯一简单模 \(S_j\) 由 \(s\mapsto1,h\mapsto\alpha^j\) 给出，全部不可分解模为 \(M_d\otimes S_j\)，其中 \(1\leq d\leq p^r\)。 \(p\)-正则元素恰为 \(h^v\)，故
	\[
		\varphi_{M_d\otimes S_j}(h^v)=d\,\iota(\zeta)^{jv},
	\]
	其中 \(\iota:\mathbb Q(\mu_{|G|})\to\mathbb C\) 是定义 Brauer 特征标时固定的嵌入。 在特征零中，所有简单表示均一维，由 \(s\) 的一个 \(p^r\) 次单位根及 \(h\mapsto\zeta^j\) 确定。前者约化后必须为 \(1\)，因为特征 \(p\) 中 \(X^{p^r}-1=(X-1)^{p^r}\)。所以每个这样的普通表示约化为 \(S_j\)，分解矩阵的相应行就是第 \(j\) 个标准行向量，且每一种行重复 \(p^r\) 次。 每个简单模的投射覆盖是该块的正则模，具有 \(p^r\) 个相同简单因子。因此 \(C=p^rI_\ell\)，也在本例直接满足 \(D^{\mathsf T}D=C\)。
\end{solution}

\begin{optionalexercise}\label{b5:ex:08-s3-characteristic-three}
	设 \(k\) 为特征 \(3\) 的域，\(S_3\) 按置换坐标作用于 \(U=\{(x,y,z)\in k^3:x+y+z=0\}\)。记 \(S_+,S_-\) 为平凡模与符号模，\(P_+\) 为 \(S_+\) 的投射覆盖。证明对角作用下有实际模同构
	\[
		U\otimes_kU\cong P_+\oplus S_-.
	\]
	要求区分由 Brauer 特征标得到的合成重数与证明上述直和分解所需的附加信息。
\end{optionalexercise}
\begin{solution}
	取 \(u=(1,-1,0)\)、\(v=(0,1,-1)\)，令 \(r=(123)\)、\(s=(12)\)。\eqref{b5:eq:s3-modular-matrix-input}给出 \(r,s\) 在这组基上的矩阵。交换两个张量因子的算子与群作用交换，且 \(2\) 在 \(k\) 中可逆，所以其 \(+1,-1\) 特征空间给出实际分解
	\[
		U\otimes U=\operatorname{Sym}^2U\oplus\Lambda^2U.
	\]
	这里将对称幂认作由平均投影得到的子空间。\(r,s\) 在一维外幂上的作用分别为其行列式 \(1,-1\)，故 \(\Lambda^2U\cong S_-\)。在对称部分的基 \(u\otimes u,\ u\otimes v+v\otimes u,\ v\otimes v\) 上，两矩阵为
	\[
		C=\begin{pmatrix}0&0&1\\0&2&1\\1&1&1\end{pmatrix},\qquad
		E=\begin{pmatrix}1&1&1\\0&2&1\\0&0&1\end{pmatrix}.
	\]
	计算得 \((C-I)^3=0\)、\((C-I)^2\ne0\)。因此这个三维空间限制到 \(\langle r\rangle=C_3\) 后只有一条长度三的幂零链，是秩一自由 \(kC_3\)-模。由\cref{b5:prop:relative-sylow-projectivity}，\(\operatorname{Sym}^2U\) 是投射 \(kS_3\)-模。

	\(U\) 的 Brauer 特征标在单位元、换位两类上为 \((2,0)\)，故张量积为 \((4,0)=2\varphi_++2\varphi_-\)。扣除外幂的 \(\varphi_-\) 后，对称部分有两个 \(S_+\) 因子和一个 \(S_-\) 因子。\cref{b5:exm:s3-three-block-local}给出两个不可分解投射覆盖 \(P_+,P_-\)，维数均为三，合成重数分别为 \((2,1),(1,2)\)。三维投射模的分解只能取其中一个，重数于是迫使对称部分同构于 \(P_+\)。这就证明所求同构；仅有 \((4,0)\) 不能独自证明投射部分的存在，也不能决定非分裂的连接方式。
\end{solution}

\begin{optionalexercise}\label{b5:ex:08-s3-three-cartan}
	设 \(p\) 为奇素数，\(k\) 为特征 \(p\) 的域，
	\[
		G=\langle r,s\mid r^p=s^2=1,\ srs=r^{-1}\rangle
	\]
	为阶 \(2p\) 的二面体群。令 \(t=r-1\)、\(e_\pm=(1\pm s)/2\)。求 \(kG\) 的根、简单模、不可分解投射覆盖及其根层，证明它只有一个块，并计算 Cartan 矩阵及其行列式。再与同一特征下的 \(C_p\times C_2\) 比较，说明群阶与 Sylow 子群同构类型相同并不足以决定块分解。
\end{optionalexercise}
\begin{solution}
	关系给出 \(t^p=0\) 及
	\[
		sts=(1+t)^{-1}-1=-t+t^2-\cdots+t^{p-1}.
	\]
	所以 \((t)\) 是双边理想，其 \(i\) 次幂由 \(t^j,t^js\)、\(i\le j<p\) 张成，\(p\) 次幂为零。商为 \(kC_2\cong k\times k\)，半单，故 \(J(kG)=(t)\)。简单模只有 \(S_+,S_-\)：\(r\) 平凡作用，\(s\) 分别作用为 \(1,-1\)。

	正则模分解为 \(P_+\oplus P_-\)，其中 \(P_\pm=kGe_\pm\) 的基为 \(e_\pm,te_\pm,\ldots,t^{p-1}e_\pm\)。它们的顶分别为 \(S_\pm\)，因而是这两个简单模的不可分解投射覆盖。在第 \(i\) 个根层，\(r\) 平凡作用，而
	\[
		s t^ie_\pm\equiv(\pm1)(-1)^it^ie_\pm
		\pmod{J^{i+1}P_\pm}.
	\]
	于是 \(P_+\) 的根层从顶到底交替为 \(S_+,S_-,S_+,\ldots,S_+\)，\(P_-\) 的交替次序相反。两种简单模同在一个不可分解投射模的合成列中，故属于同一个块；简单模已经穷尽，所以没有其他块。按简单模作行、投射覆盖作列，
	\[
		C_G=\begin{pmatrix}(p+1)/2&(p-1)/2\\(p-1)/2&(p+1)/2\end{pmatrix},
		\qquad \det C_G=p.
	\]
	对于 \(C_p\times C_2\)，两个 \(e_\pm\) 都在中心，每个块为 \(kC_p\)，故有两个块，Cartan 矩阵为 \(\operatorname{diag}(p,p)\)。两群都有阶 \(2p\)，Sylow \(p\) 子群均为循环群，但前者的反射 \(s\) 把 \(r\) 变成 \(r^{-1}\)，使根层的简单类型交替；后者的作用交换，根层类型保持不变。这项作用信息不能从群阶及 Sylow 子群类型中读出。
\end{solution}


\begin{optionalexercise}\label{b5:ex:08-naive-brauer-inner-product}
	对特征 \(3\) 的 \(S_3\) 的两个不可约 Brauer 特征标 \(\varphi_+,\varphi_-\)，令 \(\Psi_+=2\varphi_++\varphi_-\)、\(\Psi_-=\varphi_++2\varphi_-\)，并定义
	\[
		(f,g)_0=\frac1{|S_3|}\sum_{x\in(S_3)_{3'}}\overline{f(x)}g(x).
	\]
	计算 \((\Psi_i,\varphi_j)_0\) 与 \((\Psi_i,\Psi_j)_0\)。说明一组基不正交，仍可以有另一组基与它对偶，并比较后一矩阵与本例的 Cartan 矩阵。
\end{optionalexercise}
\begin{solution}
	在由恒等元、换位代表的两类上，\(\varphi_+=(1,1)\)、\(\varphi_-=(1,-1)\)，所以 \(\Psi_+=(3,1)\)、\(\Psi_-=(3,-1)\)。恒等元只有一个，换位有三个，因此
	\[
		(\Psi_+,\varphi_+)_0=(\Psi_-,\varphi_-)_0=\frac{3+3}{6}=1, \qquad (\Psi_+,\varphi_-)_0=(\Psi_-,\varphi_+)_0=0.
	\]
	所以两组基互为对偶。另一方面，\(\Psi_\pm\) 的平方长度均为 \((9+3)/6=2\)，互相的内积为 \((9-3)/6=1\)，故它们自身的 Gram 矩阵是 \(\left(\begin{smallmatrix}2&1\\1&2\end{smallmatrix}\right)\)，恰与\cref{b5:exm:s3-three-block-local}算出的 Cartan 矩阵相同。这是本例的计算，并不是从普通不可约特征标的正交公式推出的。
\end{solution}

\begin{optionalexercise}\label{b5:ex:08-s3-two-tensor}
	在特征 \(2\) 的 \(S_3\) 模中，记 \(S_1\) 为平凡模，\(S_2\) 为二维简单投射模，\(P_1\) 为 \(S_1\) 的投射覆盖。先由 Brauer 特征标求 \(S_2\otimes S_2\) 的简单重数，再证明实际模同构
	\[
		S_2\otimes S_2\cong P_1\oplus S_2.
	\]
\end{optionalexercise}
\begin{solution}
	两类 \(2\)-正则元素由 \(1,r\) 代表，\(\varphi_{S_1}=(1,1)\)、\(\varphi_{S_2}=(2,-1)\)。张量特征标为 \((4,1)\)，解 \((4,1)=a(1,1)+b(2,-1)\) 得到 \(a=2,b=1\)。这一步只确定两个平凡因子与一个二维简单因子。 由于 \(S_2\) 投射，\cref{b5:ex:08-projective-tensor}保证张量积也投射。全部不可分解投射类型为 \(P_1,S_2\)，它们的简单重数分别是 \((2,0)\)、\((0,1)\)。因此投射直和分解只能各取一份，得到所求同构。这里从合成重数走到实际直和，额外使用了投射性，不能仅凭特征标相同就作这一步推断。
\end{solution}

\begin{optionalexercise}\label{b5:ex:08-s3-two-lattices}
	设 \((K,\mathcal O,k)\) 为本章特征 \(3\) 的分裂模系统，\(\pi\) 为统一化元，并让 \(S_3\) 置换三个坐标。在 \(V=K^3/K(1,1,1)\) 中取满格
	\[
		\Lambda=\mathcal O^3/\mathcal O(1,1,1),\qquad L=\text{\(\bigl\{(x,y,z)\in\mathcal O^3:x+y+z=0\bigr\}\) 在 \(V\) 中的像}.
	\]
	证明 \(\Lambda/L\cong\mathcal O/3\mathcal O\)，并证明两个约化不同构，却有相同的简单合成因子及 Brauer 特征标。
\end{optionalexercise}
\begin{solution}
	对角向量有一个单位坐标，故商 \(\Lambda\) 为秩二自由模；和为零的子模也自由秩二，在 \(V\) 中的映射单射，因为特征零中常向量的坐标和为 \(3c\)，只有 \(c=0\) 时为零。它们都是被群保持的满格。 求和映射识别 \(\mathcal O^3/\ker\Sigma\cong\mathcal O\)，而对角子模在其下的像为 \(3\mathcal O\)，所以 \(\Lambda/L\cong\mathcal O/3\mathcal O\)。 约化后，\(L/\pi L\) 是 \(k^3\) 的二维增广子模，含一维平凡子模，商为符号模。另一方面，\(\Lambda/\pi\Lambda\cong k^3/k(1,1,1)\)。它的增广子模像是一维符号模，商为平凡模，因此简单重数相同。 但第二个约化没有非零群不动元：若向量 \(v\in k^3\) 的类不动，则每个换位使 \(v\) 的变化都在对角直线上。这个变化有一个坐标为零，故只能是零；所有换位因此固定 \(v\)，使 \(v\) 本身为常向量，其商类为零。第一个约化的不动空间却由 \((1,1,1)\) 生成，维数为一，故两个约化不同构。它们的 Brauer 特征标都为 \(\varphi_++\varphi_-\)。
\end{solution}

\begin{optionalexercise}\label{b5:ex:08-decomposition-column-rank}
	设 \(G\) 为有限群，\(p\) 为素数，\((K,\mathcal O,k)\) 为 \(G\) 的分裂 \(p\)-模系统，\(K\) 含有全部 \(|G|\) 次单位根，并固定这些根到复数的对应。记 \(d:R_K(G)\to G_0(kG)\) 为分解映射，\(D\) 为分解矩阵。证明 \(D\) 的列在 \(\mathbb C\) 上线性无关，并据此计算 \(\mathbb Q\otimes\ker d\) 的维数：用普通共轭类数与 \(p\)-正则共轭类数表示。说明这一秩结论本身为何不能证明分解映射在整数上满射。
\end{optionalexercise}
\begin{solution}
	普通不可约特征标限制到 \(G_{p'}\) 后生成全部复类函数，而不可约 Brauer 特征标是一组基。恒等式 \(\chi_\alpha|_{G_{p'}}=\sum_i d_{\alpha i}\varphi_i\) 因此说明 \(D\) 的行向量生成整个 \(\mathbb C^s\)，其中 \(s\) 为简单模数，即 \(p\)-正则类数。所以 \(D\) 的秩为 \(s\)，各列线性无关。 普通表示群的秩为全部共轭类数 \(c\)，分解映射的矩阵是 \(D^{\mathsf T}\)。该整数矩阵在 \(\mathbb Q\) 与 \(\mathbb C\) 上秩相同，故有理化后的核维数为 \(c-s\)。但满秩的整数矩阵未必在整数格上满射，例如一阶矩阵 \((2)\) 的有理秩为一，整数像却为 \(2\mathbb Z\)。因此若要证明整数上的更强结论，还需额外论证。
\end{solution}

\begin{optionalexercise}\label{b5:ex:08-brauer-domain}
	设 \(k\) 为特征 \(2\) 的域，\(G=S_3=\langle r,s\mid r^3=s^2=1,\ srs=r^{-1}\rangle\)，\(P=\langle s\rangle\)。求共轭固定子代数 \((kG)^P\) 的一组基及 \(\operatorname{Br}_P^G\) 的核，判断这个核是否幂零，再将它看成有限维含幺代数，计算它自身的 Jacobson 根。
\end{optionalexercise}
\begin{solution}
	\(P\) 的共轭轨道是 \(\{1\},\{s\},\{r,r^2\},\{rs,r^2s\}\)，所以固定子代数以 \(1,s,a=r+r^2,b=rs+r^2s\) 为基。\(a\) 与 \(s\) 交换，且 \(b=as\)。特征 \(2\) 中
	\[
		a^2=r^2+r=a,\qquad ab=a^2s=b,\qquad b^2=a^2s^2=a.
	\]
	Brauer 映射将 \(1,s\) 保留、将 \(a,b\) 送到零，所以核为 \(ka\oplus kb\)，是理想；它不幂零，因为 \(a^n=a\ne0\) 对每个 \(n\ge1\) 成立。这个核以 \(a\) 为单位，\(b^2=a\)，所以同构于 \(kC_2\)。其自身的根为 \(k(a+b)\)，平方为零。Brauer 映射的核与定义域的 Jacobson 根因而不能直接混为一谈。
\end{solution}

\begin{optionalexercise}\label{b5:ex:08-defect-change-prime}
	在特征 \(3\) 中，求\cref{b5:ex:08-c6-blocks}的 \(C_6\) 两个块的缺陷群，以及\cref{b5:exm:s3-three-block-local}的 \(S_3\) 唯一块的缺陷群。将后者与正文中特征 \(2\) 的 \(S_3\) 比较。
\end{optionalexercise}
\begin{solution}
	\(C_6\) 交换，所以任一子群的中心化子都是全群，Brauer 截断在群代数上就是恒等。因此两个非零块幂等元对 Sylow \(3\) 子群 \(C_3\) 的 Brauer 像都非零；没有更大的 \(3\) 子群，两个块的缺陷群都是 \(C_3\)。 特征 \(3\) 的 \(S_3\) 只有一个块，幂等元为 \(1\)。对 \(P=\langle r\rangle\) 有 \(\operatorname{Br}_P^G(1)=1\neq0\)，且 \(P\) 是 Sylow \(3\) 子群，所以它是缺陷群。 特征 \(2\) 时，\(S_3\) 则有一个缺陷群为 \(C_2\) 的二维局部块和一个缺陷群平凡的四维矩阵块。改变模特征会同时改变块的数目、块的结构以及相应的缺陷群。
\end{solution}

\begin{optionalexercise}\label{b5:ex:08-induction-reduction-example}
	设 \(k\) 为特征 \(3\) 的域，\(G=S_3=\langle r,s\mid r^3=s^2=1,\ srs=r^{-1}\rangle\)，\(H=\langle r\rangle\cong C_3\)、\(t=r-1\)。令 \(M_d=kH/(t^d)\)，其中 \(d=1,2,3\)。求三个诱导模 \(\operatorname{Ind}_H^G M_d\) 的不可分解直和类型及简单合成因子，说明诱导一个不可分解模可能得到可分解模。
\end{optionalexercise}
\begin{solution}
	记 \(S_+,S_-\) 为平凡与符号简单模，\(P_\pm=kGe_\pm\)，其中 \(e_\pm=(1\pm s)/2\)。由\cref{b5:exm:s3-three-block-local}，\(J=J(kG)=(t)\)，且 \(P_+\) 的根层从顶到底为 \(S_+,S_-,S_+\)，\(P_-\) 的为 \(S_-,S_+,S_-\)。 \(M_1\) 是 \(H\) 的平凡模，诱导后为两陪集上的置换模；它通过 \(G/H\cong C_2\) 作用，故
	\[
		\operatorname{Ind}_H^G M_1\cong S_+\oplus S_-.
	\]
	\(M_3=kH\) 自由，乘法映射给出 \(\operatorname{Ind}_H^G M_3\cong kG=P_+\oplus P_-\)。 对正合列 \(0\to t kH\to kH\to M_1\to0\)，其核 \(t kH\cong M_2\)。诱导保持正合，得到
	\[
		0\to\operatorname{Ind}_H^G M_2\to kG\to k[G/H]\to0.
	\]
	右端是商代数 \(kG/J\) 的正则模，故中间满射的核就是 \(J\)。于是 \(\operatorname{Ind}_H^G M_2\cong JP_+\oplus JP_-\)。两项均为二维不可分解模：在基 \(te_\pm,t^2e_\pm\) 上，\(t\) 非零而平方为零；若分成两个一维模，正规 \(C_3\) 将在两项上都平凡作用，与此矛盾。 三种诱导模的简单重数依次为 \((1,1),(2,2),(3,3)\)，其中第二个来自两个根模各有一个 \(S_+\) 和一个 \(S_-\)。源端的 \(M_d\) 都不可分解，但这里三个诱导模都分解为两个非零项。
\end{solution}

\begin{optionalexercise}\label{b5:ex:08-symmetric-center-dual}
	设 \(A\) 为有限维对称 \(k\) 代数，\(\lambda:A\to k\) 为定义中的线性泛函，并记 \([A,A]\) 为全部交换子 \(ab-ba\) 张成的 \(k\) 子空间，不把它改成由交换子生成的双边理想。证明配对诱导同构
	\[
		Z(A)\cong\operatorname{Hom}_k(A/[A,A],k),
		\qquad z\longmapsto\bigl(x+[A,A]\mapsto\lambda(zx)\bigr).
	\]
	对群代数 \(A=kG\)，取 \(\lambda\) 为单位群元素的系数，证明 \([A,A]\) 正是同一共轭类中群元素之差张成的空间，并说明上述配对如何将共轭类和与商空间的共轭类基配对。
\end{optionalexercise}
\begin{solution}
	令 \([A,A]^\perp=\{z\in A:\lambda(zc)=0\text{ 对所有 }c\in[A,A]\}\)。由 \(\lambda(uv)=\lambda(vu)\)，对任意 \(a,b,z\in A\) 有
	\[
		\lambda(z(ab-ba))=\lambda((za-az)b).
	\]
	若 \(z\) 在中心，右边为零。反过来，若左边对所有 \(a,b\) 为零，配对非退化使 \(za-az=0\) 对每个 \(a\) 成立。因此 \([A,A]^\perp=Z(A)\)。非退化配对把 \(A\) 同构到其线性对偶，\([A,A]\) 的零化子在该同构下正好对应商空间的对偶，得到所述同构。这一步只对向量空间取商，没有使用 \([A,A]\) 为理想。

	群代数中的交换子由 \(gh-hg\) 张成，而 \(gh\) 与 \(hg\) 共轭。反过来，对 \(x,h\in G\)，有
	\[
		xhx^{-1}-h=[x,hx^{-1}],
	\]
	所以同一共轭类内的差都属于交换子空间。商空间因此每个共轭类保留一维：不同类的支集不相交，每一类内部模去系数和为零的子空间。若 \(C,D\) 为共轭类，\(c_C=\sum_{g\in C}g\)，取 \(h\in D\)，则
	\[
		\lambda(c_Ch)=
		\begin{cases}1,&D=C^{-1},\\0,&D\ne C^{-1}.
		\end{cases}
	\]
	故共轭类和与逆共轭类的商基互为对偶；此计算在模特征下也成立，没有除以共轭类大小。
\end{solution}

\begin{optionalexercise}\label{b5:ex:08-cyclic-vertices-sources}
	设 \(p\) 为素数，\(k\) 为特征 \(p\) 的域，\(G=C_{p^r}=\langle g\rangle\)、\(r\ge1\)，令 \(t=g-1\)、\(M_d=k[t]/(t^d)\)，其中 \(1\le d\le p^r\)。对 \(0\le s\le r\)，记 \(H_s=\langle g^{p^{r-s}}\rangle\)。证明 \(M_d\) 相对于 \(H_s\) 投射，当且仅当 \(p^{r-s}\mid d\)。据此求出每个 \(M_d\) 的顶点和源，并计算 \(G=C_8\)、\(d=6\) 的例子。
\end{optionalexercise}
\begin{solution}
	写 \(q=p^{r-s}\)，令 \(u=g^q-1\)。特征 \(p\) 中有 \(u=t^q\)，所以 \(kH_s\cong k[u]/(u^{p^s})\)。这个子群的不可分解模恰为 \(U_e=k[u]/(u^e)\)、\(1\le e\le p^s\)，包括 \(s=0\) 时的唯一模 \(U_1=k\)。诱导给出
	\[
		kG\otimes_{kH_s}U_e\cong k[t]/(t^{qe})=M_{qe}.
	\]
	每个 \(U_e\) 的诱导仍不可分解。若 \(M_d\) 相对投射，则它是某个有限维 \(H_s\)-模诱导所得的直和因子。把该模分解为各 \(U_e\)，诱导后是各 \(M_{qe}\) 的直和；\cref{b5:thm:cyclic-p-modules}的分解唯一性迫使 \(d=qe\) 对某个 \(e\) 成立。反过来，若 \(q\mid d\)，上述同构直接使 \(M_d\) 从 \(H_s\) 诱导得到。

	令 \(d=p^a e\)，其中 \(p\nmid e\)、\(0\le a\le r\)。最小的可用子群为 \(H_{r-a}\)，阶为 \(p^{r-a}\)，故它是顶点。对应的源是 \(k[u]/(u^e)\)，其中 \(u=t^{p^a}\)：诱导它恰好得到 \(M_d\)。也可直接验证限制模为 \(p^a\) 个这样的源的直和，因为
	\[
		M_d=\bigoplus_{i=0}^{p^a-1}
		\operatorname{span}_k\{t^i,t^{i+p^a},\ldots,t^{i+(e-1)p^a}\},
	\]
	各项对乘 \(u\) 稳定，都是长度 \(e\) 的链。对 \(C_8\)、\(d=6=2\cdot3\)，顶点是 \(\langle g^2\rangle\cong C_4\)，源为 \(k[u]/(u^3)\)，限制模是该源的两份直和。特别地，源的维数不必为一。
\end{solution}

\begin{optionalexercise}\label{b5:ex:08-tensor-vertex-intersection}
	设 \(G\) 为有限群，\(k\) 为特征 \(p>0\) 的域，\(M,N\) 为非零有限维不可分解左 \(kG\)-模，分别有顶点 \(P,Q\)。张量积取对角群作用。证明 \(M\otimes_kN\) 的每个非零不可分解直和因子都有一个顶点包含于某个 \(P\cap gQg^{-1}\)。据此证明：若 \(P\cap gQg^{-1}=1\) 对所有 \(g\in G\) 成立，则 \(M\otimes_kN\) 投射。用 \(G=C_p\times C_p\) 两个不同直因子的诱导平凡模给出两个非投射模张量后投射的例子。
\end{optionalexercise}
\begin{solution}
	取 \(M\) 的一个 \(P\)-源 \(S\)，使 \(M\) 为 \(\operatorname{Ind}_P^G S\) 的直和因子。张量保持分裂映射，而\cref{b5:thm:tensor-identity}给出
	\[
		(\operatorname{Ind}_P^G S)\otimes_kN
		\cong\operatorname{Ind}_P^G(S\otimes_k\operatorname{Res}_P^GN).
	\]
	因此 \(M\otimes_kN\) 相对于 \(P\) 投射。交换两个张量因子后，同理知道它相对于 \(Q\) 投射。每个不可分解直和因子 \(V\) 也具有这两种相对投射性，由\cref{b5:lem:relative-intersection}可知 \(V\) 相对于某个 \(P\cap gQg^{-1}\) 投射；\cref{b5:thm:vertices-sources}使该交集包含 \(V\) 的一个顶点。若这些交集均平凡，所有不可分解项都有平凡顶点，故各项投射，其有限直和也投射。

	具体取 \(G=P\times Q\)、\(P\cong Q\cong C_p\)，令 \(M=k[G/P]\)、\(N=k[G/Q]\)。\(M\) 的作用通过 \(G/P\cong C_p\) 给出正则模，因而不可分解；\(N\) 同样不可分解。两模维数均为 \(p\)，而群代数 \(kG\) 是局部代数，\cref{b5:thm:split-projective-covers}说明其有限维投射模都是自由模，维数为 \(p^2\) 的倍数，所以两模均不投射。张量基由 \((G/P)\times(G/Q)\) 给出，映射
	\[
		G\longrightarrow(G/P)\times(G/Q),\qquad x\longmapsto(xP,xQ)
	\]
	是 \(G\)-集合双射：核条件为 \(P\cap Q=1\)，两边大小都为 \(p^2\)。因此 \(M\otimes_kN\cong kG\)，是秩一自由模。
\end{solution}

\begin{optionalexercise}\label{b5:ex:08-syzygy-tensor}
	设 \(G\) 为有限群，\(k\) 为域，\(M\) 为有限维左 \(kG\)-模。以对角群作用证明稳定范畴中有同构
	\[
		M\otimes_k\Omega k\cong\Omega M.
	\]
	这里的 \(k\) 为平凡模；\(\Omega\) 允许从任意有限维投射满射取核。取 \(G=C_p\)，说明这个稳定同构不总是实际模同构。
\end{optionalexercise}
\begin{solution}
	选 \(0\to\Omega k\to P\to k\to0\)，在域上与 \(M\) 张量保持正合，得到
	\[
		0\longrightarrow M\otimes_k\Omega k
		\longrightarrow M\otimes_kP\longrightarrow M\longrightarrow0.
	\]
	中间项投射：将 \(P\) 补成有限自由模，张量后的中间项便是若干 \(M\otimes_k kG\) 的直和因子；\cref{b5:thm:tensor-identity}在平凡子群处的同构说明每个 \(M\otimes_k kG\) 都是有限自由模。故这个核是由一个投射满射定义的合冲模，与任意另选的 \(\Omega M\) 通过\eqref{b5:eq:syzygy-schanuel}只相差投射直和项，得到稳定同构。

	取特征 \(p\) 的 \(k\)、\(G=C_p\)、\(M=kG\)，并用恒等映射作为 \(M\) 的投射覆盖，则 \(\Omega M=0\)。而由增广覆盖得到的 \(\Omega k\) 维数为 \(p-1\)，故 \(M\otimes_k\Omega k\) 的维数为 \(p(p-1)>0\)，实际模不为零。它是有限自由模，因而在稳定范畴中同构于零，恰与公式一致。
\end{solution}

\begin{optionalexercise}\label{b5:ex:08-finite-projective-resolution}
	设 \(G\) 为有限群，\(k\) 为域，\(M\) 为有限维左 \(kG\)-模。若存在有限正合列
	\[
		0\longrightarrow P_n\longrightarrow P_{n-1}\longrightarrow\cdots
		\longrightarrow P_0\longrightarrow M\longrightarrow0,
	\]
	其中各 \(P_i\) 有限维且投射，证明 \(M\) 投射。由此说明，当 \(\operatorname{char}k=p\) 且 \(p\,\bigm\vert\,|G|\) 时，平凡模 \(k\) 不能有这种有限投射消解。
\end{optionalexercise}
\begin{solution}
	若 \(n=0\)，结论直接成立。若 \(n\ge1\)，\cref{b5:prop:group-algebra-symmetric}说明 \(P_n\) 同时内射，所以单射 \(P_n\to P_{n-1}\) 分裂。它的余核 \(K\) 是 \(P_{n-1}\) 的直和因子，因而投射；原正合列去掉这两个位置后变成
	\[
		0\longrightarrow K\longrightarrow P_{n-2}\longrightarrow\cdots
		\longrightarrow P_0\longrightarrow M\longrightarrow0,
	\]
	其中 \(n=1\) 时直接得到 \(M\cong K\)。逐次缩短消解，就归结到 \(n=0\) 的情形。

	若 \(p\,\bigm\vert\,|G|\)，平凡模 \(k\) 不是投射模，因为增广满射 \(kG\to k\) 的任何 \(G\)-不动提升都为 \(c\sum_{g\in G}g\)，其增广为 \(c|G|=0\)，不能提升一。因此它不可能具有上述有限投射消解。有限维模的向量空间维数有限，并不保证消解在有限步后终止。
\end{solution}

\begin{optionalexercise}\label{b5:ex:08-product-brauer-correspondence}
	设 \(k\) 为含有本原三次单位根 \(\zeta\) 的特征二域，\(G=S_3\times\langle z\rangle\)，其中 \(z^3=1\)。沿用 \(S_3\) 中的 \(r=(123)\)、\(s=(12)\)，令 \(e=1+r+r^2\)、\(f=r+r^2\)，并令
	\[
		c_j=\frac13\sum_{a=0}^2\zeta^{-ja}z^a\quad(0\le j\le2).
	\]
	求 \(kG\) 的全部块及其缺陷群。对 \(P=\langle s\rangle\times1\)，计算全部缺陷为 \(P\) 的 Brauer 对应，并求一维模 \(k\otimes k_{\zeta^j}\) 的 Green 对应；这里 \(z\) 在 \(k_{\zeta^j}\) 上作用为 \(\zeta^j\)。
\end{optionalexercise}
\begin{solution}
	\(c_jc_\ell=\delta_{j\ell}c_j\)、\(\sum_jc_j=1\)，且 \(zc_j=\zeta^jc_j\)，因此 \(kC_3=\prod_{j=0}^2kc_j\)。结合正文已经证明的 \(kS_3\cong kC_2\times M_2(k)\)，得到
	\[
		kG\cong\prod_{j=0}^2\bigl(kC_2\times M_2(k)\bigr).
	\]
	所以全部块幂等元为 \(ec_j,fc_j\)，各三枚。\(P\) 的正规化子与中心化子均为 \(H=\langle s\rangle\times\langle z\rangle\)。截断只作用于 \(S_3\) 坐标，故
	\[
		\operatorname{Br}_P^G(ec_j)=c_j,\qquad
		\operatorname{Br}_P^G(fc_j)=0.
	\]
	\(G\) 的非平凡 \(2\)-子群都是 \(P\) 的共轭，因此 \(ec_j\) 的缺陷群为 \(P\)，\(fc_j\) 的缺陷群为一。\(kHc_j\cong kC_2\) 是缺陷为 \(P\) 的块，Brauer 对应就是 \(ec_j\leftrightarrow c_j\)。

	一维模 \(V_j=k\otimes k_{\zeta^j}\) 限制到 \(H\) 后仍不可分解。它不是投射模，因为限制到 \(P\) 后是一维平凡模，而特征二下的投射 \(kP\)-模维数是二的倍数。它相对 \(P\) 投射，因为 \(P\) 是 Sylow 子群；故其顶点为 \(P\)，Green 对应为 \(f(V_j)=\operatorname{Res}_H^G V_j\)。诱导回去则有
	\[
		\operatorname{Ind}_H^G f(V_j)
		\cong V_j\oplus(S_2\otimes k_{\zeta^j}),
	\]
	第二项是矩阵块 \(fc_j\) 中的投射简单模，并不属于与 \(V_j\) 相同的块。这也说明 Green 对应的额外诱导项可以落在别的块中。
\end{solution}

\begin{optionalexercise}\label{b5:ex:08-a4-green-correspondence}
	设 \(k\) 为特征三的域，\(G=A_4=V\rtimes\langle r\rangle\)，其中 \(V=\{1,a,b,c\}\cong C_2\times C_2\)，\(r\) 以三轮换置换 \(a,b,c\)。令 \(P=\langle r\rangle\)。求 \(kG\) 的两个块及其缺陷群，并证明主块同构于 \(kC_3\)。对 \(1\le d\le3\)，将 \(J_d=k[t]/(t^d)\) 视为 \(kP\)-模，其中 \(r\) 作用为 \(1+t\)；令 \(M_d\) 为沿 \(G\to G/V\cong P\) 膨胀所得的模。求 \(M_1,M_2\) 的 Green 对应，并将 \(\operatorname{Ind}_P^GJ_d\) 完整分解。
\end{optionalexercise}
\begin{solution}
	\(V\) 的四个一维特征标都在 \(k\) 中取值。令
	\[
		e_0=\frac14(1+a+b+c),\qquad
		\varepsilon_i=\frac14\sum_{v\in V}\chi_i(v)^{-1}v\quad(i\in\mathbb Z/3\mathbb Z),
	\]
	其中 \(\chi_0,\chi_1,\chi_2\) 为三个非平凡特征标，并编号使 \(r\varepsilon_ir^{-1}=\varepsilon_{i+1}\)，下标按模三理解。它们是 \(kV\) 中的正交本原幂等元。\(e_0\) 与 \(f=\varepsilon_0+\varepsilon_1+\varepsilon_2=1-e_0\) 被 \(G\) 共轭固定，故为中心幂等元。\(kGe_0\) 以 \(e_0,re_0,r^2e_0\) 为基，满足 \(r^3e_0=e_0\)，从而同构于 \(kC_3\)。

	在另一项中，九个元素
	\[
		E_{ij}=\varepsilon_i r^{i-j}\qquad(i,j\in\mathbb Z/3\mathbb Z)
	\]
	满足 \(E_{ij}E_{\ell m}=\delta_{j\ell}E_{im}\)。这是因为 \(r^{i-j}\varepsilon_\ell r^{j-i}=\varepsilon_{\ell+i-j}\)，与 \(\varepsilon_i\) 相乘只有 \(\ell=j\) 时不为零。它们构成九维空间 \(kGf\) 的一组基，所以 \(kGf\cong M_3(k)\)。因此 \(e_0,f\) 恰为两个块幂等元，前者包含平凡模，是主块；记后者的三维简单模为 \(S\)。

	在 \(A_4\) 中 \(N_G(P)=C_G(P)=P\)，而 \(V\cap P=1\)。特征三下 \(4=1\)，故
	\[
		\operatorname{Br}_P^G(e_0)=\frac14=1,\qquad
		\operatorname{Br}_P^G(f)=1-\frac14=0.
	\]
	所以主块的缺陷群为 \(P\)，矩阵块的缺陷群为一。\(M_d\) 通过 \(kGe_0\cong kP\) 对应于 \(J_d\)，因而不可分解；\(d=1,2\) 时不是投射模，且唯一可能的非平凡顶点为 \(P\)。限制到 \(P\) 就是 \(J_d\)，故 \(f(M_d)=J_d\)。

	诱导模的主块部分为
	\[
		e_0(kG\otimes_{kP}J_d)
		\cong kGe_0\otimes_{kP}J_d\cong M_d.
	\]
	另一块为矩阵块，其中每个模都是 \(S\) 的直和。诱导模维数为 \([G:P]d=4d\)，减去主块部分的 \(d\) 维，余下 \(3d\) 维。因此对三个 \(d\) 都有
	\[
		\operatorname{Ind}_P^GJ_d\cong M_d\oplus S^{\oplus d}.
	\]
	当 \(d=3\) 时，\(J_3\) 和 \(M_3\) 均投射，顶点是一，故这时的分解不应称为顶点 \(P\) 的 Green 对应。
\end{solution}

\begin{optionalexercise}\label{b5:ex:08-generalized-decomposition-s3}
	在满足\cref{b5:thm:brauer-second-main}假设的分裂 \(p\)-模系统中，对 \(S_3\) 的普通不可约特征标按平凡、符号、二维标准特征标的次序排列。分别取 \(p=2,u=(12)\) 与 \(p=3,u=(123)\)，求相应的全部广义分解数，并用 Brauer 第二主定理解释其中的零。说明这两个例子如何区别广义分解数与普通分解数。
\end{optionalexercise}
\begin{solution}
	两个中心化子分别为 \(C_2\)、\(C_3\)，都是 \(p\)-群，故各只有一个不可约 Brauer 特征标，且其定义域中的 \(p\)-正则元素只有单位元。因此广义分解数就是三个普通特征标在 \(u\) 处的值，分别为
	\[
		D^{(12)}=\begin{pmatrix}1\\-1\\0\end{pmatrix},\qquad
		D^{(123)}=\begin{pmatrix}1\\1\\-1\end{pmatrix}.
	\]
	在特征二中，平凡与符号普通特征标属于主块，二维标准特征标属于缺陷零矩阵块。后者的 Brauer 像为零，第二主定理因此要求它的整个局部列为零，恰与标准特征标在换位处的值一致。在特征三中，全群只有一个块，且它在 \(C_3\) 的 Brauer 像为一，故定理不强迫任何条目为零。这里两列都出现负数；普通分解数记录合成因子重数，不能为负，而广义分解数还含 \(u\) 的单位根特征值。第一列的 \(-1\) 就来自符号表示在换位上的作用。
\end{solution}

\begin{optionalexercise}\label{b5:ex:08-green-normalizer-condition}
	设 \(k\) 为特征三的域，\(G=S_3\)、\(Q=H=A_3\)。证明 \(H\) 不满足 Green 对应中 \(N_G(Q)\le H\) 的假设，并计算 \(\operatorname{Ind}_H^Gk\)。证明其两个不可分解直和项都有顶点 \(Q\)，从而说明删除正规化子条件后，“唯一顶点为 \(Q\) 的诱导直和项”会失败。
\end{optionalexercise}
\begin{solution}
	\(Q\) 在 \(G\) 中正规，所以 \(N_G(Q)=G\ne H\)。两个陪集上的置换表示是商群 \(G/H\cong C_2\) 的正则表示。因为二在 \(k\) 中可逆，它分裂为
	\[
		\operatorname{Ind}_H^Gk\cong S_+\oplus S_-.
	\]
	两个模均一维，因而不可分解；\(Q\) 在它们上面都作用为恒等。对任一子群 \(L\le G\)，在这两个模的自同态环 \(k\) 上共轭作用仍平凡，Higman 判据都是 \(3\nmid[G:L]\)。因此两者的顶点都是 Sylow 三子群 \(Q\)。这里出现两个不同构、顶点相同的直和项，正是正规化子条件缺失造成的失败。
\end{solution}

\begin{optionalexercise}\label{b5:ex:08-finite-field-absolute-test}
	设 \(k=\mathbb F_2\)，群 \(C_3=\langle r\rangle\) 在 \(W=k^2\) 上按矩阵 \(A=\left(\begin{smallmatrix}0&1\\1&1\end{smallmatrix}\right)\) 作用。证明 \(W\) 不可约但不是绝对不可约，求 \(\operatorname{End}_{kC_3}(W)\)，并写出扩张到 \(\mathbb F_4\) 后的简单分解。说明计算程序给出“不可约”时为何仍须保留工作域。
\end{optionalexercise}
\begin{solution}
	直接计算 \(A^2+A+I=0\)，故 \(A^3=I\)。多项式 \(X^2+X+1\) 在 \(\mathbb F_2\) 上没有根，所以 \(A\) 没有该域中的特征向量直线；二维模的非零真子模只能是一维，故 \(W\) 不可约。若矩阵 \(X\) 与 \(A\) 交换，以 \(e_1,Ae_1=e_2\) 为基可知 \(X\) 完全由 \(Xe_1=a e_1+bAe_1\) 决定，并且 \(X=aI+bA\)。因此
	\[
		\operatorname{End}_{kC_3}(W)=k[A]
		\cong\mathbb F_2[X]/(X^2+X+1)=\mathbb F_4.
	\]
	在 \(\mathbb F_4\) 中取 \(\alpha^2+\alpha+1=0\)。两个根 \(\alpha,\alpha^2\) 不同，列向量 \((1,\alpha)^{\mathsf T}\)、\((1,\alpha^2)^{\mathsf T}\) 分别为相应特征向量，且线性无关。故扩域后为两个一维表示的直和，\(r\) 分别作用为 \(\alpha,\alpha^2\)。原模因而不是绝对不可约。程序在 \(\mathbb F_2\) 上的不可约性结论与扩域后的可分解性并不冲突；把工作域删去，就会混淆这两个不同问题。
\end{solution}

\begin{optionalexercise}\label{b5:ex:08-spin-up-cyclic-vectors}
	设 \(k=\mathbb F_2\)，\(S_3\) 在 \(W=k^2\) 上的两个生成矩阵为
	\[
		A=\begin{pmatrix}0&1\\1&1\end{pmatrix},\qquad
		B=\begin{pmatrix}0&1\\1&0\end{pmatrix}.
	\]
	在 \(M=W\oplus W\) 上取对角作用。对 \(v=(x,y)\in M\)，证明由 \(v\) 生成的子模维数为 \(2\operatorname{rank}(x\ y)\)。计算十五个非零向量中，分别有多少个生成二维子模、多少个生成整个 \(M\)。解释后者为何不能作为不可约性的证据。
\end{optionalexercise}
\begin{solution}
	正文已通过 \(I,A,B,AB\) 的线性无关性证明它们生成的代数为 \(M_2(k)\)。故由 \((x,y)\) 生成的子模是
	\[
		\{\,(Cx,Cy):C\in M_2(k)\,\}.
	\]
	若 \(x=y=0\)，维数为零。若两列张成一维，可写 \(x=a z,y=b z\)，其中 \(z\ne0\)、\((a,b)\ne(0,0)\)；矩阵 \(C\) 的像 \(Cz\) 遍历 \(k^2\)，所得子模为 \(\{(aw,bw):w\in k^2\}\)，维数二。若 \((x\ y)\) 可逆，则给定任意 \((u,w)\)，取 \(C=(u\ w)(x\ y)^{-1}\)，便有 \((Cx,Cy)=(u,w)\)，所以生成整个四维模。这证明维数公式。

	秩一矩阵唯一写成 \(z(a,b)\)，其中 \(z\) 是三个非零列之一，\((a,b)\) 是三个非零行之一；由于 \(k^\times=\{1\}\)，不存在额外的标量重复，所以共有九个。可逆矩阵有 \((2^2-1)(2^2-2)=6\) 个。因此九个非零向量生成二维子模，六个生成整个 \(M\)。但 \(W\oplus0\) 明显是非零真子模，\(M\) 仍可约。找到循环生成元只说明整个模可以从这个向量生成，未排除从其他向量生成真子模。
\end{solution}

\begin{optionalexercise}\label{b5:ex:08-character-data-permutation}
	采用\cref{b5:prop:decomposition-data-recovery}的记号，令 \(P,Q,R\) 分别为普通不可约特征标、简单模和 \(p\)-正则类次序的置换矩阵。若新数据为 \(F'=PFR\)、\(\Phi'=Q\Phi R\)，求新分解矩阵。再对特征三的 \(S_3\)，只交换 Brauer 表中的两个共轭类列而保留普通限制表不变，算出误配后的矩阵，说明检查非负整数性可以发现什么，却不能保证什么。
\end{optionalexercise}
\begin{solution}
	由 \(F=D\Phi\)，可得
	\[
		F'=PD\Phi R=(PDQ^{-1})\Phi',\qquad D'=PDQ^{-1}.
	\]
	共同的共轭类重排 \(R\) 消失，因为普通与 Brauer 函数必须在同一个群元素上比较；普通行和简单行的重排则分别重排 \(D\) 的行与列。

	在正文的 \(S_3\) 次序下，错误数据给出
	\[
		\Phi'=\begin{pmatrix}1&1\\-1&1\end{pmatrix},\qquad
		F(\Phi')^{-1}=\begin{pmatrix}1&0\\0&-1\\1&-1\end{pmatrix}.
	\]
	负数说明这不能是在所声称次序下的分解矩阵。非负整数性因此是一项有用的错误检测，但通过检查仍不能代替证明输入表完整，也不能代替共轭类与简单模的身份匹配。有些错误重排可能仍产生非负整数矩阵；必须保留列标、行标和相应模的定义，才能解释每个数的含义。
\end{solution}
